% Generated mechanically from the frozen spaces-perfect.tex target.
% Do not edit this generated file; edit the bound locale source instead.

% OK, start here.
%

\setcounter{chapter}{74}
\chapter{空间的导出范畴}



\phantomsection
\label{spaces-perfect-section-phantom}


\section{引言}
\label{spaces-perfect-section-introduction}

\noindent
本章讨论代数空间上模的导出范畴。似乎没有直接论述这一主题的良好入门文献；
现有文献通常通过引用处理代数栈上模的导出范畴的论文来涵盖这一主题，例如见
\cite{olsson_sheaves}。



\section{约定}
\label{spaces-perfect-section-conventions}

\noindent
若 $\mathcal{A}$ 是阿贝尔范畴，$M$ 是 $\mathcal{A}$ 的对象，则我们仍以
$M$ 表示 $K(\mathcal{A})$ 和（或）$D(\mathcal{A})$ 中的相应对象；它对应的复形
以 $M$ 为次数 $0$ 处的项，在其余所有次数处为零。

\medskip\noindent
若给定环 $A$，则 $K(A)$ 表示 $A$-模复形的同伦范畴，$D(A)$ 表示相应的
导出范畴。类似地，若给定带环空间 $(X, \mathcal{O}_X)$，则符号
$K(\mathcal{O}_X)$ 表示 $\mathcal{O}_X$-模复形的同伦范畴，而
$D(\mathcal{O}_X)$ 表示相应的导出范畴。





\section{一般结果}
\label{spaces-perfect-section-generalities}

\noindent
本节给出关于代数空间上模的无界复形之上同调的若干一般结果。

\begin{lemma}
\label{spaces-perfect-lemma-restrict-direct-image-open}
设 $S$ 为概形，$f : X \to Y$ 为 $S$ 上代数空间的态射。给定平展态射
$V \to Y$，令 $U = V \times_Y X$，并以 $g : U \to V$ 表示投影态射。
则 $(Rf_*E)|_V = Rg_*(E|_U)$ 对 $E$（$D(\mathcal{O}_X)$ 中的对象）成立。
\end{lemma}

\begin{proof}
以 $E$ 的一个 K-单射复形 $\mathcal{I}^\bullet$ 表示，其中各项均为
$\mathcal{O}_X$-模。
于是 $Rf_*(E) = f_*\mathcal{I}^\bullet$，并且由“位点上的上同调”引理
\ref{sites-cohomology-lemma-restrict-K-injective-to-open}，
$Rg_*(E|_U) = g_*(\mathcal{I}^\bullet|_U)$。
因此结论由“代数空间的性质”引理
\ref{spaces-properties-lemma-pushforward-etale-base-change-modules} 得出。
\end{proof}


\begin{definition}
\label{spaces-perfect-definition-supported-on}
设 $S$ 为概形。设 $X$ 为 $S$ 上的代数空间。
设 $E$ 为 $D(\mathcal{O}_X)$ 中的对象。
设 $T \subset |X|$ 为闭子集。
称 $E$ {\it 支撑在 $T$ 上}，如果其上同调层 $H^i(E)$ 都支撑在 $T$ 上。
\end{definition}

\section{小平展位点上拟凝聚模的导出范畴}
\label{spaces-perfect-section-derived-quasi-coherent-etale}

\noindent
设 $X$ 为概形。本节证明 $D_\QCoh(\mathcal{O}_X)$ 可以用记作
$X_\etale$ 的 $X$ 的小平展位点来定义。以 $\mathcal{O}_\etale$ 表示
$X_\etale$ 上的结构层。
考虑带环位点的态射
\begin{equation}
\label{spaces-perfect-equation-epsilon}
\epsilon :
(X_\etale, \mathcal{O}_\etale)
\longrightarrow
(X_{Zar}, \mathcal{O}_X).
\end{equation}
在“下降”引理 \ref{descent-lemma-compare-sites} 中，它记作
$\text{id}_{small, \etale, Zar}$。

\begin{lemma}
\label{spaces-perfect-lemma-epsilon-flat}
(\ref{spaces-perfect-equation-epsilon}) 中的态射 $\epsilon$ 是带环位点的平坦态射。特别地，
函子
$\epsilon^* : \textit{Mod}(\mathcal{O}_X) \to
\textit{Mod}(\mathcal{O}_\etale)$ 正合。此外，若
$\epsilon^*\mathcal{F} = 0$，则 $\mathcal{F} = 0$。
\end{lemma}

\begin{proof}
态射 $\epsilon$ 的平坦性即“下降”引理
\ref{descent-lemma-compare-etale-zariski-flat}。下面给出另一证明。我们须证明
$\mathcal{O}_\etale$ 是平坦的 $\epsilon^{-1}\mathcal{O}_X$-模。为此，只需验证态射
$\mathcal{O}_{X, x} \to \mathcal{O}_{\etale, \overline{x}}$
对任意几何点 $\overline{x}$（属于 $X$）是平坦的；见“位点上的模”引理
\ref{sites-modules-lemma-check-flat-stalks}、“位点”引理
\ref{sites-lemma-point-morphism-sites} 以及“平展上同调”评注
\ref{etale-cohomology-remarks-enough-points}。
由“平展上同调”引理
\ref{etale-cohomology-lemma-describe-etale-local-ring}，
$\mathcal{O}_{\etale, \overline{x}}$ 是 $\mathcal{O}_{X, x}$ 的严格亨泽尔化。
因此由“代数进阶”引理
\ref{more-algebra-lemma-dumb-properties-henselization}，
$\mathcal{O}_{X, x} \to \mathcal{O}_{\etale, \overline{x}}$
是忠实平坦的。

\medskip\noindent
由“位点上的模”引理 \ref{sites-modules-lemma-flat-pullback-exact}，
$\epsilon^*$ 的正合性由 $\epsilon$ 的平坦性得出。

\medskip\noindent
设 $\mathcal{F}$ 为 $\mathcal{O}_X$-模。若
$\epsilon^*\mathcal{F} = 0$，则采用上述记号，有
$$
0 = \epsilon^*\mathcal{F}_{\overline{x}} =
\mathcal{F}_x \otimes_{\mathcal{O}_{X, x}} \mathcal{O}_{\etale, \overline{x}}
$$
对所有几何点 $\overline{x}$ 都成立（“位点上的模”引理
\ref{sites-modules-lemma-pullback-stalk}）。由
$\mathcal{O}_{X, x} \to \mathcal{O}_{\etale, \overline{x}}$
的忠实平坦性，可推出 $\mathcal{F}_x = 0$ 对所有 $x \in X$ 都成立。
\end{proof}

\noindent
设 $X$ 为概形，记号同 (\ref{spaces-perfect-equation-epsilon})。回忆，由“下降”命题
\ref{descent-proposition-equivalence-quasi-coherent} 及评注
\ref{descent-remark-change-topologies-ringed-sites}，
$\epsilon^* : \QCoh(\mathcal{O}_X)
\to \QCoh(\mathcal{O}_\etale)$ 是等价。此外，由“下降”引理
\ref{descent-lemma-equivalence-quasi-coherent-limits}，
$\QCoh(\mathcal{O}_\etale)$ 构成
$\textit{Mod}(\mathcal{O}_\etale)$ 的 Serre 子范畴。因此可以把
$D_\QCoh(\mathcal{O}_\etale)$ 定义为 $D(\mathcal{O}_\etale)$ 的三角子范畴，
其对象是上同调层均拟凝聚的复形；见“导出范畴”第
\ref{derived-section-triangulated-sub} 节。函子 $\epsilon^*$ 正合（引理
\ref{spaces-perfect-lemma-epsilon-flat}），故诱导
$\epsilon^* :  D(\mathcal{O}_X) \to D(\mathcal{O}_\etale)$；
又因拟凝聚模的拉回仍拟凝聚，它还诱导
$\epsilon^* : D_\QCoh(\mathcal{O}_X) \to
D_\QCoh(\mathcal{O}_\etale)$。

\begin{lemma}
\label{spaces-perfect-lemma-derived-quasi-coherent-small-etale-site}
设 $X$ 为概形。上面定义的函子
$\epsilon^* : D_\QCoh(\mathcal{O}_X) \to
D_\QCoh(\mathcal{O}_\etale)$ 是等价。
\end{lemma}

\begin{proof}
我们证明函子
$R\epsilon_* : D(\mathcal{O}_\etale) \to D(\mathcal{O}_X)$
诱导一个拟逆。以下将直接使用如下事实：$\epsilon_*$ 由限制到
$X_{Zar} \subset X_\etale$ 给出，而“下降”引理
\ref{descent-lemma-compare-sites} 描述了
$\epsilon^* = \text{id}_{small, \etale, Zar}^*$。

\medskip\noindent
对拟凝聚 $\mathcal{O}_X$-模 $\mathcal{F}$，伴随映射
$\mathcal{F} \to \epsilon_*\epsilon^*\mathcal{F}$ 是同构，因为
$\mathcal{F}^a$（“下降”定义
\ref{descent-definition-structure-sheaf}）是层；这一点由“下降”引理
\ref{descent-lemma-sheaf-condition-holds} 证明。反过来，每个拟凝聚
$\mathcal{O}_\etale$-模 $\mathcal{H}$ 都具有形式
$\epsilon^*\mathcal{F}$，其中某个拟凝聚 $\mathcal{O}_X$-模记为
$\mathcal{F}$；见“下降”命题
\ref{descent-proposition-equivalence-quasi-coherent}。由刚才所述，
$\mathcal{F} = \epsilon_*\mathcal{H}$，从而伴随映射
$\epsilon^*\epsilon_*\mathcal{H} \to \mathcal{H}$ 对每个拟凝聚
$\mathcal{O}_\etale$-模 $\mathcal{H}$ 都是同构。

\medskip\noindent
设 $E$ 为 $D_\QCoh(\mathcal{O}_\etale)$ 的对象，并以
$\mathcal{H}^q = H^q(E)$ 表示其第 $q$ 个上同调层。令
$\mathcal{B}$ 为 $X_\etale$ 的仿射对象所成的集合。等式
$H^p(U, \mathcal{H}^q) = 0$ 对所有 $p > 0$、所有
$q \in \mathbf{Z}$ 及所有 $U \in \mathcal{B}$ 都成立；见“下降”命题
\ref{descent-proposition-same-cohomology-quasi-coherent} 和“概形的上同调”引理
\ref{coherent-lemma-quasi-coherent-affine-cohomology-zero}。
由“位点上的上同调”引理
\ref{sites-cohomology-lemma-cohomology-over-U-trivial}，这意味着
$$
H^q(U, E) = H^0(U, \mathcal{H}^q)
$$
对所有 $U \in \mathcal{B}$ 都成立。特别地，它对仿射开集
$U \subset X$ 成立。因此，第 $q$ 个上同调——对 $R\epsilon_*E$ 而言——
在 $U$ 上等于层 $\epsilon_*\mathcal{H}^q$ 在 $U$ 上的值。层化后得到
$$
H^q(R\epsilon_*E) = \epsilon_*\mathcal{H}^q
$$
这特别表明 $R\epsilon_*$ 诱导函子
$D_\QCoh(\mathcal{O}_\etale) \to D_\QCoh(\mathcal{O}_X)$。
由于 $\epsilon^*$ 正合，进而有
$H^q(\epsilon^*R\epsilon_*E) = \epsilon^*\epsilon_*\mathcal{H}^q =
\mathcal{H}^q$（由上述讨论）。故伴随映射
$\epsilon^*R\epsilon_*E \to E$ 是同构。

\medskip\noindent
反过来，对 $F \in D_\QCoh(\mathcal{O}_X)$，伴随映射
$F \to R\epsilon_*\epsilon^*F$ 由于同样的理由也是同构；也就是说，
$R\epsilon_*\epsilon^*F$ 的上同调层同构于
$\epsilon_*H^m(\epsilon^*F) = \epsilon_*\epsilon^*H^m(F) = H^m(F)$。
\end{proof}


\section{拟凝聚模的导出范畴}
\label{spaces-perfect-section-derived-quasi-coherent}

\noindent
设 $S$ 为概形。引理
\ref{spaces-perfect-lemma-derived-quasi-coherent-small-etale-site}
表明，范畴 $D_\QCoh(\mathcal{O}_S)$ 可以用 $\mathcal{O}_S$-模复形
（在概形 $S$ 上）来定义，也可以用 $\mathcal{O}$-模复形
（在 $S$ 的小平展位点上）来定义。因此下面的定义与概形情形下的定义相容。

\begin{definition}
\label{spaces-perfect-definition-derived-quasi-coherent}
设 $S$ 为概形，$X$ 为 $S$ 上的代数空间。{\it 上同调层拟凝聚的
$\mathcal{O}_X$-模的导出范畴}记作 $D_\QCoh(\mathcal{O}_X)$。
\end{definition}

\noindent
由“代数空间的性质”引理
\ref{spaces-properties-lemma-properties-quasi-coherent}
及“导出范畴”第 \ref{derived-section-triangulated-sub} 节，这一定义是有意义的。
于是得到典范函子
\begin{equation}
\label{spaces-perfect-equation-compare}
D(\QCoh(\mathcal{O}_X))
\longrightarrow
D_\QCoh(\mathcal{O}_X)
\end{equation}
；见“导出范畴”等式 (\ref{derived-equation-compare})。

\medskip\noindent
注意，代数空间的平坦态射 $f : Y \to X$ 诱导正合函子
$f^* : \textit{Mod}(\mathcal{O}_X) \to \textit{Mod}(\mathcal{O}_Y)$；
见“代数空间的态射”引理
\ref{spaces-morphisms-lemma-flat-morphism-sites} 和“位点上的模”引理
\ref{sites-modules-lemma-flat-pullback-exact}。特别地，
$Lf^* : D(\mathcal{O}_X) \to D(\mathcal{O}_Y)$ 可以在任意代表复形上计算
（“导出范畴”引理 \ref{derived-lemma-right-derived-exact-functor}）。
我们将写作 $Lf^* = f^*$，如果 $f$ 平坦；在这种情形下，
$H^i(f^*E) = f^*H^i(E)$ 对 $E$（$D(\mathcal{O}_X)$ 中的对象）成立。
当 $f$ 平展时，我们将经常使用这一点。当然，在平展情形下，拉回函子就是
限制到 $Y_\etale$；见“代数空间的性质”等式
(\ref{spaces-properties-equation-restrict-modules})。

\begin{lemma}
\label{spaces-perfect-lemma-check-quasi-coherence-on-covering}
设 $S$ 为概形，$X$ 为 $S$ 上的代数空间，$E$ 为
$D(\mathcal{O}_X)$ 的对象。下列条件等价：
\begin{enumerate}
\item $E$ 属于 $D_\QCoh(\mathcal{O}_X)$；
\item 对每个平展态射 $\varphi : U \to X$，其中 $U$ 为仿射概形，
$\varphi^*E$ 是 $D_\QCoh(\mathcal{O}_U)$ 的对象；
\item 对每个平展态射 $\varphi : U \to X$，其中 $U$ 为概形，
$\varphi^*E$ 是 $D_\QCoh(\mathcal{O}_U)$ 的对象；
\item 存在满平展态射 $\varphi : U \to X$，其中 $U$ 为概形，使得
$\varphi^*E$ 是 $D_\QCoh(\mathcal{O}_U)$ 的对象；以及
\item 存在代数空间的满平展态射 $f : Y \to X$，使得
$Lf^*E$ 是 $D_\QCoh(\mathcal{O}_Y)$ 的对象。
\end{enumerate}
\end{lemma}

\begin{proof}
这由引理前的讨论和“代数空间的性质”引理
\ref{spaces-properties-lemma-characterize-quasi-coherent} 立即得出。
\end{proof}

\begin{lemma}
\label{spaces-perfect-lemma-quasi-coherence-direct-sums}
设 $S$ 为概形，$X$ 为 $S$ 上的代数空间。则
$D_\QCoh(\mathcal{O}_X)$ 中存在直和。
\end{lemma}

\begin{proof}
由“单射对象”引理 \ref{injectives-lemma-derived-products}，导出范畴
$D(\mathcal{O}_X)$ 中存在直和，并且可在任意代表上逐项取直和来计算。
因此，直和的上同调层显然就是上同调层的直和，因为取直和是正合函子
（在任意 Grothendieck 阿贝尔范畴中）。拟凝聚层的直和仍拟凝聚，见
“代数空间的性质”引理
\ref{spaces-properties-lemma-properties-quasi-coherent}，故引理得证。
\end{proof}

\noindent
我们还需要一些关于导出极限的信息。提醒读者，在下面的引理中，导出极限通常
不是 $D_\QCoh$ 的对象。

\begin{lemma}
\label{spaces-perfect-lemma-Rlim-quasi-coherent}
设 $S$ 为概形，$X$ 为 $S$ 上的代数空间。设 $(K_n)$ 是
$D_\QCoh(\mathcal{O}_X)$ 中的逆系统，其导出极限
$K = R\lim K_n$ 位于 $D(\mathcal{O}_X)$ 中。假设映射
$H^q(K_{n + 1}) \to H^q(K_n)$ 对所有 $q \in \mathbf{Z}$ 和
$n \geq 1$ 都是满射。则
\begin{enumerate}
\item $H^q(K) = \lim H^q(K_n)$；
\item $R\lim H^q(K_n) = \lim H^q(K_n)$；以及
\item 对每个仿射开集 $U \subset X$，有
$H^p(U, \lim H^q(K_n)) = 0$（当 $p > 0$）。
\end{enumerate}
\end{lemma}

\begin{proof}
令 $\mathcal{B} \subset \Ob(X_\etale)$ 为仿射对象所成的集合。由于
$H^q(K_n)$ 拟凝聚，故有 $H^p(U, H^q(K_n)) = 0$，其中
$U \in \mathcal{B}$；见“代数空间的上同调”第
\ref{spaces-cohomology-section-higher-direct-image} 节和“概形的上同调”引理
\ref{coherent-lemma-quasi-coherent-affine-cohomology-zero}。此外，由类似的理由，映射
$H^0(U, H^q(K_{n + 1})) \to H^0(U, H^q(K_n))$ 对
$U \in \mathcal{B}$ 是满射。第 (1) 项由“位点上的上同调”引理
\ref{sites-cohomology-lemma-derived-limit-suitable-system} 得出，因为我们刚验证了
该引理的条件。第 (2)、(3) 项由“位点上的上同调”引理
\ref{sites-cohomology-lemma-inverse-limit-is-derived-limit} 得出。
\end{proof}

\begin{lemma}
\label{spaces-perfect-lemma-quasi-coherence-pullback}
设 $S$ 为概形。设 $f : Y \to X$ 为 $S$ 上代数空间的态射。
函子 $Lf^*$ 把 $D_\QCoh(\mathcal{O}_X)$ 映入
$D_\QCoh(\mathcal{O}_Y)$。
\end{lemma}

\begin{proof}
选择一个图
$$
\xymatrix{
U \ar[d]_a \ar[r]_h & V \ar[d]^b \\
X \ar[r]^f & Y
}
$$
% P09-ERR-0393
其中 $U$ 和 $V$ 是概形，竖直箭头平展，并且 $a$ 是满射。由于
$a^* \circ Lf^* = Lh^* \circ b^*$，结论由引理
\ref{spaces-perfect-lemma-check-quasi-coherence-on-covering} 以及概形情形得出；概形情形即
“概形的导出范畴”引理
\ref{perfect-lemma-quasi-coherence-pullback}。
\end{proof}

\begin{lemma}
\label{spaces-perfect-lemma-quasi-coherence-tensor-product}
设 $S$ 为概形，$X$ 为 $S$ 上的代数空间。对对象 $K, L$
（属于 $D_\QCoh(\mathcal{O}_X)$），导出张量积
$K \otimes^\mathbf{L} L$ 属于 $D_\QCoh(\mathcal{O}_X)$。
\end{lemma}

\begin{proof}
设 $\varphi : U \to X$ 为从概形 $U$ 出发的满平展态射。由于
$\varphi^*(K \otimes_{\mathcal{O}_X}^\mathbf{L} L) =
\varphi^*K \otimes_{\mathcal{O}_U}^\mathbf{L} \varphi^*L$，由引理
\ref{spaces-perfect-lemma-check-quasi-coherence-on-covering} 可知，结论归结为概形情形；后者即
“概形的导出范畴”引理
\ref{perfect-lemma-quasi-coherence-tensor-product}。
\end{proof}

\noindent
下面的引理将帮助我们在 $D_\QCoh(\mathcal{O}_X)$ 的对象上“计算”右导出函子。

\begin{lemma}
\label{spaces-perfect-lemma-nice-K-injective}
设 $S$ 为概形，$X$ 为 $S$ 上的代数空间，$E$ 为
$D_\QCoh(\mathcal{O}_X)$ 的对象。则“导出范畴”评注
\ref{derived-remark-map-into-derived-limit-truncations} 中的映射
$E \to R\lim \tau_{\geq -n}E$ 是同构\footnote{特别地，
$E$ 有 K-单射代表；见“导出范畴”引理
\ref{derived-lemma-difficulty-K-injectives}。}。
\end{lemma}

\begin{proof}
以 $\mathcal{H}^i = H^i(E)$ 表示第 $i$ 个上同调层（对 $E$ 而言）。令
$\mathcal{B}$ 为 $X_\etale$ 的仿射对象所成的集合。于是
$H^p(U, \mathcal{H}^i) = 0$ 对所有 $p > 0$、所有
$i \in \mathbf{Z}$ 以及所有 $U \in \mathcal{B}$ 都成立，因为
$U$ 是仿射概形。见“代数空间的上同调”第
\ref{spaces-cohomology-section-higher-direct-image} 节和“概形的上同调”引理
\ref{coherent-lemma-quasi-coherent-affine-cohomology-zero}。
因此，本引理由“位点上的上同调”引理
\ref{sites-cohomology-lemma-is-limit-dimension}（取 $d = 0$）得出。
\end{proof}

\begin{lemma}
\label{spaces-perfect-lemma-application-nice-K-injective}
设 $S$ 为概形，$X$ 为 $S$ 上的代数空间。设
$F : \textit{Mod}(\mathcal{O}_X) \to \textit{Ab}$
为函子，$N \geq 0$ 为整数。假设
\begin{enumerate}
\item $F$ 左正合；
\item $F$ 与可数直积可交换；
\item $R^pF(\mathcal{F}) = 0$ 对所有 $p \geq N$ 及拟凝聚
$\mathcal{F}$ 都成立。
\end{enumerate}
则对 $E \in D_\QCoh(\mathcal{O}_X)$，有
\begin{enumerate}
\item $H^i(RF(\tau_{\leq a}E) \to H^i(RF(E))$ 在
$i \leq a$ 时为同构； % P09-ERR-0394
\item $H^i(RF(E)) \to H^i(RF(\tau_{\geq b - N + 1}E))$ 在
$i \geq b$ 时为同构；
\item 若 $H^i(E) = 0$ 在 $i \not \in [a, b]$ 时成立，其中某个
$-\infty \leq a \leq b \leq \infty$，则 $H^i(RF(E)) = 0$
在 $i \not \in [a, b + N - 1]$ 时成立。
\end{enumerate}
\end{lemma}

\begin{proof}
陈述 (1) 即“导出范畴”引理 \ref{derived-lemma-negative-vanishing}。

\medskip\noindent
证明陈述 (2)。写 $E_n = \tau_{\geq -n}E$。我们有
$E = R\lim E_n$，见引理 \ref{spaces-perfect-lemma-nice-K-injective}。于是由“单射对象”引理
\ref{injectives-lemma-RF-commutes-with-Rlim}，有
$RF(E) = R\lim RF(E_n)$（在 $D(\textit{Ab})$ 中）。因此对每个
$i \in \mathbf{Z}$，存在短正合列
$$
0 \to R^1\lim H^{i - 1}(RF(E_n)) \to H^i(RF(E)) \to \lim H^i(RF(E_n)) \to 0
$$
；见“代数进阶”评注
\ref{more-algebra-remark-compare-derived-limit}。为证明 (2)，我们将证明左端项为零，
且右端项等于 $H^i(RF(E_{-b + N - 1})$；这里 $b$ 任取并满足
$i \geq b$。 % P09-ERR-0395

\medskip\noindent
对每个 $n$，有特异三角
$$
H^{-n}(E)[n] \to E_n \to E_{n - 1} \to H^{-n}(E)[n + 1]
$$
（“导出范畴”评注
\ref{derived-remark-truncation-distinguished-triangle}），它位于
$D(\mathcal{O}_X)$ 中。由于 $H^{-n}(E)$ 拟凝聚，有
$$
H^i(RF(H^{-n}(E)[n])) = R^{i + n}F(H^{-n}(E)) = 0
$$
，当 $i + n \geq N$；并且有
$$
H^i(RF(H^{-n}(E)[n + 1])) = R^{i + n + 1}F(H^{-n}(E)) = 0
$$
，当 $i + n + 1 \geq N$。由此可知
$$
H^i(RF(E_n)) \to H^i(RF(E_{n - 1}))
$$
在 $n \geq N - i$ 时是同构。因此，各系统 $H^i(RF(E_n))$ 都满足
ML 条件，而短正合列中的 $R^1\lim$ 项为零（见“代数进阶”第
\ref{more-algebra-section-Rlim} 节中的讨论）。此外，系统
$H^i(RF(E_n))$ 从 $n = N - i - 1$ 起为常值系统，正如所需。

\medskip\noindent
证明 (3)。在关于 $E$ 的假设下，有 $\tau_{\leq a - 1}E = 0$，
由 (1) 得到 $H^i(RF(E))$ 在 $i \leq a - 1$ 时消失。类似地，我们有
$\tau_{\geq b + 1}E = 0$，因而由第 (2) 项得到
$H^i(RF(E))$ 在 $i \geq b + n$ 时消失。 % P09-ERR-0396
\end{proof}


\section{总正像}
\label{spaces-perfect-section-total-direct-image}

\noindent
下面的引理是“代数空间的上同调”引理
\ref{spaces-cohomology-lemma-vanishing-higher-direct-images} 的类似形式。

\begin{lemma}
\label{spaces-perfect-lemma-quasi-coherence-direct-image}
设 $S$ 为概形，$f : X \to Y$ 为 $S$ 上代数空间的拟分离拟紧态射。
\begin{enumerate}
\item 函子 $Rf_*$ 把 $D_\QCoh(\mathcal{O}_X)$ 映入
$D_\QCoh(\mathcal{O}_Y)$。
\item 若 $Y$ 拟紧，则存在整数 $N = N(X, Y, f)$，使得对对象 $E$
（属于 $D_\QCoh(\mathcal{O}_X)$），如果 $H^m(E) = 0$ 对
$m > 0$ 成立，那么 $H^m(Rf_*E) = 0$ 对 $m \geq N$ 成立。
\item 事实上，若 $Y$ 拟紧，则可找到 $N = N(X, Y, f)$，使得对代数空间的
每个态射 $Y' \to Y$，函子 $R(f')_*$ 都满足同一结论，其中
$f' : X' \to Y'$ 是 $f$ 的基变换。
\end{enumerate}
\end{lemma}

\begin{proof}
设 $E$ 为 $D_\QCoh(\mathcal{O}_X)$ 的对象。为证明 (1)，须证明
$Rf_*E$ 的上同调层拟凝聚。这个问题在 $Y$ 上是局部的，因此可以假设
$Y$ 拟紧。取“代数空间的上同调”引理
\ref{spaces-cohomology-lemma-vanishing-higher-direct-images} 中的
$N = N(X, Y, f)$。于是 $R^pf_*\mathcal{F} = 0$ 对所有拟凝聚
$\mathcal{O}_X$-模 $\mathcal{F}$ 及所有 $p \geq N$ 都成立。
此外，$R^pf_*\mathcal{F}$ 对所有 $p$ 均拟凝聚；见“代数空间的上同调”引理
\ref{spaces-cohomology-lemma-higher-direct-image}。这些陈述在基变换后仍成立。

\medskip\noindent
首先假设 $E$ 下有界。我们将证明，对这样的 $E$，以所选的 $N$，
(1)、(2)、(3) 均成立。在这种情形下，例如可以使用谱序列
$$
R^pf_*H^q(E) \Rightarrow R^{p + q}f_*E
$$
（“导出范畴”引理 \ref{derived-lemma-two-ss-complex-functor}），并使用
$R^pf_*H^q(E)$ 的拟凝聚性以及 $R^pf_*H^q(E)$ 在 $p \geq N$ 时的消失性，
从而看出 (1)、(2)、(3) 在这种情形下成立。

\medskip\noindent
接下来证明 (2)、(3)。设 $H^m(E) = 0$ 对 $m > 0$ 成立。令
$V$ 为 $Y_\etale$ 的仿射对象。我们有
$H^p(V \times_Y X, \mathcal{F}) = 0$（当 $p \geq N$）；见
“代数空间的上同调”引理
\ref{spaces-cohomology-lemma-quasi-coherence-higher-direct-images-application}。
因此可把引理 \ref{spaces-perfect-lemma-application-nice-K-injective} 应用于函子
$\Gamma(V \times_Y X, -)$，从而看出
$$
R\Gamma(V, Rf_*E) = R\Gamma(V \times_Y X, E)
$$
在次数 $\geq N$ 处的上同调消失。由于这对所有 $V$
（$Y_\etale$ 中仿射）均成立，可得 $H^m(Rf_*E) = 0$ 对
$m \geq N$ 成立。

\medskip\noindent
最后证明一般情形下的 (1)。回忆，有特异三角
$$
\tau_{\leq -n - 1}E \to E \to \tau_{\geq -n}E \to
(\tau_{\leq -n - 1}E)[1]
$$
，它位于 $D(\mathcal{O}_X)$ 中；见“导出范畴”评注
\ref{derived-remark-truncation-distinguished-triangle}。由 (2) 可知，
$Rf_*\tau_{\leq -n - 1}E$ 在次数 $\geq -n + N$ 处的上同调层消失。
因此，给定整数 $q$，可见 $R^qf_*E$ 等于
$R^qf_*\tau_{\geq -n}E$（对某个 $n$），于是可以应用上面的结果。
\end{proof}

\begin{lemma}
\label{spaces-perfect-lemma-quasi-coherence-pushforward-direct-sums}
设 $S$ 为概形，$f : X \to Y$ 为 $S$ 上代数空间的拟分离拟紧态射。则
$Rf_* : D_\QCoh(\mathcal{O}_X) \to D_\QCoh(\mathcal{O}_Y)$
与直和可交换。
\end{lemma}

\begin{proof}
设 $E_i$ 为 $D_\QCoh(\mathcal{O}_X)$ 的一族对象，并令
$E = \bigoplus E_i$。我们要证明映射
$$
\bigoplus Rf_*E_i \longrightarrow Rf_*E
$$
是同构。我们将证明它在次数 $0$ 的上同调层上诱导同构，这便蕴含引理。
选取引理 \ref{spaces-perfect-lemma-quasi-coherence-direct-image} 中的整数 $N$。
于是由所引引理，有 $R^0f_*E = R^0f_*\tau_{\geq -N}E$ 和
$R^0f_*E_i = R^0f_*\tau_{\geq -N}E_i$。注意
$\tau_{\geq -N}E = \bigoplus \tau_{\geq -N}E_i$。
因此可以假设所有 $E_i$ 在次数 $< -N$ 处的上同调层均消失。
接着使用谱序列
$$
R^pf_*H^q(E) \Rightarrow R^{p + q}f_*E
\quad\text{以及}\quad
R^pf_*H^q(E_i) \Rightarrow R^{p + q}f_*E_i
$$
（“导出范畴”引理 \ref{derived-lemma-two-ss-complex-functor}），把问题归结为
拟凝聚层的直和情形。该情形由“代数空间的上同调”引理
\ref{spaces-cohomology-lemma-colimit-cohomology} 处理。
\end{proof}

\begin{remark}
\label{spaces-perfect-remark-match-total-direct-images}
设 $S$ 为概形，$f : X \to Y$ 为可表代数空间 $X$ 与 $Y$
（在 $S$ 上）之间的态射。
设 $f_0 : X_0 \to Y_0$ 为表示 $f$ 的概形态射（这个记号很别扭，但只是临时的）。
则图
$$
\xymatrix{
D_\QCoh(\mathcal{O}_{X_0})
\ar@{=}[rrrrrr]_{\text{引理
\ref{spaces-perfect-lemma-derived-quasi-coherent-small-etale-site}}}
& & & & & &
D_\QCoh(\mathcal{O}_X) \\
D_\QCoh(\mathcal{O}_{Y_0})
\ar[u]^{Lf^*_0}
\ar@{=}[rrrrrr]^{\text{引理
\ref{spaces-perfect-lemma-derived-quasi-coherent-small-etale-site}}}
& & & & & &
D_\QCoh(\mathcal{O}_Y) \ar[u]_{Lf^*}
}
$$
可交换（引理 \ref{spaces-perfect-lemma-quasi-coherence-pullback} 和
“概形的导出范畴”引理
\ref{perfect-lemma-quasi-coherence-pullback}）。这是因为引理
\ref{spaces-perfect-lemma-derived-quasi-coherent-small-etale-site} 中的等价
$D_\QCoh(\mathcal{O}_{X_0}) \to D_\QCoh(\mathcal{O}_X)$ 和
$D_\QCoh(\mathcal{O}_{Y_0}) \to D_\QCoh(\mathcal{O}_Y)$
来自沿带环位点的（平坦）态射
$\epsilon : X_\etale \to X_{0, Zar}$ 和
$\epsilon : Y_\etale \to Y_{0, Zar}$ 的拉回，并且带环位点的图
$$
\xymatrix{
X_{0, Zar} \ar[d]_{f_0} & X_\etale \ar[l]^\epsilon \ar[d]^f \\
Y_{0, Zar} & Y_\etale \ar[l]_\epsilon
}
$$
可交换（略去细节）。若 $f$ 拟紧且拟分离——等价地，若 $f_0$ 拟紧且
拟分离——则我们断言
$$
\xymatrix{
D_\QCoh(\mathcal{O}_{X_0})
\ar[d]_{Rf_{0, *}} \ar@{=}[rrrrrr]_{\text{引理
\ref{spaces-perfect-lemma-derived-quasi-coherent-small-etale-site}}}
& & & & & &
D_\QCoh(\mathcal{O}_X) \ar[d]^{Rf_*} \\
D_\QCoh(\mathcal{O}_{Y_0})
\ar@{=}[rrrrrr]^{\text{引理
\ref{spaces-perfect-lemma-derived-quasi-coherent-small-etale-site}}}
& & & & & &
D_\QCoh(\mathcal{O}_Y)
}
$$
也可交换（引理 \ref{spaces-perfect-lemma-quasi-coherence-direct-image} 和
“概形的导出范畴”引理
\ref{perfect-lemma-quasi-coherence-direct-image}）。这也由上面所示的位点交换图得出，
因为引理 \ref{spaces-perfect-lemma-derived-quasi-coherent-small-etale-site} 的证明表明，函子
$R\epsilon_*$ 给出等价
$D_\QCoh(\mathcal{O}_X) \to D_\QCoh(\mathcal{O}_{X_0})$ 和
$D_\QCoh(\mathcal{O}_Y) \to D_\QCoh(\mathcal{O}_{Y_0})$。
\end{remark}

\begin{lemma}
\label{spaces-perfect-lemma-affine-morphism}
设 $S$ 为概形，$f : X \to Y$ 为 $S$ 上代数空间的仿射态射。则
$Rf_* : D_\QCoh(\mathcal{O}_X) \to D_\QCoh(\mathcal{O}_Y)$
反映同构。
\end{lemma}

\begin{proof}
该陈述的含义是：态射 $\alpha : E \to F$
（属于 $D_\QCoh(\mathcal{O}_X)$）若满足 $Rf_*\alpha$ 是同构，
则其本身也是同构。
这可以在上同调层上检验。特别地，这个问题在 $Y$ 上是平展局部的。
因此可以假设 $Y$，从而 $X$，是仿射的。在这种情形下，问题归结为概形情形
（“概形的导出范畴”引理 \ref{perfect-lemma-affine-morphism}）；这一归结通过
引理 \ref{spaces-perfect-lemma-derived-quasi-coherent-small-etale-site} 和评注
\ref{spaces-perfect-remark-match-total-direct-images} 实现。
\end{proof}

\begin{lemma}
\label{spaces-perfect-lemma-affine-morphism-pull-push}
设 $S$ 为概形，$f : X \to Y$ 为 $S$ 上代数空间的仿射态射。对 $E$
（$D_\QCoh(\mathcal{O}_Y)$ 中的对象），有
$Rf_* Lf^* E = E \otimes^\mathbf{L}_{\mathcal{O}_Y} f_*\mathcal{O}_X$。
\end{lemma}

\begin{proof}
由于 $f$ 仿射，映射 $f_*\mathcal{O}_X \to Rf_*\mathcal{O}_X$
是同构（“代数空间的上同调”引理
\ref{spaces-cohomology-lemma-affine-vanishing-higher-direct-images}）。
存在典范映射
$E \otimes^\mathbf{L} f_*\mathcal{O}_X =
E \otimes^\mathbf{L} Rf_*\mathcal{O}_X \to Rf_* Lf^* E$，
它与映射
$$
Lf^*(E \otimes^\mathbf{L} Rf_*\mathcal{O}_X) =
Lf^*E \otimes^\mathbf{L} Lf^*Rf_*\mathcal{O}_X \longrightarrow
Lf^* E \otimes^\mathbf{L} \mathcal{O}_X = Lf^* E
$$
伴随；后一个映射来自 $1 : Lf^*E \to Lf^*E$ 和典范映射
$Lf^*Rf_*\mathcal{O}_X \to \mathcal{O}_X$。为检验所构造的映射是同构，
可以在 $Y$ 上局部地工作。因此可以假设 $Y$，从而 $X$，是仿射的。
在这种情形下，问题归结为概形情形（“概形的导出范畴”引理
\ref{perfect-lemma-affine-morphism-pull-push}）；这一归结通过引理
\ref{spaces-perfect-lemma-derived-quasi-coherent-small-etale-site} 和评注
\ref{spaces-perfect-remark-match-total-direct-images} 实现。
\end{proof}


\section{在基底上为固有}
\label{spaces-perfect-section-proper-over-base}

\noindent
本节是“概形的上同调”第
\ref{coherent-section-proper-over-base} 节的类似形式。和通常一样，
只要材料涉及点集上的拓扑，把它转移到代数空间时就必须谨慎。

\begin{lemma}
\label{spaces-perfect-lemma-closed-proper-over-base}
设 $S$ 为概形，$f : X \to Y$ 为 $S$ 上代数空间的局部有限型态射。
设 $T \subset |X|$ 为闭子集。下列条件等价：
\begin{enumerate}
\item 态射 $Z \to Y$ 固有，如果 $Z$ 是 $T$ 上的约化诱导代数空间结构
（“代数空间的性质”定义
\ref{spaces-properties-definition-reduced-induced-space}）；
\item 对某个闭子空间 $Z \subset X$，若 $|Z| = T$，则态射
$Z \to Y$ 固有；以及
\item 对任意闭子空间 $Z \subset X$，若 $|Z| = T$，则态射
$Z \to Y$ 固有。
\end{enumerate}
\end{lemma}

\begin{proof}
蕴含关系 (3) $\Rightarrow$ (1) 和 (1) $\Rightarrow$ (2) 是直接的。
故只需证明 (2) 蕴含 (3)。我们力劝读者自行找出这个事实的证明。
设 $Z'$ 和 $Z''$ 为闭子空间，满足 $T = |Z'| = |Z''|$，并且
$Z' \to Y$ 是代数空间的固有态射。须证明 $Z'' \to Y$ 也固有。
令 $Z''' = Z' \cup Z''$ 为概形论并；见“代数空间的态射”定义
\ref{spaces-morphisms-definition-scheme-theoretic-intersection-union}。
于是 $Z'''$ 是另一个满足 $|Z'''| = T$ 的闭子空间。例如，这由
“代数空间的态射”引理
\ref{spaces-morphisms-lemma-scheme-theoretic-union} 中概形论并的描述得出。
由于 $Z'' \to Z'''$ 是闭浸入，只需证明 $Z''' \to Y$ 固有
（见“代数空间的态射”引理
\ref{spaces-morphisms-lemma-closed-immersion-proper} 和
\ref{spaces-morphisms-lemma-composition-proper}）。态射
$Z' \to Z'''$ 是双射的闭浸入，特别地，它是满的且普遍闭的。
于是 $Z' \to Y$ 分离这一事实蕴含 $Z''' \to Y$ 分离；见
“代数空间的态射”引理
\ref{spaces-morphisms-lemma-image-universally-closed-separated}。
此外，$Z''' \to Y$ 局部有限型，因为 $X \to Y$ 局部有限型
（“代数空间的态射”引理
\ref{spaces-morphisms-lemma-immersion-locally-finite-type} 和
\ref{spaces-morphisms-lemma-composition-finite-type}）。由于
$Z' \to Y$ 拟紧，且 $Z' \to Z'''$ 是普遍同胚，可见
$Z''' \to Y$ 拟紧。最后，因为 $Z' \to Y$ 普遍闭，由
“代数空间的态射”引理
\ref{spaces-morphisms-lemma-image-proper-is-proper} 可见
$Z''' \to Y$ 也普遍闭。这就完成了证明。
\end{proof}

\begin{definition}
\label{spaces-perfect-definition-proper-over-base}
设 $S$ 为概形。设 $f : X \to Y$ 为 $S$ 上代数空间的局部有限型态射。
设 $T \subset |X|$ 为闭子集。如果引理
\ref{spaces-perfect-lemma-closed-proper-over-base} 中的等价条件成立，我们称
{\it $T$ 在 $Y$ 上固有}。
\end{definition}

\noindent
若态射 $f : X \to Y$ 不是局部有限型的，则上面定义中使用的引理不成立。
因此我们力劝读者：若 $f$ 不是局部有限型的，就不要使用这个术语。

\begin{lemma}
\label{spaces-perfect-lemma-closed-closed-proper-over-base}
设 $S$ 为概形。设 $f : X \to Y$ 为 $S$ 上代数空间的局部有限型态射。
设 $T' \subset T \subset |X|$ 为闭子集。若 $T$ 在 $Y$ 上固有，
则 $T'$ 也是如此。
\end{lemma}

\begin{proof}
略。
\end{proof}

\begin{lemma}
\label{spaces-perfect-lemma-base-change-closed-proper-over-base}
设 $S$ 为概形。考虑 $S$ 上代数空间的笛卡尔图
$$
\xymatrix{
X' \ar[d]_{f'} \ar[r]_{g'} & X \ar[d]^f \\
Y' \ar[r]^g & Y
}
$$
，其中 $f$ 局部有限型。若 $T$ 是 $|X|$ 的闭子集并在 $Y$ 上固有，
则 $|g'|^{-1}(T)$ 是 $|X'|$ 的闭子集并在 $Y'$ 上固有。
\end{lemma}

\begin{proof}
注意，由“代数空间的态射”引理
\ref{spaces-morphisms-lemma-base-change-finite-type}，$f'$ 局部有限型，
故该陈述是有意义的。设 $Z \subset X$ 为 $T$ 上的约化诱导闭子空间结构。
以 $Z' = (g')^{-1}(Z)$ 表示概形论逆像。于是
$Z' = X' \times_X Z = (Y' \times_Y X) \times_X Z = Y' \times_Y Z$
在 $Y'$ 上固有，因为它是 $Z$ 在 $Y$ 上的基变换
（“代数空间的态射”引理
\ref{spaces-morphisms-lemma-base-change-proper}）。另一方面，
我们有 $T' = |Z'|$。 % P09-ERR-0397
故引理成立。
\end{proof}

\begin{lemma}
\label{spaces-perfect-lemma-functoriality-closed-proper-over-base}
设 $S$ 为概形，$B$ 为 $S$ 上的代数空间。设 $f : X \to Y$ 为代数空间的态射，
其源与靶均在 $B$ 上局部有限型。
\begin{enumerate}
\item 若 $Y$ 在 $B$ 上分离，且 $T \subset |X|$ 是在 $B$ 上固有的闭子集，
则 $|f|(T)$ 是 $|Y|$ 的闭子集并在 $B$ 上固有。
\item 若 $f$ 普遍闭，且 $T \subset |X|$ 是在 $B$ 上固有的闭子集，
则 $|f|(T)$ 是 $Y$ 的闭子集并在 $B$ 上固有。 % P09-ERR-0398
\item 若 $f$ 固有，且 $T \subset |Y|$ 是在 $B$ 上固有的闭子集，
则 $|f|^{-1}(T)$ 是 $|X|$ 的闭子集并在 $B$ 上固有。
\end{enumerate}
\end{lemma}

\begin{proof}
证明 (1)。假设 $Y$ 在 $B$ 上分离，且 $T \subset |X|$ 是在 $B$ 上固有的
闭子集。设 $Z$ 为 $T$ 上的约化诱导闭子空间结构，并把
“代数空间的态射”引理
\ref{spaces-morphisms-lemma-scheme-theoretic-image-is-proper}
应用于 $Z \to Y$（在 $B$ 上），即得结论。

\medskip\noindent
证明 (2)。假设 $f$ 普遍闭，且 $T \subset |X|$ 是在 $B$ 上固有的闭子集。
设 $Z$ 为 $T$ 上的约化诱导闭子空间结构，$Z'$ 为 $|f|(T)$ 上的
约化诱导闭子空间结构。于是得到诱导态射 $Z \to Z'$。
以 $Z'' = f^{-1}(Z')$ 表示概形论逆像。由于 $Z'' \to Z'$ 是 $f$ 的基变换，
它普遍闭（“代数空间的态射”引理
\ref{spaces-morphisms-lemma-base-change-proper}）。因此 $Z \to Z'$ 作为闭浸入
$Z \to Z''$ 与 $Z'' \to Z'$ 的复合而普遍闭
（“代数空间的态射”引理
\ref{spaces-morphisms-lemma-closed-immersion-proper} 和
\ref{spaces-morphisms-lemma-composition-proper}）。由“代数空间的态射”引理
\ref{spaces-morphisms-lemma-image-universally-closed-separated}，可得
$Z' \to B$ 分离。由于 $Z \to B$ 拟紧且 $Z \to Z'$ 满，可见
$Z' \to B$ 拟紧。因为 $Z' \to B$ 是 $Z' \to Y$ 与 $Y \to B$ 的复合，
可见 $Z' \to B$ 局部有限型
（“代数空间的态射”引理
\ref{spaces-morphisms-lemma-immersion-locally-finite-type} 和
\ref{spaces-morphisms-lemma-composition-finite-type}）。最后，因为
$Z \to B$ 普遍闭，由“代数空间的态射”引理
\ref{spaces-morphisms-lemma-image-proper-is-proper} 可见
$Z' \to B$ 也普遍闭。这就完成了证明。

\medskip\noindent
证明 (3)。假设 $f$ 固有，且 $T \subset |Y|$ 是在 $B$ 上固有的闭子集。
设 $Z$ 为 $T$ 上的约化诱导闭子空间结构。以
$Z' = f^{-1}(Z)$ 表示概形论逆像。由于 $Z' \to Z$ 是 $f$ 的基变换，
它固有（“代数空间的态射”引理
\ref{spaces-morphisms-lemma-base-change-proper}）。从而 $Z' \to B$ 作为
$Z' \to Z$ 与 $Z \to B$ 的复合而固有
（“代数空间的态射”引理
\ref{spaces-morphisms-lemma-composition-proper}）。这就完成了证明。
\end{proof}

\begin{lemma}
\label{spaces-perfect-lemma-union-closed-proper-over-base}
设 $S$ 为概形。设 $f : X \to Y$ 为 $S$ 上代数空间的局部有限型态射。
设 $T_i \subset |X|$（$i = 1, \ldots, n$）为闭子集。若
$T_i$（$i = 1, \ldots, n$）均在 $Y$ 上固有，则
$T_1 \cup \ldots \cup T_n$ 也是如此。
\end{lemma}

\begin{proof}
设 $Z_i$ 为 $T_i$ 上的约化诱导闭子概形结构。 % P09-ERR-0399
态射
$$
Z_1 \amalg \ldots \amalg Z_n \longrightarrow X
$$
是有限的；见“代数空间的态射”引理
\ref{spaces-morphisms-lemma-closed-immersion-finite} 和
\ref{spaces-morphisms-lemma-finite-union-finite}。有限态射普遍闭
（“代数空间的态射”引理 \ref{spaces-morphisms-lemma-finite-proper}），
并且 $Z_1 \amalg \ldots \amalg Z_n$ 在 $S$ 上固有， % P09-ERR-0400
故由引理 \ref{spaces-perfect-lemma-functoriality-closed-proper-over-base} 第 (2) 项可得，
像 $Z_1 \cup \ldots \cup Z_n$ 在 $S$ 上固有。
\end{proof}

\noindent
设 $S$ 为概形。设 $f : X \to Y$ 为 $S$ 上代数空间的局部有限型态射。
设 $\mathcal{F}$ 为有限型拟凝聚 $\mathcal{O}_X$-模。于是支撑
$\text{Supp}(\mathcal{F})$（属于 $\mathcal{F}$）是 $|X|$ 的闭子集；
见“代数空间的态射”引理
\ref{spaces-morphisms-lemma-support-finite-type}。因此，说
“$\mathcal{F}$ 的支撑在 $Y$ 上固有”是有意义的。

\begin{lemma}
\label{spaces-perfect-lemma-module-support-proper-over-base}
设 $S$ 为概形，$f : X \to Y$ 为 $S$ 上代数空间的局部有限型态射。
设 $\mathcal{F}$ 为有限型拟凝聚 $\mathcal{O}_X$-模。下列条件等价：
\begin{enumerate}
\item $\mathcal{F}$ 的支撑在 $Y$ 上固有；
\item $\mathcal{F}$ 的概形论支撑（“代数空间的态射”定义
\ref{spaces-morphisms-definition-scheme-theoretic-support}）在 $Y$ 上固有；以及
\item 存在闭子空间 $Z \subset X$ 和有限型拟凝聚
$\mathcal{O}_Z$-模 $\mathcal{G}$，使得 (a) $Z \to Y$ 固有，并且 (b)
$(Z \to X)_*\mathcal{G} = \mathcal{F}$。
\end{enumerate}
\end{lemma}

\begin{proof}
支撑 $\text{Supp}(\mathcal{F})$（属于 $\mathcal{F}$）是 $|X|$ 的闭子集；
见“代数空间的态射”引理
\ref{spaces-morphisms-lemma-support-finite-type}。故可应用定义
\ref{spaces-perfect-definition-proper-over-base}。由于 $\mathcal{F}$ 的概形论支撑是闭子空间，
其底闭子集为 $\text{Supp}(\mathcal{F})$，由定义
\ref{spaces-perfect-definition-proper-over-base} 可见 (1)、(2) 等价。显然 (2) 蕴含 (3)。
反过来，若 (3) 成立，则
$\text{Supp}(\mathcal{F}) \subset |Z|$，因而例如由引理
\ref{spaces-perfect-lemma-closed-closed-proper-over-base}，$\text{Supp}(\mathcal{F})$
在 $Y$ 上固有。
\end{proof}

\begin{lemma}
\label{spaces-perfect-lemma-base-change-module-support-proper-over-base}
设 $S$ 为概形。考虑 $S$ 上代数空间的笛卡尔图
$$
\xymatrix{
X' \ar[d]_{f'} \ar[r]_{g'} & X \ar[d]^f \\
Y' \ar[r]^g & Y
}
$$
，其中 $f$ 局部有限型。设 $\mathcal{F}$ 为有限型拟凝聚
$\mathcal{O}_X$-模。若 $\mathcal{F}$ 的支撑在 $Y$ 上固有，则
$(g')^*\mathcal{F}$ 的支撑在 $Y'$ 上固有。
\end{lemma}

\begin{proof}
注意，该陈述是有意义的，因为 $(g')*\mathcal{F}$ 由“位点上的模”引理
\ref{sites-modules-lemma-local-pullback} 为有限型。 % P09-ERR-0401
由“代数空间的态射”引理
\ref{spaces-morphisms-lemma-support-finite-type}，有
$\text{Supp}((g')^*\mathcal{F}) = |g'|^{-1}(\text{Supp}(\mathcal{F}))$。
于是本引理由引理
\ref{spaces-perfect-lemma-base-change-closed-proper-over-base} 得出。
\end{proof}

\begin{lemma}
\label{spaces-perfect-lemma-cat-module-support-proper-over-base}
设 $S$ 为概形。设 $f : X \to Y$ 为 $S$ 上代数空间的局部有限型态射。
设 $\mathcal{F}$、$\mathcal{G}$ 为有限型拟凝聚 $\mathcal{O}_X$-模。
\begin{enumerate}
\item 若 $\mathcal{F}$、$\mathcal{G}$ 的支撑在 $Y$ 上固有，则下列各模也具有
这一性质：$\mathcal{F} \oplus \mathcal{G}$；
$\mathcal{G}$ 被 $\mathcal{F}$ 扩张所得的任意扩张；
$\Im(u)$ 和 $\Coker(u)$（给定任意 $\mathcal{O}_X$-模映射
$u : \mathcal{F} \to \mathcal{G}$）；以及 $\mathcal{F}$ 或
$\mathcal{G}$ 的任意拟凝聚商。
\item 若 $Y$ 局部诺特，则凝聚 $\mathcal{O}_X$-模中支撑在 $Y$ 上固有者
所成的范畴，是凝聚 $\mathcal{O}_X$-模的阿贝尔范畴的
Serre 子范畴（“同调”定义
\ref{homology-definition-serre-subcategory}）。
\end{enumerate}
\end{lemma}

\begin{proof}
证明 (1)。设 $T$、$T'$ 分别为 $\mathcal{F}$、$\mathcal{G}$ 的支撑。
则 (1) 中提到的所有层，其支撑均包含于 $T \cup T'$。因此，只要检验这些层
有限型且拟凝聚，断言本身便由引理
\ref{spaces-perfect-lemma-closed-closed-proper-over-base} 和
\ref{spaces-perfect-lemma-union-closed-proper-over-base} 直接得出。关于拟凝聚性，参见
“代数空间的性质”第 \ref{spaces-properties-section-quasi-coherent} 节。
关于“有限型”，参见“代数空间的性质”第
\ref{spaces-properties-section-properties-modules} 节。

\medskip\noindent
证明 (2)。证明与 (1) 的证明相同。注意，这些断言是有意义的，因为由
“代数空间的态射”引理
\ref{spaces-morphisms-lemma-locally-finite-type-locally-noetherian}，
$X$ 局部诺特；还要使用“代数空间的上同调”第
\ref{spaces-cohomology-section-coherent} 节中对凝聚模范畴的描述。
\end{proof}

\begin{lemma}
\label{spaces-perfect-lemma-support-proper-over-base-pushforward}
设 $S$ 为概形，$f : X \to Y$ 为 $S$ 上代数空间的态射。假设 $f$ 局部有限型，
且 $Y$ 局部诺特。设 $\mathcal{F}$ 为凝聚 $\mathcal{O}_X$-模，其支撑在
$Y$ 上固有。则 $R^pf_*\mathcal{F}$ 是凝聚 $\mathcal{O}_Y$-模，
对所有 $p \geq 0$ 都成立。
\end{lemma}

\begin{proof}
由引理 \ref{spaces-perfect-lemma-module-support-proper-over-base}，存在闭浸入
$i : Z \to X$，使得 $g = f \circ i : Z \to Y$ 固有，并且有
$\mathcal{F} = i_*\mathcal{G}$，其中 $\mathcal{G}$ 是凝聚模（在 $Z$ 上）。
由“代数空间的上同调”引理
\ref{spaces-cohomology-lemma-proper-pushforward-coherent}，可见
$R^pg_*\mathcal{G}$ 在 $S$ 上凝聚。 % P09-ERR-0402
另一方面，有 $R^qi_*\mathcal{G} = 0$（当 $q > 0$）
（“代数空间的上同调”引理
\ref{spaces-cohomology-lemma-finite-pushforward-coherent}）。
由“位点上的上同调”引理
\ref{sites-cohomology-lemma-relative-Leray}，得到
$R^pf_*\mathcal{F} = R^pg_*\mathcal{G}$，引理得证。
\end{proof}

\section{凝聚模的导出范畴}
\label{spaces-perfect-section-derived-coherent}

\noindent
设 $S$ 为概形。设 $X$ 为 $S$ 上的局部诺特代数空间。
此时，范畴
$\textit{Coh}(\mathcal{O}_X) \subset \textit{Mod}(\mathcal{O}_X)$
（即凝聚 $\mathcal{O}_X$-模所成的范畴）是一个弱 Serre 子范畴；参见“同调”第
\ref{homology-section-serre-subcategories} 节以及
“代数空间的上同调”引理
\ref{spaces-cohomology-lemma-coherent-abelian-Noetherian}。
记
$$
D_{\textit{Coh}}(\mathcal{O}_X) \subset D(\mathcal{O}_X)
$$
为上同调层均凝聚的复形所成的子范畴；参见“导出范畴”第
\ref{derived-section-triangulated-sub} 节。因此得到典范函子
\begin{equation}
\label{spaces-perfect-equation-compare-coherent}
D(\textit{Coh}(\mathcal{O}_X))
\longrightarrow
D_{\textit{Coh}}(\mathcal{O}_X)
\end{equation}
参见“导出范畴”等式 (\ref{derived-equation-compare})。

\begin{lemma}
\label{spaces-perfect-lemma-direct-image-coherent}
设 $S$ 为概形。设 $f : X \to Y$ 为 $S$ 上代数空间的态射。
假设 $f$ 局部有限型，且 $Y$ 诺特。设 $E$ 为
$D^b_{\textit{Coh}}(\mathcal{O}_X)$ 的对象，并且
$H^i(E)$ 的支撑在 $Y$ 上固有，对所有 $i$ 均成立。则 $Rf_*E$ 是
$D^b_{\textit{Coh}}(\mathcal{O}_Y)$ 的对象。
\end{lemma}

\begin{proof}
考虑谱序列
$$
R^pf_*H^q(E) \Rightarrow R^{p + q}f_*E
$$
参见“导出范畴”引理 \ref{derived-lemma-two-ss-complex-functor}。
由假设和引理 \ref{spaces-perfect-lemma-support-proper-over-base-pushforward}，
层 $R^pf_*H^q(E)$ 是凝聚的。因此
$R^{p + q}f_*E$ 凝聚，即 $E \in D_{\textit{Coh}}(\mathcal{O}_Y)$。 % P09-ERR-0403
下有界性显然。上有界性由“代数空间的上同调”引理
\ref{spaces-cohomology-lemma-vanishing-higher-direct-images} 推出，
亦可由引理 \ref{spaces-perfect-lemma-quasi-coherence-direct-image} 推出。
\end{proof}

\begin{lemma}
\label{spaces-perfect-lemma-direct-image-coherent-bdd-below}
设 $S$ 为概形。设 $f : X \to Y$ 为 $S$ 上代数空间的态射。
假设 $f$ 局部有限型，且 $Y$ 诺特。设 $E$ 为
$D^+_{\textit{Coh}}(\mathcal{O}_X)$ 的对象，并且
$H^i(E)$ 的支撑在 $S$ 上固有，对所有 $i$ 均成立。 % P09-ERR-0404
则 $Rf_*E$ 是 $D^+_{\textit{Coh}}(\mathcal{O}_Y)$ 的对象。
\end{lemma}

\begin{proof}
证明与引理 \ref{spaces-perfect-lemma-direct-image-coherent} 的证明相同。
考察正合函子 $Rf_*$ 对特异三角
$\tau_{\leq a}E \to E \to \tau_{\geq a + 1}E \to \tau_{\leq a}E[1]$
的作用，也可由引理 \ref{spaces-perfect-lemma-direct-image-coherent} 推出此结论。
\end{proof}

\begin{lemma}
\label{spaces-perfect-lemma-coherent-internal-hom}
设 $S$ 为概形。设 $X$ 为 $S$ 上的局部诺特代数空间。
若 $L$ 属于 $D^+_{\textit{Coh}}(\mathcal{O}_X)$，且
$K$ 属于 $D^-_{\textit{Coh}}(\mathcal{O}_X)$，则
$R\SheafHom(K, L)$ 属于 $D^+_{\textit{Coh}}(\mathcal{O}_X)$。
\end{lemma}

\begin{proof}
一个 $D(\mathcal{O}_X)$ 的对象是否属于
$D_{\textit{Coh}}(\mathcal{O}_X)$，可在 $X$ 上平展局部地检验；
参见“代数空间的上同调”引理
\ref{spaces-cohomology-lemma-coherent-Noetherian}。
因此，本引理由概形情形推出；参见“概形的导出范畴”引理
\ref{perfect-lemma-coherent-internal-hom}。
\end{proof}

\begin{lemma}
\label{spaces-perfect-lemma-ext-finite}
设 $A$ 为诺特环。设 $X$ 为 $A$ 上的固有代数空间。
若 $L$ 属于 $D^+_{\textit{Coh}}(\mathcal{O}_X)$，且 $K$ 属于
$D^-_{\textit{Coh}}(\mathcal{O}_X)$，则 $A$-模
$\Ext_{\mathcal{O}_X}^n(K, L)$ 均有限生成。
\end{lemma}

\begin{proof}
回顾
$$
\Ext_{\mathcal{O}_X}^n(K, L) =
H^n(X, R\SheafHom_{\mathcal{O}_X}(K, L)) =
H^n(\Spec(A), Rf_*R\SheafHom_{\mathcal{O}_X}(K, L))
$$
参见“位点上的上同调”引理
\ref{sites-cohomology-lemma-section-RHom-over-U} 以及“位点上的上同调”第
\ref{sites-cohomology-section-leray} 节。因此，结论由引理
\ref{spaces-perfect-lemma-coherent-internal-hom} 和
\ref{spaces-perfect-lemma-direct-image-coherent-bdd-below} 推出。
\end{proof}

\section{归纳原理}
\label{spaces-perfect-section-induction}

\noindent
本节讨论代数空间的一条归纳原理，它类似于概形情形的“概形的上同调”引理
\ref{coherent-lemma-induction-principle}。为表述这条原理，我们引入
{\it 初等特异方形}的概念；这一术语取自 \cite{MV}。
这里所表述的原理隐含于 Raynaud 与 Gruson 的论文 \cite{GruRay} 中。
David Rydh 的 \cite[Theorem D]{rydh_etale_devissage}
给出了代数栈的一条相关原理。

\begin{definition}
\label{spaces-perfect-definition-elementary-distinguished-square}
设 $S$ 为概形。交换图
$$
\xymatrix{
U \times_W V \ar[r] \ar[d] & V \ar[d]^f \\
U \ar[r]^j & W
}
$$
是 $S$ 上代数空间的图，称为{\it 初等特异方形}，如果
\begin{enumerate}
\item $U$ 是 $W$ 的开子空间，且 $j$ 是包含态射；
\item $f$ 平展；并且
\item 令 $T = W \setminus U$（赋予约化诱导子空间结构），则态射
$f^{-1}(T) \to T$ 是同构。
\end{enumerate}
我们将用如下说法表示这一点：“设 $(U \subset W, f : V \to W)$
为初等特异方形。”
\end{definition}

\noindent
注意，若 $(U \subset W, f : V \to W)$ 为初等特异方形，则
$W = U \cup f(V)$。因此，$\{U \to W, V \to W\}$ 是 $W$ 的平展覆盖。
事实表明，这些平展覆盖具有良好性质，而且在某种意义下它们“足够多”。

\begin{lemma}
\label{spaces-perfect-lemma-make-more-elementary-distinguished-squares}
设 $S$ 为概形。设 $(U \subset W, f : V \to W)$ 为 $S$ 上代数空间的
初等特异方形。
\begin{enumerate}
\item 若 $V' \subset V$ 和 $U \subset U' \subset W$ 是开子空间，且
$W' = U' \cup f(V')$，则
$(U' \subset W', f|_{V'} : V' \to W')$ 是初等特异方形。
\item 若 $p : W' \to W$ 是代数空间的态射，则
$(p^{-1}(U) \subset W', V \times_W W' \to W')$ 是初等特异方形。
\item 若 $S' \to S$ 是概形的态射，则
$(S' \times_S U \subset S' \times_S W, S' \times_S V \to S' \times_S W)$
是初等特异方形。
\end{enumerate}
\end{lemma}

\begin{proof}
略。
\end{proof}

\begin{lemma}
\label{spaces-perfect-lemma-induction-principle}
设 $S$ 为概形。设 $X$ 为 $S$ 上拟紧且拟分离的代数空间。设 $P$ 是
$X_{spaces, \etale}$ 的拟紧拟分离对象的一项性质。假设
\begin{enumerate}
\item $P$ 对 $X_{spaces, \etale}$ 的每个仿射对象成立；
\item 对每个初等特异方形 $(U \subset W, f : V \to W)$，若
\begin{enumerate}
\item $W$ 是 $X_{spaces, \etale}$ 的拟紧拟分离对象；
\item $U$ 拟紧；
\item $V$ 仿射；并且
\item $P$ 对 $U$、$V$ 和 $U \times_W V$ 成立，
\end{enumerate}
则 $P$ 对 $W$ 成立。
\end{enumerate}
于是，$P$ 对 $X_{spaces, \etale}$ 的每个拟紧拟分离对象成立，特别地对
$X$ 成立。
\end{lemma}

\begin{proof}
首先断言，$P$ 对 $X_{spaces, \etale}$ 的每个可表拟紧拟分离对象成立。
事实上，设 $U \to X$ 平展，且 $U$ 是拟紧拟分离概形。由假设 (1)，
性质 $P$ 对 $U$ 的每个仿射开集成立。此外，若 $W, V \subset U$
是拟紧开集，$V$ 仿射，而且 $P$ 对 $W$、$V$ 和 $W \cap V$ 成立，
则由 (2)，$P$ 对 $W \cup V$ 成立（因为
$(W \subset W \cup V, V \to W \cup V)$ 这一对是初等特异方形）。
因此，由概形的归纳原理，$P$ 对 $U$ 成立；参见“概形的上同调”引理
\ref{coherent-lemma-induction-principle}。

\medskip\noindent
为完成证明，只需证明 $P$ 对 $X$ 成立（因为可将当前的 $X$ 替换为
$X_{spaces, \etale}$ 中我们想证明结论的任意拟紧拟分离对象）。
我们将使用“良态空间”引理
\ref{decent-spaces-lemma-filter-quasi-compact-quasi-separated} 中的滤过
$$
\emptyset = U_{n + 1} \subset
U_n \subset U_{n - 1} \subset \ldots \subset U_1 = X
$$
以及态射 $f_p : V_p \to U_p$。我们将证明 $P$ 对 $U_p$ 成立，并对
$p$ 作降序归纳。注意，由 (1)，$P$ 对 $U_{n + 1}$ 成立，
因为空代数空间是仿射的。假设 $P$ 对 $U_{p + 1}$ 成立。
注意，$(U_{p + 1} \subset U_p, f_p : V_p \to U_p)$ 是初等特异方形，
但 (2) 可能不适用，因为 $V_p$ 未必仿射。然而，由于 $V_p$ 是拟紧概形，
可选取有限仿射开覆盖
$V_p = V_{p, 1} \cup \ldots \cup V_{p, m}$。令
$W_{p, 0} = U_{p + 1}$，并令
$$
W_{p, i} = U_{p + 1} \cup f_p(V_{p, 1} \cup \ldots \cup V_{p, i})
$$
其中 $i = 1, \ldots, m$。这些是 $X$ 的拟紧开子空间。于是
$$
U_{p + 1} = W_{p, 0} \subset
W_{p, 1} \subset \ldots \subset
W_{p, m} = U_p
$$
且各对
$$
(W_{p, 0} \subset W_{p, 1}, f_p|_{V_{p, 1}}),
(W_{p, 1} \subset W_{p, 2}, f_p|_{V_{p, 2}}),\ldots,
(W_{p, m - 1} \subset W_{p, m}, f_p|_{V_{p, m}})
$$
由引理 \ref{spaces-perfect-lemma-make-more-elementary-distinguished-squares}
均为初等特异方形。注意，$P$ 对每个 $V_{p, 1}$ 均成立（因为它们是
仿射概形）； % P09-ERR-0405
而 $W_{p, i} \times_{W_{p, i + 1}} V_{p, i + 1}$ 是
$V_{p, i + 1}$ 的拟紧开集，所以由本证明第一段，$P$ 对它也成立。
因此，(2) 可用于其中每一对；归纳可得 $P$ 对
$W_{p, 1}, \ldots, W_{p, m} = U_p$ 成立。
\end{proof}

\begin{lemma}
\label{spaces-perfect-lemma-induction-principle-separated}
设 $S$ 为概形。设 $X$ 为 $S$ 上拟紧且拟分离的代数空间。设
$\mathcal{B} \subset \Ob(X_{spaces, \etale})$。设 $P$ 是
$\mathcal{B}$ 中元素的一项性质。假设
\begin{enumerate}
\item 每个 $W \in \mathcal{B}$ 都拟紧且拟分离；
\item 若 $W \in \mathcal{B}$，且 $U \subset W$ 是拟紧开集，则
$U \in \mathcal{B}$；
\item 若 $V \in \Ob(X_{spaces, \etale})$ 仿射，则
(a) $V \in \mathcal{B}$，并且 (b) $P$ 对 $V$ 成立；
\item 对每个初等特异方形 $(U \subset W, f : V \to W)$，若
\begin{enumerate}
\item $W \in \mathcal{B}$；
\item $U$ 拟紧；
\item $V$ 仿射；并且
\item $P$ 对 $U$、$V$ 和 $U \times_W V$ 成立，
\end{enumerate}
则 $P$ 对 $W$ 成立。
\end{enumerate}
于是，$P$ 对每个 $W \in \mathcal{B}$ 成立。
\end{lemma}

\begin{proof}
证明与引理 \ref{spaces-perfect-lemma-induction-principle} 的证明完全相同。
（补充说明，(4)(d) 是有意义的，因为 $U \times_W V$ 是 $V$ 的拟紧开集，
所以由条件 (2) 和 (3)，它是 $\mathcal{B}$ 的元素。）
\end{proof}

\begin{remark}
\label{spaces-perfect-remark-how-to}
在引理 \ref{spaces-perfect-lemma-induction-principle-separated} 中，应如何选择集合
$\mathcal{B}$？以下是一些例子：
\begin{enumerate}
\item 若 $X$ 拟紧且分离，则可取 $\mathcal{B}$ 为
$X_{spaces, \etale}$ 的拟紧分离对象所成的集合。这时
$X \in \mathcal{B}$，且 $\mathcal{B}$ 满足 (1)、(2) 和 (3)(a)。
在这样选择 $\mathcal{B}$ 时，引理
\ref{spaces-perfect-lemma-induction-principle-separated} 重现引理
\ref{spaces-perfect-lemma-induction-principle}。
\item 若 $X$ 拟紧，且其在 $\mathbf{Z}$ 上的对角态射仿射（见“代数空间的性质”定义
\ref{spaces-properties-definition-separated}），则可取 $\mathcal{B}$ 为
$X_{spaces, \etale}$ 中拟紧且在 $\mathbf{Z}$ 上具有仿射对角态射的对象所成的集合。
同样，$X \in \mathcal{B}$，且 $\mathcal{B}$ 满足 (1)、(2) 和 (3)(a)。
\item 若 $X$ 拟紧且拟分离，则满足 (1)、(2) 和 (3)(a) 的最小子集
$\mathcal{B}$（它包含 $X$）由下述规则给出：$W \in \mathcal{B}$ 当且仅当，或者
$W$ 是 $X$ 的拟紧开子空间，或者 $W$ 是 $X_{spaces, \etale}$
某个仿射对象的拟紧开集。
\end{enumerate}
\end{remark}

\noindent
下面给出一个变式，将某个开子空间上的成立性扩张到更大的开子空间。

\begin{lemma}
\label{spaces-perfect-lemma-induction-principle-enlarge}
设 $S$ 为概形。设 $X$ 为 $S$ 上拟紧且拟分离的代数空间。
设 $W \subset X$ 为拟紧开子空间。设 $P$ 是 $X$ 的拟紧开子空间的一项性质。
假设
\begin{enumerate}
\item $P$ 对 $W$ 成立；并且
\item 对每个初等特异方形
$(W_1 \subset W_2, f : V \to W_2)$，若满足下列条件： % P09-ERR-0406
\begin{enumerate}
\item $W_1$、$W_2$ 是 $X$ 的拟紧开子空间；
\item $W \subset W_1$；
\item $V$ 仿射；并且
\item $P$ 对 $W_1$ 成立，
\end{enumerate}
则 $P$ 对 $W_2$ 成立。
\end{enumerate}
于是，$P$ 对 $X$ 成立。
\end{lemma}

\begin{proof}
这可由引理 \ref{spaces-perfect-lemma-induction-principle-separated} 推出；不过，我们将直接
重做引理 \ref{spaces-perfect-lemma-induction-principle} 的证明，给出一个直接论证。
我们将使用滤过
$$
\emptyset = U_{n + 1} \subset
U_n \subset U_{n - 1} \subset \ldots \subset U_1 = X
$$
以及“良态空间”引理
\ref{decent-spaces-lemma-filter-quasi-compact-quasi-separated} 中的态射
$f_p : V_p \to U_p$。我们将证明 $P$ 对
$W_p = W \cup U_p$ 成立，并对 $p$ 作降序归纳。由于 $W_1 = X$，这将完成证明。
注意，由 (1)，$P$ 对 $W_{n + 1} = W \cap U_{n + 1} = W$ 成立。 % P09-ERR-0407
假设 $P$ 对 $W_{p + 1}$ 成立。注意，
$W_p \setminus W_{p + 1}$（赋予约化诱导子空间结构）是
$U_p \setminus U_{p + 1}$ 的闭子空间。由于
$(U_{p + 1} \subset U_p, f_p : V_p \to U_p)$ 是初等特异方形，
$(W_{p + 1} \subset W_p, f_p : V_p \to W_p)$ 也是如此。
然而，(2) 可能不适用，因为 $V_p$ 未必仿射。
然而，由于 $V_p$ 是拟紧概形，可选取有限仿射开覆盖
$V_p = V_{p, 1} \cup \ldots \cup V_{p, m}$。令
$W_{p, 0} = W_{p + 1}$，并令
$$
W_{p, i} = W_{p + 1} \cup f_p(V_{p, 1} \cup \ldots \cup V_{p, i})
$$
其中 $i = 1, \ldots, m$。这些是 $X$ 的拟紧开子空间，并包含 $W$。于是
$$
W_{p + 1} = W_{p, 0} \subset
W_{p, 1} \subset \ldots \subset
W_{p, m} = W_p
$$
且各对
$$
(W_{p, 0} \subset W_{p, 1}, f_p|_{V_{p, 1}}),
(W_{p, 1} \subset W_{p, 2}, f_p|_{V_{p, 2}}),\ldots,
(W_{p, m - 1} \subset W_{p, m}, f_p|_{V_{p, m}})
$$
由引理 \ref{spaces-perfect-lemma-make-more-elementary-distinguished-squares}
均为初等特异方形。现在，(2) 可用于其中每一对；归纳可得 $P$ 对
$W_{p, 1}, \ldots, W_{p, m} = W_p$ 成立。
\end{proof}

\section{Mayer--Vietoris}
\label{spaces-perfect-section-mayer-vietoris}

\noindent
本节证明，一个初等特异三角会导出各种 Mayer--Vietoris 序列。 % P09-ERR-0408

\medskip\noindent
设 $S$ 为概形。设 $U \to X$ 为 $S$ 上代数空间的平展态射。在
“代数空间的性质”第 \ref{spaces-properties-section-localize} 节中已经证明，

$U_{spaces, \etale} = X_{spaces, \etale}/U$

并且与结构层相容。因此，在这种情形下，我们常把态射 $j_U : U \to X$
看作局部化态射（参见“位点上的模”定义
\ref{sites-modules-definition-localize-ringed-site}）。特别地，我们把拉回
$j_U^*$ 看作限制到 $U$，并常将其记为 ${}|_U$；这与
“代数空间的性质”等式
(\ref{spaces-properties-equation-restrict-modules}) 相容。特别地，我们看到，
\begin{equation}
\label{spaces-perfect-equation-stalk-restriction}
(\mathcal{F}|_U)_{\overline{u}} = \mathcal{F}_{\overline{x}}
\end{equation}
若 $\overline{u}$ 是 $U$ 的几何点，且 $\overline{x}$ 是
$\overline{u}$ 在 $X$ 中的像。此外，限制有一个正合左伴随 $j_{U!}$；参见“位点上的模”引理
\ref{sites-modules-lemma-extension-by-zero} 和
\ref{sites-modules-lemma-extension-by-zero-exact}。
最后，回顾若 $\mathcal{G}$ 是一个 $\mathcal{O}_X$-模，则有
\begin{equation}
\label{spaces-perfect-equation-stalk-j-shriek}
(j_{U!}\mathcal{G})_{\overline{x}} =
\bigoplus\nolimits_{\overline{u}} \mathcal{G}_{\overline{u}}
\end{equation}
对任意几何点 $\overline{x} : \Spec(k) \to X$，其中直和取遍态射
$\overline{u} : \Spec(k) \to U$，这些态射满足 $j_U \circ \overline{u} = \overline{x}$；参见“位点上的模”引理
\ref{sites-modules-lemma-stalk-j-shriek} 以及
“代数空间的性质”引理
\ref{spaces-properties-lemma-points-small-etale-site}。

\begin{lemma}
\label{spaces-perfect-lemma-exact-sequence-lower-shriek}
设 $S$ 为概形。设 $(U \subset X, V \to X)$ 为 $S$ 上代数空间的初等特异方形。
\begin{enumerate}
\item 对于 $\mathcal{O}_X$-模层 $\mathcal{F}$，有短正合列
$$
0 \to j_{U \times_X V!}\mathcal{F}|_{U \times_X V} \to
j_{U!}\mathcal{F}|_U \oplus j_{V!}\mathcal{F}|_V \to \mathcal{F} \to 0
$$
\item 对于对象 $E$（属于 $D(\mathcal{O}_X)$），有特异三角
$$
j_{U \times_X V!}E|_{U \times_X V} \to
j_{U!}E|_U \oplus j_{V!}E|_V \to E \to
j_{U \times_X V!}E|_{U \times_X V}[1]
$$
在 $D(\mathcal{O}_X)$ 中。
\end{enumerate}
\end{lemma}

\begin{proof}
为证明 (1) 中的列正合，根据“代数空间的性质”定理
\ref{spaces-properties-theorem-exactness-stalks}，可在几何点处取茎来检验。
设 $\overline{x}$ 为 $X$ 的几何点。由等式
(\ref{spaces-perfect-equation-stalk-restriction}) 和 (\ref{spaces-perfect-equation-stalk-j-shriek})，
在 $\overline{x}$ 处取茎得到序列
$$
0 \to
\bigoplus\nolimits_{(\overline{u}, \overline{v})} \mathcal{F}_{\overline{x}}
\to
\bigoplus\nolimits_{\overline{u}} \mathcal{F}_{\overline{x}}
\oplus
\bigoplus\nolimits_{\overline{v}} \mathcal{F}_{\overline{x}}
\to
\mathcal{F}_{\overline{x}} \to 0
$$
此序列正合，因为对每个 $\overline{x}$，要么恰有一个
$\overline{u}$ 映到 $\overline{x}$，要么没有 $\overline{u}$，且恰有一个
$\overline{v}$ 映到 $\overline{x}$。

\medskip\noindent
证明 (2)。我们在“位点上的上同调”第
\ref{sites-cohomology-section-properties-K-injective} 节中看到，
导出范畴上的限制函子和零延拓函子，只需将函子逐项作用于任意复形即可计算。
设 $\mathcal{E}^\bullet$ 为 $\mathcal{O}_X$-模复形，表示对象 $E$。
本引理的特异三角是下述 $\mathcal{O}_X$-模复形短正合列所关联的特异三角，
其中所用关联来自“导出范畴”第
\ref{derived-section-canonical-delta-functor} 节，尤其是引理
\ref{derived-lemma-derived-canonical-delta-functor}：
$$
0 \to j_{U \times_X V!}\mathcal{E}^\bullet|_{U \times_X V} \to
j_{U!}\mathcal{E}^\bullet|_U \oplus j_{V!}\mathcal{E}^\bullet|_V
\to \mathcal{E}^\bullet \to 0
$$
由 (1)，此列短正合。
\end{proof}

\begin{lemma}
\label{spaces-perfect-lemma-exact-sequence-j-star}
设 $S$ 为概形。设 $(U \subset X, V \to X)$ 为 $S$ 上代数空间的初等特异方形。
\begin{enumerate}
\item 对每个 $\mathcal{O}_X$-模层 $\mathcal{F}$，有短正合列
$$
0 \to \mathcal{F} \to
j_{U, *}\mathcal{F}|_U \oplus j_{V, *}\mathcal{F}|_V \to
j_{U \times_X V, *}\mathcal{F}|_{U \times_X V} \to 0
$$
\item 对于对象 $E$（属于 $D(\mathcal{O}_X)$），有特异三角
$$
E \to
Rj_{U, *}E|_U \oplus Rj_{V, *}E|_V \to
Rj_{U \times_X V, *}E|_{U \times_X V} \to
E[1]
$$
在 $D(\mathcal{O}_X)$ 中。
\end{enumerate}
\end{lemma}

\begin{proof}
设 $W$ 为 $X_\etale$ 中的对象。我们断言，序列
$$
0 \to
\mathcal{F}(W) \to
\mathcal{F}(W \times_X U) \oplus \mathcal{F}(W \times_X V) \to
\mathcal{F}(W \times_X U \times_X V)
$$
正合，并且末项中的元素可在 $W$ 上局部地提升到中间项。
由引理 \ref{spaces-perfect-lemma-make-more-elementary-distinguished-squares}，
$(W \times_X U \subset W, V \times_X W \to W)$ 是初等特异方形。
因此可设 $W = X$，只需证明
$$
0 \to
\mathcal{F}(X) \to
\mathcal{F}(U) \oplus \mathcal{F}(V) \to
\mathcal{F}(U \times_X V)
$$
我们已经在引理 \ref{spaces-perfect-lemma-exact-sequence-lower-shriek} 中看到，
$$
0 \to j_{U \times_X V!}\mathcal{O}_{U \times_X V}
\to j_{U!}\mathcal{O}_U \oplus
j_{V!}\mathcal{O}_V \to
\mathcal{O}_X \to 0
$$
是 $\mathcal{O}_X$-模的正合列，而对其施加右正合函子
$\Hom_{\mathcal{O}_X}(- , \mathcal{F})$，得到上面的序列。 % P09-ERR-0409
这也意味着，将 $s \in \mathcal{F}(U \times_X V)$ 提升为
$\mathcal{F}(U) \oplus \mathcal{F}(V)$ 中元素的障碍属于
$\Ext^1_{\mathcal{O}_X}(\mathcal{O}_X, \mathcal{F}) = H^1(X, \mathcal{F})$。
由上同调的局部性（“位点上的上同调”引理
\ref{sites-cohomology-lemma-kill-cohomology-class-on-covering}），
此障碍在 $X$ 上平展局部地消失，故 (1) 的证明完成。

\medskip\noindent
证明 (2)。选取 K-单射复形 $\mathcal{I}^\bullet$ 表示对象 $E$，其各项
$\mathcal{I}^n$ 是 $\textit{Mod}(\mathcal{O}_X)$ 中的单射对象；参见“单射”定理
\ref{injectives-theorem-K-injective-embedding-grothendieck}。
于是 $\mathcal{I}^\bullet|U$ 是 K-单射复形（“位点上的上同调”引理
\ref{sites-cohomology-lemma-restrict-K-injective-to-open}）。因此，
$Rj_{U, *}E|_U$ 由 $j_{U, *}\mathcal{I}^\bullet|_U$ 表示。
对 $V$ 及 $U \times_X V$ 也同理。因此，本引理的特异三角是下述复形短正合列
所关联的特异三角，其中所用关联来自“导出范畴”第
\ref{derived-section-canonical-delta-functor} 节，尤其是引理
\ref{derived-lemma-derived-canonical-delta-functor}：
$$
0 \to
\mathcal{I}^\bullet \to
j_{U, *}\mathcal{I}^\bullet|_U \oplus j_{V, *}\mathcal{I}^\bullet|_V \to
j_{U \times_X V, *}\mathcal{I}^\bullet|_{U \times_X V} \to
0.
$$
由 (1)，此列正合。
\end{proof}
\begin{lemma}
\label{spaces-perfect-lemma-unbounded-relative-mayer-vietoris}
设 $S$ 为概形。设 $f : X \to Y$ 为
$S$ 上代数空间之间的态射。设 $(U \subset X, V \to X)$ 为初等特异方形。
记
$a = f|_U : U \to Y$、$b = f|_V : V \to Y$，以及
$c = f|_{U \times_X V} : U \times_X V \to Y$ 为相应限制。
对于对象 $E$（属于 $D(\mathcal{O}_X)$），存在特异三角
$$
Rf_*E \to
Ra_*(E|_U) \oplus Rb_*(E|_V) \to
Rc_*(E|_{U \times_X V}) \to
Rf_*E[1]
$$
在 $D(\mathcal{O}_Y)$ 中。该三角关于 $E$ 具有函子性。
\end{lemma}

\begin{proof}
选取 K-单射复形 $\mathcal{I}^\bullet$ 表示对象 $E$。不妨设
$\mathcal{I}^n$ 都是 $\textit{Mod}(\mathcal{O}_X)$ 中的单射对象，对所有 $n$ 均成立；参见“单射”定理
\ref{injectives-theorem-K-injective-embedding-grothendieck}。
于是 $Rf_*E$ 由 $f_*\mathcal{I}^\bullet$ 计算。
对 $U$、$V$ 以及 $U \cap V$ 也同理 % P09-ERR-0410
，参见“位点上的上同调”引理
\ref{sites-cohomology-lemma-restrict-K-injective-to-open}。
因此，本引理的特异三角是下述短正合复形列所关联的特异三角，其中所用关联来自
“导出范畴”第 \ref{derived-section-canonical-delta-functor} 节以及尤其是引理
\ref{derived-lemma-derived-canonical-delta-functor}：
$$
0 \to
f_*\mathcal{I}^\bullet \to
a_*\mathcal{I}^\bullet|_U \oplus b_*\mathcal{I}^\bullet|_V \to
c_*\mathcal{I}^\bullet|_{U \times_X V} \to
0.
$$
为说明这是复形的短正合列，论证如下。取单射对象 $\mathcal{I}$（属于
$\textit{Mod}(\mathcal{O}_X)$）。对下列短正合列施加 $f_*$：
$$
0 \to \mathcal{I} \to
j_{U, *}\mathcal{I}|_U \oplus j_{V, *}\mathcal{I}|_V \to
j_{U \times_X V, *}\mathcal{I}|_{U \times_X V} \to 0
$$
该列来自引理 \ref{spaces-perfect-lemma-exact-sequence-j-star}；再利用
$R^1f_*\mathcal{I} = 0$，得到短正合列
$$
0 \to f_*\mathcal{I} \to
f_*j_{U, *}\mathcal{I}|_U \oplus f_*j_{V, *}\mathcal{I}|_V \to
f_*j_{U \times_X V, *}\mathcal{I}|_{U \times_X V} \to 0
$$
最后注意到 $a_* = f_*j_{U, *}$，对 $b_*$ 和 $c_*$ 也同样成立，即得证明。
\end{proof}

\begin{lemma}
\label{spaces-perfect-lemma-mayer-vietoris-hom}
设 $S$ 为概形。设 $(U \subset X, V \to X)$ 为 $S$ 上代数空间的初等特异方形。
对于对象 $E$、$F$（属于 $D(\mathcal{O}_X)$），有 Mayer--Vietoris 序列
$$
\xymatrix{
& \ldots \ar[r] &
\Ext^{-1}(E_{U \times_X V}, F_{U \times_X V}) \ar[lld] \\
\Hom(E, F) \ar[r] &
\Hom(E_U, F_U) \oplus
\Hom(E_V, F_V) \ar[r] &
\Hom(E_{U \times_X V}, F_{U \times_X V})
}
$$
其中下标表示限制到相应开子空间，并且各个 $\Hom$ 都是在相应导出范畴中取的。
\end{lemma}

\begin{proof}
使用引理 \ref{spaces-perfect-lemma-exact-sequence-lower-shriek} 的特异三角，
得到 $\Hom$ 的长正合序列（由“导出范畴”引理
\ref{derived-lemma-representable-homological} 得到），并利用
\[
\Hom(j_{U!}E|_U, F) = \Hom(E|_U, F|_U)
\]
以及“位点上的上同调”引理
\ref{sites-cohomology-lemma-adjoint-lower-shriek-restrict}。
\end{proof}

\begin{lemma}
\label{spaces-perfect-lemma-unbounded-mayer-vietoris}
设 $S$ 为概形。设 $(U \subset X, V \to X)$ 为 $S$ 上代数空间的初等特异方形。
对于对象 $E$（属于 $D(\mathcal{O}_X)$），有特异三角
$$
R\Gamma(X, E) \to R\Gamma(U, E) \oplus R\Gamma(V, E) \to
R\Gamma(U \times_X V, E) \to R\Gamma(X, E)[1]
$$
以及特别的长正合上同调序列
$$
\ldots \to
H^n(X, E) \to
H^n(U, E) \oplus H^n(V, E) \to
H^n(U \times_X V, E) \to
H^{n + 1}(X, E) \to \ldots
$$
特异三角和长正合序列的构造关于 $E$ 具有函子性。
\end{lemma}

\begin{proof}
选取 K-单射复形 $\mathcal{I}^\bullet$ 表示对象 $E$，其各项
$\mathcal{I}^n$ 都是 $\textit{Mod}(\mathcal{O}_X)$ 中的单射对象；参见“单射”定理
\ref{injectives-theorem-K-injective-embedding-grothendieck}。
在引理 \ref{spaces-perfect-lemma-exact-sequence-j-star} 的证明中，我们得到复形的短正合列
$$
0 \to \mathcal{I}^\bullet \to
j_{U, *}\mathcal{I}^\bullet|_U \oplus j_{V, *}\mathcal{I}^\bullet|_V \to
j_{U \times_X V, *}\mathcal{I}^\bullet|_{U \times_X V} \to 0
$$
由于 $H^1(X, \mathcal{I}^n) = 0$，可知取整体截面得到复形的正合列
$$
0 \to \Gamma(X, \mathcal{I}^\bullet) \to
\Gamma(U, \mathcal{I}^\bullet) \oplus
\Gamma(V, \mathcal{I}^\bullet) \to
\Gamma(U \times_X V, \mathcal{I}^\bullet) \to 0
$$
这些复形分别表示
$R\Gamma(X, E)$、$R\Gamma(U, E)$、$R\Gamma(V, E)$ 以及
$R\Gamma(U \times_X V, E)$，故由“导出范畴”第
\ref{derived-section-canonical-delta-functor} 节以及尤其是引理
\ref{derived-lemma-derived-canonical-delta-functor}，
得到特异三角。
\end{proof}

\begin{lemma}
\label{spaces-perfect-lemma-restrict-lower-shriek}
设 $S$ 为概形。设 $j : U \to X$ 为 $S$ 上代数空间之间的平展态射。
给定平展态射 $V \to Y$，令 $W = V \times_X U$，并记投影态射 % P09-ERR-0411
$j_W : W \to V$。则
$(j_!E)|_V = j_{W!}(E|_W)$，其中 $E$ 属于 $D(\mathcal{O}_U)$。
\end{lemma}

\begin{proof}
这是因为
$(j_!\mathcal{F})|_V = j_{W!}(\mathcal{F}|_W)$
对于 $\mathcal{O}_X$-模 $\mathcal{F}$ 成立；这直接来自函子 $j_!$ 和 $j_{W!}$ 的构造；参见“位点上的模”引理
\ref{sites-modules-lemma-extension-by-zero}。
\end{proof}

\begin{lemma}
\label{spaces-perfect-lemma-pushforward-with-support-in-open}
设 $S$ 为概形。设 $(U \subset X, j : V \to X)$ 为 $S$ 上代数空间的初等特异方形。
置 $T = |X| \setminus |U|$。
\begin{enumerate}
\item 若 $E$ 是 $D(\mathcal{O}_X)$ 中支撑在 $T$ 上的对象，则
(a) $E \to Rj_*(E|_V)$ 以及 (b) $j_!(E|_V) \to E$ 都是同构。
\item 若 $F$ 是 $D(\mathcal{O}_V)$ 中支撑在 $j^{-1}T$ 上的对象，则
(a) $F \to (j_!F)|_V$、(b) $(Rj_*F)|_V \to F$，以及 (c)
$j_!F \to Rj_*F$ 都是同构。
\end{enumerate}
\end{lemma}

\begin{proof}
设 $E$ 为 $D(\mathcal{O}_X)$ 中的对象，且其上同调层支撑在 $T$ 上。
于是 $E|_U = 0$ 且 $E|_{U \times_X V} = 0$，因为 $T$ 不与 $U$ 相交，
而 $j^{-1}T$ 不与 $U \times_X V$ 相交。因此 (1)(a) 由引理
\ref{spaces-perfect-lemma-exact-sequence-j-star} 得出。完全同样地，(1)(b) 由引理
\ref{spaces-perfect-lemma-exact-sequence-lower-shriek} 得出。

\medskip\noindent
设 $F$ 为 $D(\mathcal{O}_V)$ 中的对象，且其上同调层支撑在 $j^{-1}T$ 上。
由引理 \ref{spaces-perfect-lemma-restrict-direct-image-open}，有
$(Rj_*F)|_U = Rj_{W, *}(F|_W) = 0$，因为按假设 $F|_W = 0$。
同理，由引理 \ref{spaces-perfect-lemma-restrict-lower-shriek}，
$(j_!F)|_U = j_{W!}(F|_W) = 0$。
因此 $j_!F$ 和 $Rj_*F$ 支撑在 $T$ 上，而
$(j_!F)|_V$ 和 $(Rj_*F)|_V$ 支撑在 $j^{-1}(T)$ 上。根据“代数空间的性质”定理
\ref{spaces-properties-theorem-exactness-stalks}，要在导出范畴中检验映射
(2)(a)、(b)、(c) 为同构，只需检验它们在 $T$ 和 $j^{-1}(T)$ 的几何点处
诱导的上同调层茎上的映射为同构。 % P09-ERR-0412
为此，可作如下替换：将 $X$ 换为 $V$，将 $U$ 换为
$U \times_X V$，将 $V$ 换为 $V \times_X V$，并将 $F$ 换为
$F|_{V \times_X V}$（通过第一投影限制）；参见引理
\ref{spaces-perfect-lemma-restrict-direct-image-open}、
\ref{spaces-perfect-lemma-restrict-lower-shriek} 以及
\ref{spaces-perfect-lemma-make-more-elementary-distinguished-squares}。
由于 $V \times_X V \to V$ 有截面，这将 (2) 归结为
$j : V \to X$ 有截面的情形。

\medskip\noindent
设 $j$ 有截面 $\sigma : X \to V$。置 $V' = \sigma(X)$，这是 $V$ 的开子空间。
置 $U' = j^{-1}(U)$，这是 $V$ 的另一个开子空间。
于是 $(U' \subset V, V' \to V)$ 是初等特异方形。
注意，$F|_{U'} = 0$ 且 $F|_{V' \cap U'} = 0$，因为 $F$ 支撑在
$j^{-1}(T)$ 上。记开浸入 $j' : V' \to V$，并记
$j_{V'} : V' \to X$ 为复合态射 $V' \to V \to X$，它是 $\sigma$ 的逆。
置 $F' = \sigma^*F$。引理
\ref{spaces-perfect-lemma-exact-sequence-lower-shriek} 和
\ref{spaces-perfect-lemma-exact-sequence-j-star} 的特异三角表明
$F = j'_!(F|_{V'})$ 且 $F = Rj'_*(F|_{V'})$。
于是
$j_!F = j_!j'_!(F|_{V'}) = j_{V'!}F = F'$，
因为 $j_{V'} : V' \to X$ 是同构且是 $\sigma$ 的逆。
同理，$Rj_*F = Rj_*Rj'_*F = Rj_{V', *}F = F'$。
这证明了 (2)(c)。为证明 (2)(a) 和 (2)(b)，只需证明
$F = F'|_V$。这是显然的，因为 $F$ 和 $F'|_V$ 在 $U'$ 及
$U' \cap V'$ 上都限制为零，并且在 $V'$ 上是同一个对象。
\end{proof}

\noindent
我们可以粘合复形！

\begin{lemma}
\label{spaces-perfect-lemma-glue}
设 $S$ 为概形。设 $(U \subset X, V \to X)$ 为 $S$ 上代数空间的初等特异方形。
给定
\begin{enumerate}
\item 对象 $A$，属于 $D(\mathcal{O}_U)$，
\item 对象 $B$，属于 $D(\mathcal{O}_V)$，以及
\item 同构 $c : A|_{U \times_X V} \to B|_{U \times_X V}$。
\end{enumerate}
则存在对象 $F$，属于 $D(\mathcal{O}_X)$，以及同构
$f : F|_U \to A$、$g : F|_V \to B$，使得
$c = g|_{U \times_X V} \circ f^{-1}|_{U \times_X V}$。
此外，给定
\begin{enumerate}
\item 对象 $E$，属于 $D(\mathcal{O}_X)$，
\item 态射 $a : A \to E|_U$，属于 $D(\mathcal{O}_U)$，
\item 态射 $b : B \to E|_V$，属于 $D(\mathcal{O}_V)$，
\end{enumerate}
满足
$$
a|_{U \times_X V}  = b|_{U \times_X V} \circ c.
$$
则存在态射 $F \to E$，属于 $D(\mathcal{O}_X)$，
其限制到 $U$ 为 $a \circ f$，限制到 $V$ 为 $b \circ g$。
\end{lemma}

\begin{proof}
记 $j_U$、$j_V$、$j_{U \times_X V}$ 为指向 $X$ 的相应态射。选取特异三角
$$
F \to Rj_{U, *}A \oplus Rj_{V, *}B \to
Rj_{U \times_X V, *}(B|_{U \times_X V}) \to F[1]
$$
其中映射 $Rj_{V, *}B \to Rj_{U \times_X V, *}(B|_{U \times_X V})$
是显然的。映射
$Rj_{U, *}A \to Rj_{U \times_X V, *}(B|_{U \times_X V})$
是复合映射：先取
$Rj_{U, *}A \to Rj_{U \times_X V, *}(A|_{U \times_X V})$，
再与 $Rj_{U \times_X V, *}c$ 复合。将其限制到 $U$，得到
$$
F|_U \to A \oplus (Rj_{V, *}B)|_U \to
(Rj_{U \times_X V, *}(B|_{U \times_X V}))|_U \to F|_U[1]
$$
记 $j : U \times_X V \to U$。根据引理
\ref{spaces-perfect-lemma-restrict-direct-image-open}，限制与总正像相容，故
$(Rj_{V, *}B)|_U$ 以及
$(Rj_{U \times_X V, *}(B|_{U \times_X V}))|_U$
都典范同构于 $Rj_*(B|_{U \times_X V})$。
因此上面最后一个显示式中的第二个箭头有截面，从而
$F|_U \to A$ 是同构。

\medskip\noindent
为证明 $F|_V \to B$ 是同构，我们使用一个技巧。也就是说，选取特异三角
$$
F|_V \to B \to B' \to F[1]|_V
$$
在 $D(\mathcal{O}_V)$ 中。由于 $F|_U \to A$ 是同构，且有同构
$c : A|_{U \times_X V} \to B|_{U \times_X V}$，
$F|_V \to B$ 在 $U \times_X V$ 上的限制是同构。
因此 $B'$ 支撑在 $j_V^{-1}(T)$ 上，其中 $T = |X| \setminus |U|$。
另一方面，有特异三角的态射
$$
\xymatrix{
F \ar[r] \ar[d] &
Rj_{U, *}F|_U \oplus Rj_{V, *}F|_V \ar[r] \ar[d] &
Rj_{U \times_X V, *}F|_{U \times_X V} \ar[r] \ar[d] &
F[1] \ar[d] \\
F \ar[r] &
Rj_{U, *}A \oplus Rj_{V, *}B \ar[r] &
Rj_{U \times_X V, *}(B|_{U \times_X V}) \ar[r] &
F[1]
}
$$
该图中所有竖直映射都是同构，除了
$Rj_{V, *}F|_V \to Rj_{V, *}B$；因此由“导出范畴”引理
\ref{derived-lemma-third-isomorphism-triangle}，该映射也为同构。 % P09-ERR-0413
这意味着 $Rj_{V, *}B' = 0$。于是由引理
\ref{spaces-perfect-lemma-pushforward-with-support-in-open} 得到 $B' = 0$。

\medskip\noindent
态射 $F \to E$ 的存在性由 Mayer--Vietoris 的 $\Hom$ 序列得出，参见引理
\ref{spaces-perfect-lemma-mayer-vietoris-hom}。
\end{proof}

\section{凝聚化函子}
\label{spaces-perfect-section-coherator}

\noindent
设 $S$ 为概形。设 $X$ 为 $S$ 上的代数空间。
{\it 凝聚化函子}是函子
$$
Q_X :
\textit{Mod}(\mathcal{O}_X)
\longrightarrow
\QCoh(\mathcal{O}_X)
$$
它是包含函子
$\QCoh(\mathcal{O}_X) \to \textit{Mod}(\mathcal{O}_X)$ 的右伴随。
它对任意代数空间 $X$ 都存在；此外，伴随映射
$Q_X(\mathcal{F}) \to \mathcal{F}$ 对每个拟凝聚模 $\mathcal{F}$ 都是同构，参见“代数空间的性质”命题
\ref{spaces-properties-proposition-coherator}。
由于 $Q_X$ 是左正合的（作为右伴随），我们可以考虑其右导出延拓
$$
RQ_X :
D(\mathcal{O}_X)
\longrightarrow
D(\QCoh(\mathcal{O}_X)).
$$
由于 $Q_X$ 是包含函子
$\QCoh(\mathcal{O}_X) \to \textit{Mod}(\mathcal{O}_X)$ 的右伴随，
由“导出范畴”引理
\ref{derived-lemma-derived-adjoint-functors} 可知，$RQ_X$ 是典范函子
$D(\QCoh(\mathcal{O}_X)) \to D(\mathcal{O}_X)$ 的右伴随。

\medskip\noindent
本节研究函子 $RQ_X$。在
\ref{spaces-perfect-section-better-coherator} 节中，我们将研究包含函子
$D_\QCoh(\mathcal{O}_X) \to D(\mathcal{O}_X)$ 的（密切相关的）右伴随（若其存在）。

\begin{lemma}
\label{spaces-perfect-lemma-affine-pushforward}
设 $S$ 为概形。设 $f : X \to Y$ 为 $S$ 上代数空间的仿射态射。
则 $f_*$ 定义导出函子
$f_* : D(\QCoh(\mathcal{O}_X)) \to D(\QCoh(\mathcal{O}_Y))$。
该函子具有如下性质：
$$
\xymatrix{
D(\QCoh(\mathcal{O}_X)) \ar[d]_{f_*} \ar[r] &
D_\QCoh(\mathcal{O}_X) \ar[d]^{Rf_*} \\
D(\QCoh(\mathcal{O}_Y)) \ar[r] &
D_\QCoh(\mathcal{O}_Y)
}
$$
该图交换。
\end{lemma}

\begin{proof}
函子
$f_* : \QCoh(\mathcal{O}_X) \to \QCoh(\mathcal{O}_Y)$
是正合的，参见“空间的上同调”引理
\ref{spaces-cohomology-lemma-affine-vanishing-higher-direct-images}。
因此 $f_*$ 定义导出函子
$f_* : D(\QCoh(\mathcal{O}_X)) \to D(\QCoh(\mathcal{O}_Y))$，
只需将 $f_*$ 逐项作用于任意代表复形即可；参见“导出范畴”引理
\ref{derived-lemma-right-derived-exact-functor}。
对于任意 $\mathcal{O}_X$-模复形 $\mathcal{F}^\bullet$，有典范映射
$f_*\mathcal{F}^\bullet \to Rf_*\mathcal{F}^\bullet$。
为完成证明，我们证明当 $\mathcal{F}^\bullet$ 是每一项
$\mathcal{F}^n$ 都拟凝聚的复形时，该映射是拟同构。
该断言在 $Y$ 上平展局部成立，因此不妨设 $Y$ 仿射。
由于仿射态射是可表的，利用评注
\ref{spaces-perfect-remark-match-total-direct-images} 的相容性，可将问题归结为概形情形。
概形情形是“概形的导出范畴”引理
\ref{perfect-lemma-affine-pushforward}。
\end{proof}

\begin{lemma}
\label{spaces-perfect-lemma-flat-pushforward-coherator}
设 $S$ 为概形。设 $f : X \to Y$ 为 $S$ 上代数空间的态射。
假设 $f$ 拟紧、拟分离且平坦。记
$$
\Phi : D(\QCoh(\mathcal{O}_X)) \to D(\QCoh(\mathcal{O}_Y))
$$
为函子
$f_* : \QCoh(\mathcal{O}_X) \to \QCoh(\mathcal{O}_Y)$ 的右导出函子，
则有 $RQ_Y \circ Rf_* = \Phi \circ RQ_X$。
\end{lemma}

\begin{proof}
我们将证明 $RQ_Y \circ Rf_*$ 和 $\Phi \circ RQ_X$
是同一个函子的右伴随，该函子为
$D(\QCoh(\mathcal{O}_Y)) \to D(\mathcal{O}_X)$。

\medskip\noindent
由于 $f$ 拟紧且拟分离，$f_*$ 保持拟凝聚性，参见“空间的态射”引理
\ref{spaces-morphisms-lemma-pushforward}。
回顾 $\QCoh(\mathcal{O}_X)$ 是 Grothendieck 阿贝尔范畴（“代数空间的性质”命题
\ref{spaces-properties-proposition-coherator}）。
因此，任意对象 $K$ 属于 $D(\QCoh(\mathcal{O}_X))$，都可由
K-单射复形 $\mathcal{I}^\bullet$ 表示，该复形属于 $\QCoh(\mathcal{O}_X)$，参见“单射”定理
\ref{injectives-theorem-K-injective-embedding-grothendieck}。
于是可定义 $\Phi(K) = f_*\mathcal{I}^\bullet$。

\medskip\noindent
由于 $f$ 平坦，函子 $f^*$ 正合。因此 $f^*$ 定义
$f^* : D(\mathcal{O}_Y) \to D(\mathcal{O}_X)$ 以及
$f^* : D(\QCoh(\mathcal{O}_Y)) \to D(\QCoh(\mathcal{O}_X))$。
函子 $f^* = Lf^* : D(\mathcal{O}_Y) \to D(\mathcal{O}_X)$
是 $Rf_* : D(\mathcal{O}_X) \to D(\mathcal{O}_Y)$ 的左伴随，
参见“位点上的上同调”引理 \ref{sites-cohomology-lemma-adjoint}。
同理，函子
$f^* : D(\QCoh(\mathcal{O}_Y)) \to D(\QCoh(\mathcal{O}_X))$
是 $\Phi : D(\QCoh(\mathcal{O}_X)) \to D(\QCoh(\mathcal{O}_Y))$ 的左伴随，
参见“导出范畴”引理 \ref{derived-lemma-derived-adjoint-functors}。

\medskip\noindent
设 $A$ 为 $D(\QCoh(\mathcal{O}_Y))$ 中的对象，设 $E$ 为
$D(\mathcal{O}_X)$ 中的对象。则
\begin{align*}
\Hom_{D(\QCoh(\mathcal{O}_Y))}(A, RQ_Y(Rf_*E))
& =
\Hom_{D(\mathcal{O}_Y)}(A, Rf_*E) \\
& =
\Hom_{D(\mathcal{O}_X)}(f^*A, E) \\
& =
\Hom_{D(\QCoh(\mathcal{O}_X))}(f^*A, RQ_X(E)) \\
& =
\Hom_{D(\QCoh(\mathcal{O}_Y))}(A, \Phi(RQ_X(E)))
\end{align*}
这就证明了所需结论。
\end{proof}

\begin{lemma}
\label{spaces-perfect-lemma-affine-coherator}
设 $S$ 为概形。设 $X$ 为 $S$ 上的仿射代数空间。
置 $A = \Gamma(X, \mathcal{O}_X)$。则
\begin{enumerate}
\item $Q_X : \textit{Mod}(\mathcal{O}_X) \to \QCoh(\mathcal{O}_X)$
是如下函子：它把 $\mathcal{F}$ 送到拟凝聚 $\mathcal{O}_X$-模，
该模与 $A$-模 $\Gamma(X, \mathcal{F})$ 相关联；
\item $RQ_X : D(\mathcal{O}_X) \to D(\QCoh(\mathcal{O}_X))$
是如下函子：它把 $E$ 送到拟凝聚 $\mathcal{O}_X$-模复形，
该复形与对象 $R\Gamma(X, E)$（属于 $D(A)$）相关联；
\item 在 $D_\QCoh(\mathcal{O}_X)$ 上限制时，函子 $RQ_X$ 定义
(\ref{spaces-perfect-equation-compare}) 的拟逆。
\end{enumerate}
\end{lemma}

\begin{proof}
令 $X_0 = \Spec(A)$ 为表示 $X$ 的仿射概形。
回顾存在带环位点态射
$\epsilon : X_\etale \to X_{0, Zar}$，它诱导等价
$$
\xymatrix{
\QCoh(\mathcal{O}_X) \ar@<1ex>[r]^{{\epsilon_*}} &
\QCoh(\mathcal{O}_{X_0}) \ar@<1ex>[l]^{{\epsilon^*}}
}
$$
参见引理
\ref{spaces-perfect-lemma-derived-quasi-coherent-small-etale-site}。
因此，由伴随函子的唯一性，$Q_X = \epsilon^* \circ Q_{X_0} \circ \epsilon_*$。
于是，根据“概形的导出范畴”引理
\ref{perfect-lemma-affine-coherator} 对 $Q_{X_0}$ 的描述以及
$\Gamma(X_0, \epsilon_*\mathcal{F}) = \Gamma(X, \mathcal{F})$，得到 (1)。
断言 (2) 由 (1) 以及从 $A$-模到拟凝聚
$\mathcal{O}_X$-模的函子是正合的这一事实推出。
第三个断言现在由概形情形的结果
（“概形的导出范畴”引理 \ref{perfect-lemma-affine-coherator}）
以及引理
\ref{spaces-perfect-lemma-derived-quasi-coherent-small-etale-site} 得出。
\end{proof}

\noindent
接下来，我们证明函子
$D(\QCoh(\mathcal{O}_X)) \to D_\QCoh(\mathcal{O}_X)$
何时为等价。

\begin{lemma}
\label{spaces-perfect-lemma-argument-proves}
设 $S$ 为概形。设 $X$ 为 $S$ 上拟紧且拟分离的代数空间。
假设对于每个平展态射
$j : V \to W$，其中 $W \subset X$ 为拟紧开子空间且 $V$ 为仿射，
右导出函子
$$
\Phi : D(\QCoh(\mathcal{O}_U)) \to D(\QCoh(\mathcal{O}_W))
$$ % P09-ERR-0414
（左正合函子
$j_* : \QCoh(\mathcal{O}_V) \to \QCoh(\mathcal{O}_W)$ 的右导出函子）
都包含在如下交换图中：
$$
\xymatrix{
D(\QCoh(\mathcal{O}_V)) \ar[d]_\Phi \ar[r]_{i_V} &
D_\QCoh(\mathcal{O}_V) \ar[d]^{Rj_*} \\
D(\QCoh(\mathcal{O}_W)) \ar[r]^{i_W} &
D_\QCoh(\mathcal{O}_W)
}
$$
则函子
(\ref{spaces-perfect-equation-compare})
$$
D(\QCoh(\mathcal{O}_X))
\longrightarrow
D_\QCoh(\mathcal{O}_X)
$$
是等价，并且拟逆由 $RQ_X$ 给出。
\end{lemma}

\begin{proof}
首先使用归纳原理证明 $i_X$ 是全忠实的。
更确切地说，我们将使用引理
\ref{spaces-perfect-lemma-induction-principle-enlarge}。
设 $(U \subset W, V \to W)$ 为初等特异方形，且 $V$ 为仿射、
$U, W$ 为 $X$ 中的拟紧开子空间。假设 $i_U$ 是全忠实的。
我们必须证明 $i_W$ 是全忠实的。
可以将 $X$ 替换为 $W$，即假设 $W = X$（这样做只是为了简化记号——
注意引理陈述中的条件在此操作下保持不变）。

\medskip\noindent
设 $A, B$ 为 $D(\QCoh(\mathcal{O}_X))$ 中的对象。
我们要证明
$$
\Hom_{D(\QCoh(\mathcal{O}_X))}(A, B)
\longrightarrow
\Hom_{D(\mathcal{O}_X)}(i_X(A), i_X(B))
$$
是双射。置 $T = |X| \setminus |U|$。

\medskip\noindent
先假设 $i_X(B)$ 支撑在 $T$ 上。在此情形，映射
$$
i_X(B) \to Rj_{V, *}(i_X(B)|_V) = Rj_{V, *}(i_V(B|_V))
$$
是拟同构（引理
\ref{spaces-perfect-lemma-pushforward-with-support-in-open}）。
由假设有同构
$i_X(\Phi(B|_V)) \to Rj_{V, *}(i_V(B|_V))$，其位于
$D(\mathcal{O}_X)$ 中。
此外，$\Phi$ 与 ${-}|_V$ 是拟凝聚模导出范畴上的伴随函子（参见
“导出范畴”引理 \ref{derived-lemma-derived-adjoint-functors}）。
伴随映射 $B \to \Phi(B|_V)$ 在施加 $i_X$ 后变为同构，
因此它在 $D(\QCoh(\mathcal{O}_X))$ 中是同构。
于是
\begin{align*}
\Mor_{D(\QCoh(\mathcal{O}_X))}(A, B)
& =
\Mor_{D(\QCoh(\mathcal{O}_X))}(A, \Phi(B|_V)) \\
& =
\Mor_{D(\QCoh(\mathcal{O}_V))}(A|_V, B|_V) \\
& =
\Mor_{D(\mathcal{O}_V)}(i_V(A|_V), i_V(B|_V)) \\
& =
\Mor_{D(\mathcal{O}_X)}(i_X(A), Rj_{V, *}(i_V(B|_V))) \\
& =
\Mor_{D(\mathcal{O}_X)}(i_X(A), i_X(B))
\end{align*}
如所需。这里使用了 $i_V$ 是全忠实的这一事实
（引理 \ref{spaces-perfect-lemma-affine-coherator}）。

\medskip\noindent
一般地，选取表示 B 的复形 $\mathcal{B}^\bullet$，其各项为拟凝聚
$\mathcal{O}_X$-模，表示 $B$。接着，选取复形的任意拟同构
$s : \mathcal{B}^\bullet|_U \to \mathcal{C}^\bullet$，其中复形由
$U$ 上的拟凝聚模组成。由于 $j_U : U \to X$ 拟紧且拟分离，
函子 $j_{U, *}$ 将拟凝聚模变为拟凝聚模
（空间的态射，引理 \ref{spaces-morphisms-lemma-pushforward}）。
因此有典范映射
$\mathcal{B}^\bullet \to j_{U, *}(\mathcal{B}^\bullet|_U) \to
j_{U, *}\mathcal{C}^\bullet$，它是 $X$ 上拟凝聚模复形之间的映射。
置 $B'' = j_{U, *}\mathcal{C}^\bullet$ 于
$D(\QCoh(\mathcal{O}_X))$ 中，并选取特异三角
$$
B \to B'' \to B' \to B[1]
$$
于 $D(\QCoh(\mathcal{O}_X))$ 中。
由于该三角的第一个箭头在 $U$ 上的限制是同构，故 $B'$ 支撑在 $T$ 上。
于是，在下图中
$$
\xymatrix{
\Hom_{D(\QCoh(\mathcal{O}_X))}(A, B'[-1]) \ar[r] \ar[d] &
\Hom_{D(\mathcal{O}_X)}(i_X(A), i_X(B')[-1]) \ar[d] \\
\Hom_{D(\QCoh(\mathcal{O}_X))}(A, B) \ar[r] \ar[d] &
\Hom_{D(\mathcal{O}_X)}(i_X(A), i_X(B)) \ar[d] \\
\Hom_{D(\QCoh(\mathcal{O}_X))}(A, B'') \ar[r] \ar[d] &
\Hom_{D(\mathcal{O}_X)}(i_X(A), i_X(B'')) \ar[d] \\
\Hom_{D(\QCoh(\mathcal{O}_X))}(A, B') \ar[r] &
\Hom_{D(\mathcal{O}_X)}(i_X(A), i_X(B'))
}
$$
两列是正合的，且顶端与底端的水平箭头是双射。
最后，选取拟凝聚模复形 $\mathcal{A}^\bullet$ 表示 $A$。

\medskip\noindent
设 $\alpha : i_X(A) \to i_X(B)$ 是
$D(\mathcal{O}_X)$ 中的一个态射。 % P09-ERR-0415
限制 $\alpha|_U$ 来自 $D(\QCoh(\mathcal{O}_U))$ 中的态射，
因为 $i_U$ 是全忠实的。 % P09-ERR-0416
因此，可以选择如上的
$s : \mathcal{B}^\bullet|_U \to \mathcal{C}^\bullet$，
使得 $\alpha|_U$ 由复形的实际映射
$\mathcal{A}^\bullet|_U \to \mathcal{C}^\bullet$ 表示。
这对应于复形的映射
$\mathcal{A} \to j_{U, *}\mathcal{C}^\bullet$。
换言之，$\alpha$ 在
$\Hom_{D(\mathcal{O}_X)}(i_X(A), i_X(B''))$ 中的像
来自 $\Hom_{D(\QCoh(\mathcal{O}_X))}(A, B'')$ 中的一个元素。
通过图追逐可知，$\alpha$ 来自态射 $A \to B$，
其属于 $D(\QCoh(\mathcal{O}_X))$。
最后，设 $a : A \to B$ 是
$D(\QCoh(\mathcal{O}_X))$ 中的态射，且在 $D(\mathcal{O}_X)$ 中变为零。
适当选择 $\mathcal{B}^\bullet$ 后，可以假设 $a$ 由复形的态射
$a^\bullet : \mathcal{A}^\bullet \to \mathcal{B}^\bullet$ 表示。
由于 $i_U$ 是全忠实的，限制 $a^\bullet|_U$ 在
$D(\QCoh(\mathcal{O}_U))$ 中为零。因此，可以选择 $s$，使得
$s \circ a^\bullet|_U : \mathcal{A}^\bullet|_U \to \mathcal{C}^\bullet$
同伦于零。施加函子 $j_{U, *}$ 后，我们得到
$\mathcal{A}^\bullet \to j_{U, *}\mathcal{C}^\bullet$ 同伦于零。
因此，$a$ 在
$\Hom_{D(\QCoh(\mathcal{O}_X))}(A, B'')$ 中的像为零。
于是可以假设 $a$ 是某个元素
$b \in \Hom_{D(\QCoh(\mathcal{O}_X))}(A, B'[-1])$ 的像。
$b$ 在
$\Hom_{D(\mathcal{O}_X)}(i_X(A), i_X(B')[-1])$ 中的像
来自
$\gamma \in \Hom_{D(\mathcal{O}_X)}(A, B''[-1])$。 % P09-ERR-0417
（这是因为 $a$ 在右侧的群中映射为零。）如上所见，水平箭头是满射的，
因此 $\gamma$ 来自某个 $c$，属于
$\Hom_{D(\QCoh(\mathcal{O}_X))}(A, B''[-1])$，
从而 $a = 0$，如所需。

\medskip\noindent
此时我们知道，$i_X$ 对于原来的 $X$ 是全忠实的。
由于 $RQ_X$ 是它的右伴随，我们看到
$RQ_X \circ i_X = \text{id}$（范畴，引理
\ref{categories-lemma-adjoint-fully-faithful}）。
为完成证明，我们证明对于任意 $E$（其属于
$D_\QCoh(\mathcal{O}_X)$），映射
$i_X(RQ_X(E)) \to E$ 是同构。选取特异三角
$$
i_X(RQ_X(E)) \to E \to E' \to i_X(RQ_X(E))[1]
$$
于 $D_\QCoh(\mathcal{O}_X)$ 中。
利用上述结果的形式论证表明 $i_X(RQ_X(E')) = 0$。
因此，只需证明对于 $E \in D_\QCoh(\mathcal{O}_X)$，条件
$i_X(RQ_X(E)) = 0$ 蕴含 $E = 0$。
考虑平展态射 $j : V \to X$，其中 $V$ 为仿射。根据引理
\ref{spaces-perfect-lemma-affine-coherator} 和
\ref{spaces-perfect-lemma-flat-pushforward-coherator}，并利用我们的假设，有
$$
Rj_*(E|_V) = Rj_*(i_V(RQ_V(E|_V))) = i_X(\Phi(RQ_V(E|_V))) =
i_X(RQ_X(Rj_*(E|_V)))
$$
选取特异三角
$$
E \to Rj_*(E|_V) \to E' \to E[1]
$$
对其施加 $RQ_X$，得到特异三角
$$
0 \to RQ_X(Rj_*(E|_V)) \to RQ_X(E') \to 0[1]
$$
换言之，中间的映射是同构。
结合上面的等式串，我们发现第一个特异三角变成特异三角
$$
E \to i_X(RQ_X(E')) \to E' \to E[1]
$$
其中中间态射是伴随映射。然而，复合映射
$E \to E'$ 为零，因此由伴随性 $E \to i_X(RQ_X(E'))$ 为零！
由于该态射同构于映射
$E \to Rj_*(E|_V)$，后者是
$\text{id} : E|_V \to E|_V$ 的伴随，我们得到 $E|_V$ 为零。
这对所有平展的仿射 $V$（其位于 $X$ 上）都成立，故 $E$ 为零，如所需。
\end{proof}

\begin{proposition}
\label{spaces-perfect-proposition-quasi-compact-affine-diagonal}
设 $S$ 为概形。设 $X$ 为 $S$ 上的拟紧代数空间，且在
$\mathbf{Z}$ 上具有仿射对角态射（如“代数空间的性质”定义
\ref{spaces-properties-definition-separated} 中所述）。则函子
(\ref{spaces-perfect-equation-compare})
$$
D(\QCoh(\mathcal{O}_X))
\longrightarrow
D_\QCoh(\mathcal{O}_X)
$$
是等价，并且拟逆由 $RQ_X$ 给出。
\end{proposition}

\begin{proof}
设 $V \to W$ 为平展态射，其中 $V$ 为仿射，$W$ 为 $X$ 的拟紧开子空间。
则态射 $V \to W$ 是仿射的，因为 $W$ 在 $\mathbf{Z}$ 上具有仿射对角态射，
且 $V$ 为仿射（空间的态射，引理
\ref{spaces-morphisms-lemma-affine-permanence}）。
于是引理 \ref{spaces-perfect-lemma-affine-pushforward} 保证
引理 \ref{spaces-perfect-lemma-argument-proves} 的假设成立。
因此结论成立。
\end{proof}

\begin{lemma}
\label{spaces-perfect-lemma-direct-image-coherator}
设 $S$ 为概形，设 $f : X \to Y$ 为 $S$ 上代数空间的态射。
假设 $X$ 和 $Y$ 都拟紧，并在 $\mathbf{Z}$ 上具有仿射对角态射
（如“代数空间的性质”定义
\ref{spaces-properties-definition-separated} 中所述）。记
$$
\Phi : D(\QCoh(\mathcal{O}_X)) \to D(\QCoh(\mathcal{O}_Y))
$$
为函子
$f_* : \QCoh(\mathcal{O}_X) \to \QCoh(\mathcal{O}_Y)$ 的右导出函子，
则下图
$$
\xymatrix{
D(\QCoh(\mathcal{O}_X)) \ar[d]_\Phi \ar[r] &
D_\QCoh(\mathcal{O}_X) \ar[d]^{Rf_*} \\
D(\QCoh(\mathcal{O}_Y)) \ar[r] &
D_\QCoh(\mathcal{O}_Y)
}
$$
是交换的。 % P09-ERR-0418
\end{lemma}

\begin{proof}
注意，由命题
\ref{spaces-perfect-proposition-quasi-compact-affine-diagonal}，图中的水平箭头
都是范畴等价。因此可以将这些范畴等同（对于其他具有仿射对角态射的
拟紧代数空间也同样处理），于是引理的陈述变为：典范映射
$\Phi(K) \to Rf_*(K)$ 对所有对象 $K$ 属于
$D(\QCoh(\mathcal{O}_X))$ 都是同构。
注意，若
$K_1 \to K_2 \to K_3 \to K_1[1]$
是 $D(\QCoh(\mathcal{O}_X))$ 中的特异三角，且该陈述对其中两个成立，
则对第三个也成立。

\medskip\noindent
令 $\mathcal{B} \subset \Ob(X_{spaces, \etale})$ 为由如下对象组成的集合：
这些对象拟紧且具有仿射对角态射。
对 $U \in \mathcal{B}$ 以及任意态射 $g : U \to Z$，其中 $Z$ 是
在 $S$ 上具有仿射对角态射的拟紧代数空间，记
$$
\Phi_g : D(\QCoh(\mathcal{O}_U)) \to D(\QCoh(\mathcal{O}_Z))
$$
为 $g_*$ 的导出延拓。令
$P(U) =$ “对于任意 $K$ 属于 $D(\QCoh(\mathcal{O}_U))$，以及上述任意
$g : U \to Z$，映射 $\Phi_g(K) \to Rg_*K$ 是同构”。
根据评注 \ref{spaces-perfect-remark-how-to}，归纳原理引理
\ref{spaces-perfect-lemma-induction-principle-separated} 的条件 (1)、(2) 和 (3)(a)
成立，因此还需证明 (3)(b) 和 (4)。

\medskip\noindent
检验条件 (3)(b)。设 $U$ 为在 $X$ 上平展的仿射概形。
设 $g : U \to Z$ 如上。由于 $Z$ 的对角态射是仿射的，
态射 $g : U \to Z$ 是仿射的（空间的态射，引理
\ref{spaces-morphisms-lemma-affine-permanence}）。
因此，由引理 \ref{spaces-perfect-lemma-affine-pushforward}，$P(U)$ 成立。

\medskip\noindent
检验条件 (4)。设 $(U \subset W, V \to W)$ 为
$X_{spaces, \etale}$ 中的初等特异方形，其中 $U, W, V$ 属于
$\mathcal{B}$ 且 $V$ 为仿射。假设 $P$ 对 $U$、$V$ 以及
$U \times_W V$ 成立。
我们必须证明 $P$ 对 $W$ 成立。设 $g : W \to Z$ 为到一个
具有仿射对角态射的拟紧代数空间的态射。
设 $K$ 为 $D(\QCoh(\mathcal{O}_W))$ 中的对象。考虑特异三角
$$
K \to Rj_{U, *}K|_U \oplus Rj_{V, *}K|_V \to
Rj_{U \times_W V, *}K|_{U \times_W V} \to K[1]
$$
于 $D(\mathcal{O}_W)$ 中。由上面提到的两出三性质，
只需证明
$\Phi_g(Rj_{U, *}K|_U) \to Rg_*(Rj_{U, *}K|_U)$
是同构，并且对 $V$ 和 $U \times_W V$ 也同样成立。
这将在下一段中讨论。

\medskip\noindent
设 $j : U \to W$ 为 $X_{spaces, \etale}$ 中的态射，其中
$U, W$ 属于 $\mathcal{B}$，且 $P$ 对 $U$ 成立。 % P09-ERR-0419
设 $g : W \to Z$ 为到一个具有仿射对角态射的拟紧代数空间的态射。
为完成证明，必须证明
$\Phi_g(Rj_*K) \to Rg_*(Rj_*K)$
对于任意 $K$ 属于 $D(\QCoh(\mathcal{O}_U))$ 都是同构。
设 $\mathcal{I}^\bullet$ 为 $\QCoh(\mathcal{O}_U)$ 中表示 $K$ 的
K-单射复形。
将 $P(U)$ 应用于 $j$，可知
$j_*\mathcal{I}^\bullet$ 表示 $Rj_*K$。
由于
$j_* : \QCoh(\mathcal{O}_U) \to \QCoh(\mathcal{O}_X)$
具有正合左伴随
$j^* : \QCoh(\mathcal{O}_X) \to \QCoh(\mathcal{O}_U)$， % P09-ERR-0420
可知 $j_*\mathcal{I}^\bullet$ 是
$\QCoh(\mathcal{O}_W)$ 中的 K-单射复形；参见“导出范畴”引理
\ref{derived-lemma-adjoint-preserve-K-injectives}。
因此 $\Phi_g(Rj_*K)$ 由
$g_*j_*\mathcal{I}^\bullet = (g \circ j)_*\mathcal{I}^\bullet$ 表示。
将 $P(U)$ 应用于 $g \circ j$，可知它表示
$R_{g \circ j, *}(K) = Rg_*(Rj_*K)$。证明完成。
\end{proof}

\section{诺特空间上的凝聚化函子}
\label{spaces-perfect-section-coherator-Noetherian}

\noindent
为处理诺特代数空间的情形，我们还需要一些关于单射模的事实。

\begin{lemma}
\label{spaces-perfect-lemma-affine-injective-colimit-direct-sum-pushforwards-artin}
设 $S$ 为诺特仿射概形。$\QCoh(\mathcal{O}_S)$ 的每个单射对象都是
拟凝聚层的滤余极限 $\colim_i \mathcal{F}_i$，其中
$$
\mathcal{F}_i = (Z_i \to S)_*\mathcal{G}_i
$$
且 $Z_i$ 是某个阿廷环的谱，$\mathcal{G}_i$ 是 $Z_i$ 上的相干模。
\end{lemma}

\begin{proof}
令 $S = \Spec(A)$。设 $\mathcal{J}$ 是
$\QCoh(\mathcal{O}_S)$ 中的单射对象。由于 $\QCoh(\mathcal{O}_S)$
等价于 $A$-模的范畴，我们看到 $\mathcal{J}$ 等于某个单射模所对应的层
$\widetilde{J}$，其来源是某个单射 $A$-模 $J$。
根据《对偶化复形》命题
\ref{dualizing-proposition-structure-injectives-noetherian}，我们可以写成
$J = \bigoplus E_\alpha$，其中 $E_\alpha$ 不可分解，
因而同构于某个点处的 reside 域的单射包。 % P09-ERR-0421
因此（因为阿廷概形的有限不交并仍是阿廷的），不妨设 $J$ 是
某个素理想处 $\kappa(\mathfrak p)$ 的单射包，其中该素理想为
$\mathfrak p$，且属于 $A$。
于是 $J = \bigcup J[\mathfrak p^n]$，其中 $J[\mathfrak p^n]$ 是
$\kappa(\mathfrak p)$ 在 $A_\mathfrak/\mathfrak p^nA_\mathfrak p$ 上的单射包；见
《对偶化复形》引理 \ref{dualizing-lemma-union-artinian}。 % P09-ERR-0422
因此 $\widetilde{J}$ 是下列层的余极限：
$(Z_n \to X)_*\mathcal{G}_n$，其中 % P09-ERR-0423
$Z_n = \Spec(A_\mathfrak p/\mathfrak p^nA_\mathfrak p)$，且
$\mathfrak G_n$ 是与有限的
$A_\mathfrak/\mathfrak p^nA_\mathfrak p$-模 $J[\mathfrak p^n]$ 关联的凝聚层。
\par % P09-ERR-0424
有限性由《对偶化复形》引理
\ref{dualizing-lemma-finite} 得到。
\end{proof}

\begin{lemma}
\label{spaces-perfect-lemma-injective-colimit-direct-sum-pushforwards-artin}
设 $S$ 为仿射概形，$X$ 为 $S$ 上的诺特代数空间。$\QCoh(\mathcal{O}_X)$
中的每个单射对象都是滤余极限 $\colim_i \mathcal{F}_i$ 的直接因子，
其中这些层具有如下拟凝聚层的形式：
$$
\mathcal{F}_i = (Z_i \to X)_*\mathcal{G}_i
$$
其中 $Z_i$ 是某个阿廷环的谱，$\mathcal{G}_i$ 是 $Z_i$ 上的凝聚模。
\end{lemma}

\begin{proof}
取一个仿射概形 $U$ 以及满射平展态射 $j : U \to X$（《空间的性质》引理
\ref{spaces-properties-lemma-quasi-compact-affine-cover}）。
于是 $U$ 是诺特仿射概形。取一个单射对象 $\mathcal{J}'$，它属于
$\QCoh(\mathcal{O}_U)$，使得存在一个单射 $\mathcal{J}|_U \to \mathcal{J}'$。于是
$$
\mathcal{J} \to j_*\mathcal{J}'
$$
是 $\QCoh(\mathcal{O}_X)$ 中的单射态射，
从而将 $\mathcal{J}$ 识别为 $j_*\mathcal{J}'$ 的一个直接因子。
因此，结论由关于 $\mathcal{J}'$ 的相应结果推出；该结果已在引理
\ref{spaces-perfect-lemma-affine-injective-colimit-direct-sum-pushforwards-artin} 中证明。
\end{proof}

\begin{lemma}
\label{spaces-perfect-lemma-flat-pullback-injective-quasi-coherent}
设 $S$ 为概形。设 $f : X \to Y$ 为 $S$ 上代数空间之间的平坦、拟紧且
拟分离态射。若 $\mathcal{J}$ 是 $\QCoh(\mathcal{O}_X)$ 中的单射对象，
则 $f_*\mathcal{J}$ 是 $\QCoh(\mathcal{O}_Y)$ 中的单射对象。
\end{lemma}

\begin{proof}
由于 $f$ 拟紧且拟分离，函子 $f_*$ 将拟凝聚层变为拟凝聚层
（《空间的态射》引理 \ref{spaces-morphisms-lemma-pushforward}）。
函子 $f^*$ 是 $f_*$ 的左伴随，并且将单射变为单射。
因此结论由《同调》引理
\ref{homology-lemma-adjoint-preserve-injectives} 得到 % P09-ERR-0425
\end{proof}

\begin{lemma}
\label{spaces-perfect-lemma-injective-pushforward}
设 $S$ 为概形。设 $X$ 为 $S$ 上的诺特代数空间。若
$\mathcal{J}$ 是 $\QCoh(\mathcal{O}_X)$ 中的单射对象，则
\begin{enumerate}
\item 有 $H^p(U, \mathcal{J}|_U) = 0$，其中 $p > 0$，且 $U$ 是
在 $X$ 上平展的拟紧拟分离代数空间；
\item 对于代数空间之间的任意 $f : X \to Y$（均为 $S$ 上的代数空间），
若 $Y$ 拟分离，则有 $R^pf_*\mathcal{J} = 0$（$p > 0$）。
\end{enumerate}
\end{lemma}

\begin{proof}
第 (1) 部分的证明。将 $\mathcal{J}$ 写成
$\colim \mathcal{F}_i$ 的一个直接因子，其中
$\mathcal{F}_i = (Z_i \to X)_*\mathcal{G}_i$，
如引理 \ref{spaces-perfect-lemma-injective-colimit-direct-sum-pushforwards-artin} 所述。
显然，只需对 $\colim \mathcal{F}_i$ 证明消灭即可。
由于拉回与余极限交换，且 $U$ 拟紧拟分离，只需证明
$H^p(U, \mathcal{F}_i|_U) = 0$（$p > 0$）；见
《空间的上同调》引理 \ref{spaces-cohomology-lemma-colimits}。
注意 $Z_i \to X$ 是仿射态射；见《空间的态射》引理
\ref{spaces-morphisms-lemma-Artinian-affine}。因此
$$
\mathcal{F}_i|_U = (Z_i \times_X U \to U)_*\mathcal{G}'_i =
R(Z_i \times_X U \to U)_*\mathcal{G}'_i
$$
其中 $\mathcal{G}'_i$ 是 $\mathcal{G}_i$ 拉回到
$Z_i \times_X U$ 所得的层；见《空间的上同调》引理
\ref{spaces-cohomology-lemma-affine-base-change}。
由于 $Z_i \times_X U$ 是仿射的，我们得到
$\mathcal{G}'_i$ 在 $Z_i \times_X U$ 上没有高阶上同调。
由 Leray 谱序列可知，对 $\mathcal{F}_i|_U$ 也有同样结论
（《位点上的上同调》引理
\ref{sites-cohomology-lemma-apply-Leray}）。

\medskip\noindent
第 (2) 部分的证明。设 $f : X \to Y$ 为 $S$ 上代数空间之间的态射。
设 $V \to Y$ 为一个平展态射，且 $V$ 仿射。则
$V \times_Y X \to X$ 是平展态射，并且 $V \times_Y X$ 是一个在 $X$ 上
平展的拟紧拟分离代数空间（细节略）。因此，由第 (1) 部分
$H^p(V \times_Y X, \mathcal{J})$ 为零。
由于 $R^pf_*\mathcal{J}$ 是下述预层所关联的层：
$V \mapsto H^p(V \times_Y X, \mathcal{J})$，结论得证。
\end{proof}

\begin{lemma}
\label{spaces-perfect-lemma-Noetherian-pushforward}
设 $S$ 为概形。
设 $f : X \to Y$ 为 $S$ 上诺特代数空间之间的态射。
则拟凝聚层上的 $f_*$ 具有右导出延拓
$\Phi : D(\QCoh(\mathcal{O}_X)) \to D(\QCoh(\mathcal{O}_Y))$，
使得下图交换：
$$
\xymatrix{
D(\QCoh(\mathcal{O}_X)) \ar[d]_{\Phi} \ar[r] &
D_\QCoh(\mathcal{O}_X) \ar[d]^{Rf_*} \\
D(\QCoh(\mathcal{O}_Y)) \ar[r] &
D_\QCoh(\mathcal{O}_Y)
}
$$
\end{lemma}

\begin{proof}
由于 $X$ 和 $Y$ 是诺特的，该态射是拟紧且拟分离的（见
《空间的态射》引理
\ref{spaces-morphisms-lemma-quasi-compact-quasi-separated-permanence}）。
因此 $f_*$ 保持拟凝聚性；见《空间的态射》引理
\ref{spaces-morphisms-lemma-pushforward}。 % P09-ERR-0426
接下来，设 $K$ 是 $D(\QCoh(\mathcal{O}_X))$ 中的对象。
由于 $\QCoh(\mathcal{O}_X)$ 是 Grothendieck 阿贝尔范畴
（《空间的性质》命题 \ref{spaces-properties-proposition-coherator}），
我们可以将 $K$ 表示为一个 K-单射复形 $\mathcal{I}^\bullet$，使得每个
$\mathcal{I}^n$ 都是 $\QCoh(\mathcal{O}_X)$ 中的单射对象；见
《单射》定理 \ref{injectives-theorem-K-injective-embedding-grothendieck}。
因此，函子 $\Phi$ 定义为
$$
\Phi(K) = f_*\mathcal{I}^\bullet
$$
其中右端被视为 $D(\QCoh(\mathcal{O}_Y))$ 中的对象。为了完成引理的证明，
只需证明典范映射
$$
f_*\mathcal{I}^\bullet \longrightarrow Rf_*\mathcal{I}^\bullet
$$
是 $D(\mathcal{O}_Y)$ 中的同构。为此，只需证明该映射在上同调层上诱导同构。
任取 $m \in \mathbf{Z}$。令 $N = N(X, Y, f)$ 如引理
\ref{spaces-perfect-lemma-quasi-coherence-direct-image} 中所述。
考虑如下拟凝聚层复形的短正合序列
$$
0 \to \sigma_{\geq m - N - 1}\mathcal{I}^\bullet \to
\mathcal{I}^\bullet \to \sigma_{\leq m - N - 2}\mathcal{I}^\bullet \to 0
$$
这是 $X$ 上的复形。由
由引理 \ref{spaces-perfect-lemma-quasi-coherence-direct-image}，我们看到
$Rf_*\sigma_{\leq m - N - 2}\mathcal{I}^\bullet$ 的上同调层在
$\geq m - 1$ 的次数中为零。因此
$R^mf_*\mathcal{I}^\bullet$ 同构于
$R^mf_*\sigma_{\geq m - N - 1}\mathcal{I}^\bullet$。
换言之，我们可以假设 $\mathcal{I}^\bullet$ 是
$\QCoh(\mathcal{O}_X)$ 中单射对象构成的有下界复形。
此情形由 Leray 无环性引理（《导出范畴》引理
\ref{derived-lemma-leray-acyclicity}）以及引理
\ref{spaces-perfect-lemma-injective-pushforward} 所给出的必要消灭性推出。 % P09-ERR-0427
\end{proof}

\begin{proposition}
\label{spaces-perfect-proposition-Noetherian}
设 $S$ 为概形，$X$ 为 $S$ 上的诺特代数空间。
则函子（\ref{spaces-perfect-equation-compare}）
$$
D(\QCoh(\mathcal{O}_X))
\longrightarrow
D_\QCoh(\mathcal{O}_X)
$$
是一个等价，并且准逆由 $RQ_X$ 给出。
\end{proposition}

\begin{proof}
这立即由引理 \ref{spaces-perfect-lemma-Noetherian-pushforward} 和
\ref{spaces-perfect-lemma-argument-proves} 得到。
\end{proof}

\section{拟凝聚复形与完美复形}
\label{spaces-perfect-section-spell-out}

\noindent
本节研究《位点上的上同调》第
\ref{sites-cohomology-section-strictly-perfect}、
\ref{sites-cohomology-section-pseudo-coherent}、
\ref{sites-cohomology-section-tor} 和
\ref{sites-cohomology-section-perfect} 节中为代数空间的平展位点
定义的一般概念。特别地，我们将其与概形情形相匹配。

\medskip\noindent
首先比较概形上的拟凝聚复形与其相应小平展位点上的拟凝聚复形这一概念。

\begin{lemma}
\label{spaces-perfect-lemma-descend-finite-type}
设 $X$ 为概形，$\mathcal{F}$ 为 $\mathcal{O}_X$-模。
以下条件等价：
\begin{enumerate}
\item $\mathcal{F}$ 作为 $\mathcal{O}_X$-模是有限型的；
\item $\epsilon^*\mathcal{F}$ 作为小平展位点上的
$\mathcal{O}_\etale$-模是有限型的，且该位点属于 $X$。
\end{enumerate}
这里 $\epsilon$ 如（\ref{spaces-perfect-equation-epsilon}）中所述。
\end{lemma}

\begin{proof}
(1) $\Rightarrow$ (2) 的蕴含是一般事实，见《位点上的模》引理
\ref{sites-modules-lemma-local-pullback}。
假设 (2)。按假设，存在一个平展覆盖
$\{f_i : X_i \to X\}$，使得
$\epsilon^*\mathcal{F}|_{(X_i)_\etale}$ 由有限多个截面生成。
任取 $x \in X$。我们将证明 $\mathcal{F}$
在 $x$ 的某个邻域中由有限多个截面生成。
设 $x$ 在 $X_i \to X$ 的像中，并记 $X' = X_i$。令
$s_1, \ldots, s_n \in
\Gamma(X', \epsilon^*\mathcal{F}|_{X'_\etale})$
为生成截面。由于
$\epsilon^*\mathcal{F} =
\epsilon^{-1}\mathcal{F} \otimes_{\epsilon^{-1}\mathcal{O}_X}
\mathcal{O}_\etale$
，可以找到一个平展态射 $X'' \to X'$，使得 $x$
在 $X' \to X$ 的像中，并且
$s_i|_{X''} = \sum s_{ij} \otimes a_{ij}$，其中
$s_{ij} \in \epsilon^{-1}\mathcal{F}(X'')$，且
$a_{ij} \in \mathcal{O}_\etale(X'')$。记 $U \subset X$ 为
$X'' \to X$ 的像。这是一个开子概形，因为 $f'' : X'' \to X$ 是平展的
（《态射》引理 \ref{morphisms-lemma-etale-open}）。 % P09-ERR-0428
再将 $X''$ 缩小一些后，可假设 $s_{ij}$ 来自
$t_{ij} \in \mathcal{F}(U)$，这由逆像函子 $\epsilon^{-1}$ 的构造得到。
现在断言 $t_{ij}$ 生成 $\mathcal{F}|_U$，从而完成引理的证明。事实上，
相应的映射
$\mathcal{O}_U^{\oplus N} \to \mathcal{F}|_U$
具有如下性质：其经 $f''$ 拉回到 $X''$ 后是满射。由于
$f'' : X'' \to U$ 是概形的满射平坦态射，这意味着
$\mathcal{O}_U^{\oplus N} \to \mathcal{F}|_U$ 由考察茎并利用
$\mathcal{O}_{U, f''(z)} \to \mathcal{O}_{X'', z}$
对所有 $z \in X''$ 都忠实平坦而为满射。
\end{proof}

\noindent
在上述情形中，位点态射 $\epsilon$ 是平坦的，因而在模复形上定义拉回。

\begin{lemma}
\label{spaces-perfect-lemma-descend-pseudo-coherent}
设 $X$ 为概形，$E$ 为 $D(\mathcal{O}_X)$ 中的对象。
以下条件等价：
\begin{enumerate}
\item $E$ 是 $m$-拟凝聚的；
\item $\epsilon^*E$ 是 $m$-拟凝聚的，位于 $X$ 的小平展位点上。
\end{enumerate}
这里 $\epsilon$ 如（\ref{spaces-perfect-equation-epsilon}）中所述。
\end{lemma}

\begin{proof}
(1) $\Rightarrow$ (2) 的蕴含是一般事实，见《位点上的上同调》引理
\ref{sites-cohomology-lemma-pseudo-coherent-pullback}。
假设 $\epsilon^*E$ 是 $m$-拟凝聚的。
以下不再特别说明：$\epsilon^*$ 是正合函子，因此
$$
\epsilon^*H^i(E) = H^i(\epsilon^*E).
$$
为证明 $E$ 是 $m$-拟凝聚的，我们可以在 $X$ 上局部工作，
因而不妨设 $X$ 拟紧（例如为仿射的）。
由于 $X$ 拟紧，每个平展覆盖 $\{U_i \to X\}$
都有有限细化。因此 $\epsilon^*E$ 是
$D^{-}(\mathcal{O}_\etale)$ 中的对象；见《位点上的上同调》定义
\ref{sites-cohomology-definition-pseudo-coherent} 后的评注。
由引理 \ref{spaces-perfect-lemma-epsilon-flat} 可知 $E$ 是
$D^-(\mathcal{O}_X)$ 中的对象。

\medskip\noindent
令 $n \in \mathbf{Z}$ 为使得 $H^n(E)$ 非零的最大整数；则 $n$ 也是使得
$H^n(\epsilon^*E)$ 非零的最大整数。
我们将对 $n - m$ 作归纳来证明引理。
若 $n < m$，则引理显然成立。
若 $n \geq m$，则 $H^n(\epsilon^*E)$ 是有限的
$\mathcal{O}_\etale$-模；见《位点上的上同调》引理
\ref{sites-cohomology-lemma-finite-cohomology}。
于是 $H^n(E)$ 是有限的 $\mathcal{O}_X$-模；见引理
\ref{spaces-perfect-lemma-descend-finite-type}。
将 $X$ 换成一个开覆盖的成员后，不妨设存在满射
$\mathcal{O}_X^{\oplus t} \to H^n(E)$。
在 $X$ 上局部地可将其提升为复形映射
$\alpha : \mathcal{O}_X^{\oplus t}[-n] \to E$（细节略）。
取一个特异三角
$$
\mathcal{O}_X^{\oplus t}[-n] \to E \to C \to \mathcal{O}_X^{\oplus t}[-n + 1]
$$
则 $C$ 在次数 $\geq n$ 上的上同调消失。另一方面，复形
$\epsilon^*C$ 是 $m$-拟凝聚的，见《位点上的上同调》引理
\ref{sites-cohomology-lemma-cone-pseudo-coherent}。
因此由归纳假设，$C$ 是 $m$-拟凝聚的。再次应用《位点上的上同调》引理
\ref{sites-cohomology-lemma-cone-pseudo-coherent}，即可得出结论。
\end{proof}

\begin{lemma}
\label{spaces-perfect-lemma-descend-tor-amplitude}
设 $X$ 为概形，$E$ 为 $D(\mathcal{O}_X)$ 中的对象。则
\begin{enumerate}
\item $E$ 具有 $[a, b]$ 中的 Tor 振幅，当且仅当
$\epsilon^*E$ 具有 $[a, b]$ 中的 Tor 振幅；
\item $E$ 具有有限 Tor 维数，当且仅当 $\epsilon^*E$ 具有有限
Tor 维数。
\end{enumerate}
这里 $\epsilon$ 如（\ref{spaces-perfect-equation-epsilon}）中所述。
\end{lemma}

\begin{proof}
容易的蕴含由《位点上的上同调》引理
\ref{sites-cohomology-lemma-tor-amplitude-pullback} 得到。
反之，假设 $\epsilon^*E$ 具有 $[a, b]$ 中的 Tor 振幅。
设 $\mathcal{F}$ 为 $\mathcal{O}_X$-模。由于 $\epsilon$ 是带环位点的平坦
态射（引理 \ref{spaces-perfect-lemma-epsilon-flat}），有
$$
\epsilon^*(E \otimes^\mathbf{L}_{\mathcal{O}_X} \mathcal{F})
=
\epsilon^*E
\otimes^\mathbf{L}_{\mathcal{O}_\etale}
\epsilon^*\mathcal{F}
$$
因此，右端（按假设）上同调层的消灭性，通过引理
\ref{spaces-perfect-lemma-epsilon-flat}，推出
$E \otimes^\mathbf{L}_{\mathcal{O}_X} \mathcal{F}$ 的上同调层具有所需的消灭性。
\end{proof}

\begin{lemma}
\label{spaces-perfect-lemma-tor-dimension-rel}
设 $f : X \to Y$ 为概形的态射，$E$ 为 $D(\mathcal{O}_X)$ 中的对象。则
\begin{enumerate}
\item $E$ 作为 $D(f^{-1}\mathcal{O}_Y)$ 中的对象具有 $[a, b]$ 中的
Tor 振幅，当且仅当 $\epsilon^*E$ 具有 $[a, b]$ 中的 Tor 振幅，作为
$D(f_{small}^{-1}\mathcal{O}_{Y_\etale})$ 中的对象；
\item $E$ 作为 $D(f^{-1}\mathcal{O}_Y)$ 中的对象局部具有有限 Tor 维数，
当且仅当 $\epsilon^*E$ 作为
$D(f_{small}^{-1}\mathcal{O}_{Y_\etale})$ 中的对象局部具有有限 Tor 维数。
\end{enumerate}
这里 $\epsilon$ 如（\ref{spaces-perfect-equation-epsilon}）中所述。
\end{lemma}

\begin{proof}
(1) 中容易的方向由《位点上的上同调》引理
\ref{sites-cohomology-lemma-tor-amplitude-pullback} 得到。
设 $x \in X$ 为一点，设 $\overline{x}$ 为几何点，其位于 $x$ 之上。
令 $y = f(x)$，并记 $\overline{y}$ 为 $Y$ 中由 $\overline{x}$ 给出的几何点。
则 $(f^{-1}\mathcal{O}_Y)_x = \mathcal{O}_{Y, y}$
（《层》引理 \ref{sheaves-lemma-stalk-pullback}），并且
$$
(f_{small}^{-1}\mathcal{O}_{Y_\etale})_{\overline{x}} =
\mathcal{O}_{Y_\etale, \overline{y}} =
\mathcal{O}_{Y, y}^{sh}
$$
是严格 Hensel 化
（根据《平展上同调》引理
\ref{etale-cohomology-lemma-stalk-pullback}
和 \ref{etale-cohomology-lemma-describe-etale-local-ring}）。
由于 $\mathcal{O}_{X_\etale}$ 在 $X$ 处的茎为
$\mathcal{O}_{X, x}^{sh}$，由拉回的传递性得到 % P09-ERR-0429
$$
(\epsilon^*E)_{\overline{x}} =
E_x \otimes_{\mathcal{O}_{X, x}}^\mathbf{L} \mathcal{O}_{X, x}^{sh}
$$
若 $\epsilon^*E$ 具有 $[a, b]$ 中的 Tor 振幅，作为
$f_{small}^{-1}\mathcal{O}_{Y_\etale}$-复形，则 $(\epsilon^*E)_{\overline{x}}$
具有 $[a, b]$ 中的 Tor 振幅，作为 $\mathcal{O}_{Y, y}^{sh}$-模复形
（因为取茎就是拉回，并使用上述引理）。由《平坦性补充》引理
\ref{flat-lemma-tor-amplitude-up-down-henselization}，可知
$(\epsilon^*E)_{\overline{x}}$
作为 $\mathcal{O}_{Y, y}$-模复形的 Tor 振幅属于 $[a, b]$。
由于 $\mathcal{O}_{X, x} \to \mathcal{O}_{X, x}^{sh}$ 是忠实平坦的
（《代数补充》引理
\ref{more-algebra-lemma-dumb-properties-henselization}），且
$(\epsilon^*E)_{\overline{x}} =
E_x \otimes_{\mathcal{O}_{X, x}}^\mathbf{L} \mathcal{O}_{X, x}^{sh}$，
我们可以应用《代数补充》引理
\ref{more-algebra-lemma-no-change-tor-amplitude}，从而得到
$E_x$ 作为 $\mathcal{O}_{Y, y}$-模复形的 Tor 振幅属于 $[a, b]$。
由《上同调》引理 \ref{cohomology-lemma-tor-amplitude-stalk}，得到
$E$ 作为 $D(f^{-1}\mathcal{O}_Y)$ 中的对象具有 $[a, b]$ 中的 Tor 振幅。
这给出了 (1) 的逆向蕴含。(2) 正式地由 (1) 得到。
\end{proof}

\begin{lemma}
\label{spaces-perfect-lemma-descend-perfect}
设 $X$ 为概形，$E$ 为 $D(\mathcal{O}_X)$ 中的对象。
则 $E$ 是 $D(\mathcal{O}_X)$ 中的完美对象，当且仅当
$\epsilon^*E$ 是 $D(\mathcal{O}_\etale)$ 中的完美对象。
这里 $\epsilon$ 如（\ref{spaces-perfect-equation-epsilon}）中所述。
\end{lemma}

\begin{proof}
容易的蕴含由《位点上的上同调》引理
\ref{sites-cohomology-lemma-perfect-pullback} 中的一般结果得到。
反之，可以使用《位点上的上同调》引理
\ref{sites-cohomology-lemma-perfect} 的等价刻画，
以及关于拟凝聚复形和有限 Tor 维数复形的相应结果，即引理
\ref{spaces-perfect-lemma-descend-pseudo-coherent} 和
\ref{spaces-perfect-lemma-descend-tor-amplitude}。
略去一些细节。
\end{proof}

\begin{lemma}
\label{spaces-perfect-lemma-pseudo-coherent}
设 $S$ 为概形，$X$ 为 $S$ 上的代数空间。
若 $E$ 是 $m$-拟凝聚对象，属于 $D(\mathcal{O}_X)$，
则 $H^i(E)$ 是拟凝聚 $\mathcal{O}_X$-模，其中 $i > m$。
若 $E$ 是拟凝聚的，则 $E$ 是
$D_\QCoh(\mathcal{O}_X)$ 中的对象。
\end{lemma}

\begin{proof}
局部地，$H^i(E)$ 同构于 $H^i(\mathcal{E}^\bullet)$，
其中 $\mathcal{E}^\bullet$ 严格完美。层 $\mathcal{E}^i$
是有限自由模的直接因子，因而拟凝聚。引理得证。
\end{proof}

\begin{lemma}
\label{spaces-perfect-lemma-identify-pseudo-coherent-noetherian}
设 $S$ 为概形，$X$ 为 $S$ 上的诺特代数空间。
设 $E$ 为 $D_\QCoh(\mathcal{O}_X)$ 中的对象。对于
$m \in \mathbf{Z}$，以下条件等价：
\begin{enumerate}
\item $H^i(E)$ 对于 $i \geq m$ 是凝聚的，并且当 $i \gg 0$ 时为零；
\item $E$ 是 $m$-拟凝聚的。
\end{enumerate}
特别地，$E$ 是拟凝聚的，当且仅当 $E$ 是
$D^-_{\textit{Coh}}(\mathcal{O}_X)$ 中的对象。
\end{lemma}

\begin{proof}
由于 $X$ 拟紧，可以找到一个仿射概形 $U$ 以及满射平展态射
$U \to X$（《空间的性质》引理
\ref{spaces-properties-lemma-quasi-compact-affine-cover}）。
注意 $U$ 是诺特的。
注意，$E$ 是 $m$-拟凝聚的，当且仅当 $E|_U$ 是
$m$-拟凝聚的（这由定义或《位点上的上同调》引理
\ref{sites-cohomology-lemma-pseudo-coherent-independent-representative} 推出）。
同样，$H^i(E)$ 是凝聚的，当且仅当
$H^i(E)|_U = H^i(E|_U)$ 是凝聚的（见《空间的上同调》引理
\ref{spaces-cohomology-lemma-coherent-Noetherian}）。
因此不妨设 $X$ 可表。

\medskip\noindent
若 $X$ 可由概形 $X_0$ 表示，则（引理
\ref{spaces-perfect-lemma-derived-quasi-coherent-small-etale-site}）
可写 $E = \epsilon^*E_0$，其中 $E_0$ 是
$D_\QCoh(\mathcal{O}_{X_0})$ 中的对象，且
$\epsilon : X_\etale \to (X_0)_{Zar}$ 如
（\ref{spaces-perfect-equation-epsilon}）中所述。
在此情形中，由引理 \ref{spaces-perfect-lemma-descend-pseudo-coherent}，
$E$ 是 $m$-拟凝聚的，当且仅当 $E_0$ 也是如此。
同样，由引理 \ref{spaces-perfect-lemma-descend-finite-type}，$H^i(E_0)$
是有限型的（即凝聚的），当且仅当 $H^i(E)$ 也是如此。
最后，由引理 \ref{spaces-perfect-lemma-epsilon-flat}，$H^i(E_0) = 0$
当且仅当 $H^i(E) = 0$。
因此归约到概形情形，即《概形的导出范畴》引理
\ref{perfect-lemma-identify-pseudo-coherent-noetherian}。
\end{proof}

\begin{lemma}
\label{spaces-perfect-lemma-tor-qc-qs}
设 $S$ 为概形，$X$ 为 $S$ 上的拟分离代数空间。
设 $E$ 为 $D_\QCoh(\mathcal{O}_X)$ 中的对象。设 $a \leq b$。
以下条件等价：
\begin{enumerate}
\item $E$ 具有 $[a, b]$ 中的 Tor 振幅；
\item 对所有 $\mathcal{F}$，若其属于 $\QCoh(\mathcal{O}_X)$，则均有
$H^i(E \otimes_{\mathcal{O}_X}^\mathbf{L} \mathcal{F}) = 0$，
其中 $i \not \in [a, b]$。
\end{enumerate}
\end{lemma}

\begin{proof}
显然 (1) 蕴含 (2)。假设 (2)。设 $j : U \to X$ 为平展态射，
其中 $U$ 仿射。由于 $X$ 拟分离，$j$ 既拟紧又分离。
因此 $j_*$ 将拟凝聚模变为拟凝聚模（《空间的态射》引理
\ref{spaces-morphisms-lemma-pushforward}）。
任取拟凝聚模 $\mathcal{G}$，其定义在 $U$ 上。
由假设，
$H^i(E \otimes_{\mathcal{O}_X}^\mathbf{L} j_*\mathcal{G})$
在 $i \not \in [a, b]$ 时消失。经平坦态射 $j$ 拉回，得到
$H^i(E|_U \otimes_{\mathcal{O}_U}^\mathbf{L} j^*j_*\mathcal{G})$
在 $i \not \in [a, b]$ 时消失。由《空间的上同调》引理
\ref{spaces-cohomology-lemma-etale-pull-push-split}，
$\mathcal{G}$ 是 $j^*j_*\mathcal{G}$ 的直接因子。
因此条件 (2) 推出
$H^i(E|_U \otimes_{\mathcal{O}_U}^\mathbf{L} \mathcal{G})$
在 $i \not \in [a, b]$ 时消失，这对每个拟凝聚
$\mathcal{O}_U$-模 $\mathcal{G}$ 都成立。由于只需证明
$E|_U$ 具有 $[a, b]$ 中的 Tor 振幅，故归约到 $X$ 可表的情形。

\medskip\noindent
若 $X$ 可由概形 $X_0$ 表示，则（引理
\ref{spaces-perfect-lemma-derived-quasi-coherent-small-etale-site}）
可写 $E = \epsilon^*E_0$，其中 $E_0$ 是
$D_\QCoh(\mathcal{O}_{X_0})$ 中的对象，且
$\epsilon : X_\etale \to (X_0)_{Zar}$ 如
（\ref{spaces-perfect-equation-epsilon}）中所述。对于每个拟凝聚模
$\mathcal{F}_0$（定义在 $X_0$ 上），模 $\epsilon^*\mathcal{F}_0$ 在 $X$ 上拟凝聚，并且
$$
H^i(E \otimes_{\mathcal{O}_X}^\mathbf{L} \epsilon^*\mathcal{F}_0)
=
\epsilon^*H^i(E_0 \otimes_{\mathcal{O}_{X_0}}^\mathbf{L} \mathcal{F}_0)
$$
因为 $\epsilon$ 平坦（引理 \ref{spaces-perfect-lemma-epsilon-flat}）。
此外，这些层在 $i \not \in [a, b]$ 时消失，由同一引理可推出
$H^i(E_0 \otimes_{\mathcal{O}_{X_0}}^\mathbf{L} \mathcal{F}_0)$
也具有相同的消灭性。因此问题归约到概形情形，该情形在
《概形的导出范畴》引理 \ref{perfect-lemma-tor-qc-qs} 中处理。
\end{proof}

\begin{lemma}
\label{spaces-perfect-lemma-descend-RHom}
设 $X$ 为概形，$E, F$ 为 $D(\mathcal{O}_X)$ 中的对象。
假设下列条件之一成立：
\begin{enumerate}
\item $E$ 是拟凝聚的，且 $F$ 属于 $D^+(\mathcal{O}_X)$；
\item $E$ 是完美的，而 $F$ 任意。
\end{enumerate}
则存在典范同构
$$
\epsilon^*R\SheafHom(E, F) \longrightarrow R\SheafHom(\epsilon^*E, \epsilon^*F)
$$
这里 $\epsilon$ 如（\ref{spaces-perfect-equation-epsilon}）中所述。
\end{lemma}

\begin{proof}
回忆 $\epsilon$ 是平坦的（引理 \ref{spaces-perfect-lemma-epsilon-flat}），
因而 $\epsilon^* = L\epsilon^*$。由《位点上的上同调》评注
\ref{sites-cohomology-remark-prepare-fancy-base-change}，
存在从左端到右端的典范映射。
为了证明它是同构，可以局部工作，即不妨设 $X$ 为仿射概形。

\medskip\noindent
在情形 (1) 中，可以将 $E$ 表示为有上界复形
$\mathcal{E}^\bullet$，它由有限自由 $\mathcal{O}_X$-模构成；见《概形的导出范畴》引理
\ref{perfect-lemma-lift-pseudo-coherent}。
也可将 $F$ 表示为有下界复形 $\mathcal{F}^\bullet$，
它由 $\mathcal{O}_X$-模构成。应用《上同调》引理
\ref{cohomology-lemma-Rhom-complex-of-direct-summands-finite-free}，
可知 $R\SheafHom(E, F)$ 由各项为
$$
\bigoplus\nolimits_{n = - p + q}
\SheafHom_{\mathcal{O}_X}(\mathcal{E}^p, \mathcal{F}^q)
$$
的复形表示。应用《位点上的上同调》引理
\ref{sites-cohomology-lemma-Rhom-complex-of-direct-summands-finite-free}，
可知 $R\SheafHom(\epsilon^*E, \epsilon^*F)$ 由各项为
$$
\bigoplus\nolimits_{n = - p + q}
\SheafHom_{\mathcal{O}_\etale}
(\epsilon^*\mathcal{E}^p, \epsilon^*\mathcal{F}^q)
$$
的复形表示。因此引理的陈述归结为以下事实：典范映射
$$
\epsilon^*\SheafHom_{\mathcal{O}_X}(\mathcal{E}, \mathcal{F})
\longrightarrow
\SheafHom_{\mathcal{O}_\etale}
(\epsilon^*\mathcal{E}, \epsilon^*\mathcal{F})
$$
对于任意 $\mathcal{O}_X$-模 $\mathcal{F}$ 以及有限自由
$\mathcal{O}_X$-模 $\mathcal{E}$ 都是同构。

\medskip\noindent
在情形 (2) 中，可以将 $E$ 表示为严格完美复形
$\mathcal{E}^\bullet$，它由 $\mathcal{O}_X$-模构成；使用《概形的导出范畴》引理
\ref{perfect-lemma-affine-compare-bounded} 和
\ref{perfect-lemma-perfect-affine}，以及模的完美复形可由有限生成投射模构成的
有限复形表示这一事实。于是可以逐字使用上面的证明，只需将对《上同调》引理
\ref{cohomology-lemma-Rhom-complex-of-direct-summands-finite-free}
的引用替换为对《上同调》引理
\ref{cohomology-lemma-Rhom-strictly-perfect} 的引用。
\end{proof}

\begin{lemma}
\label{spaces-perfect-lemma-quasi-coherence-internal-hom}
设 $S$ 为概形，$X$ 为 $S$ 上的代数空间。
设 $L, K$ 为 $D(\mathcal{O}_X)$ 中的对象。
若下列条件之一成立：
\begin{enumerate}
\item $L$ 属于 $D^+_\QCoh(\mathcal{O}_X)$，且 $K$ 是拟凝聚的；
\item $L$ 属于 $D_\QCoh(\mathcal{O}_X)$，且 $K$ 是完美的，
\end{enumerate} % P09-ERR-0430
则 $R\SheafHom(K, L)$ 属于 $D_\QCoh(\mathcal{O}_X)$。
\end{lemma}

\begin{proof}
这由概形的相应结果（《概形的导出范畴》引理
\ref{perfect-lemma-quasi-coherence-internal-hom}）推出，
其中使用引理 \ref{spaces-perfect-lemma-check-quasi-coherence-on-covering} 的判别准则、
引理 \ref{spaces-perfect-lemma-descend-pseudo-coherent} 和
\ref{spaces-perfect-lemma-descend-perfect} 的判别准则，以及引理
\ref{spaces-perfect-lemma-descend-RHom} 的结果。
\end{proof}

\begin{lemma}
\label{spaces-perfect-lemma-internal-hom-evaluate-tensor-isomorphism}
设 $S$ 为概形。
设 $X$ 为 $S$ 上的代数空间。
设 $K, L, M$ 为 $D_\QCoh(\mathcal{O}_X)$ 中的对象。
《位点上的上同调》引理
\ref{sites-cohomology-lemma-internal-hom-diagonal-better} 中的映射
$$
K \otimes_{\mathcal{O}_X}^\mathbf{L} R\SheafHom(M, L)
\longrightarrow
R\SheafHom(M, K \otimes_{\mathcal{O}_X}^\mathbf{L} L)
$$
在下列情形中为同构：
\begin{enumerate}
\item $M$ 是完美的；或
\item $K$ 是完美的；或
\item $M$ 是拟凝聚的，$L \in D^+(\mathcal{O}_X)$，且 $K$ 具有有限
Tor 维数。
\end{enumerate}
\end{lemma}

\begin{proof}
映射是否为同构可在平展局部检验，因而不妨设 $X$ 为仿射概形。
由引理 \ref{spaces-perfect-lemma-derived-quasi-coherent-small-etale-site}，
可将 $K, L, M$ 写成 $\epsilon^*K_0, \epsilon^*L_0, \epsilon^*M_0$，
其中某些 $K_0, L_0, M_0$ 属于 $D_\QCoh(\mathcal{O}_X)$。
于是由引理 \ref{spaces-perfect-lemma-descend-RHom} 可见，情形 (1) 和 (3)
化归为概形的情形，即《概形的导出范畴》引理
\ref{perfect-lemma-internal-hom-evaluate-tensor-isomorphism}。
如果 $K$ 是完美的而不作其他假设，则我们不知道箭头的任一端
是否属于 $D_\QCoh(\mathcal{O}_X)$；但结论仍然成立，因为进一步局部化后，
$K$ 可由有限自由模构成的有限复形表示，此时结论显然。
\end{proof}









\section{用完美复形作逼近}
\label{spaces-perfect-section-approximation}

\noindent
本节继续《概形的导出范畴》第
\ref{perfect-section-approximation} 节中开始的讨论。

\begin{definition}
\label{spaces-perfect-definition-approximation-holds}
设 $S$ 为概形，$X$ 为 $S$ 上的代数空间。
考虑三元组 $(T, E, m)$，其中
\begin{enumerate}
\item $T \subset |X|$ 是闭子集；
\item $E$ 是 $D_\QCoh(\mathcal{O}_X)$ 中的对象；且
\item $m \in \mathbf{Z}$。
\end{enumerate}
所谓三元组 $(T, E, m)$ 的{\it 逼近成立}，是指存在完美对象 $P$，
它是 $D(\mathcal{O}_X)$ 中支撑在 $T$ 上的对象，并存在映射
$\alpha : P \to E$，它诱导同构 $H^i(P) \to H^i(E)$（对 $i > m$），
以及满射 $H^m(P) \to H^m(E)$。
\end{definition}

\noindent
逼近不可能对每个三元组都成立。其原因见《概形的导出范畴》定义
\ref{perfect-definition-approximation-holds} 后的评注。

\begin{definition}
\label{spaces-perfect-definition-approximation}
设 $S$ 为概形，$X$ 为 $S$ 上的代数空间。
所谓在 $X$ 上{\it 用完美复形作逼近成立}，是指对于任意闭子集
$T \subset |X|$，只要态射
$X \setminus T \to X$ 拟紧，就存在整数 $r$，使得对定义
\ref{spaces-perfect-definition-approximation-holds} 中任意满足下列条件的三元组
$(T, E, m)$：
\begin{enumerate}
\item $E$ 是 $(m - r)$-拟凝聚的；且
\item $H^i(E)$ 支撑在 $T$ 上（对 $i \geq m - r$），
\end{enumerate}
逼近均成立。
\end{definition}

\begin{lemma}
\label{spaces-perfect-lemma-pushforward-perfect}
设 $S$ 为概形，$(U \subset X, j : V \to X)$ 为 $S$ 上代数空间的
初等特异方形。设 $E$ 为 $D(\mathcal{O}_V)$ 中支撑在
$j^{-1}(T)$ 上的完美对象，其中 $T = |X| \setminus |U|$。
则 $Rj_*E$ 是 $D(\mathcal{O}_X)$ 中的完美对象。
\end{lemma}

\begin{proof}
完美性在 $X_\etale$ 上是局部的。因此，只需检验 $Rj_*E$
在 $U$ 和 $V$ 上的限制是完美的。由引理
\ref{spaces-perfect-lemma-pushforward-with-support-in-open}，有 $Rj_*E|_V = E$，
故其完美。又有 $Rj_*E|_U = 0$，因为
$E|_{V \setminus j^{-1}(T)} = 0$（使用引理
\ref{spaces-perfect-lemma-restrict-direct-image-open}）。
\end{proof}

\begin{lemma}
\label{spaces-perfect-lemma-open}
设 $S$ 为概形，$(U \subset X, j : V \to X)$ 为 $S$ 上代数空间的
初等特异方形。设 $T$ 为 $|X| \setminus |U|$ 的闭子集，并设
$(T, E, m)$ 为定义 \ref{spaces-perfect-definition-approximation-holds} 中的三元组。
若
\begin{enumerate}
\item 三元组 $(j^{-1}T, E|_V, m)$ 的逼近成立；且
\item 层 $H^i(E)$（对 $i \geq m$）支撑在 $T$ 上，
\end{enumerate}
则三元组 $(T, E, m)$ 的逼近成立。
\end{lemma}

\begin{proof}
设 $P \to E|_V$ 是三元组 $(j^{-1}T, E|_V, m)$ 在 $V$ 上的一个逼近。
由引理 \ref{spaces-perfect-lemma-pushforward-perfect}，$Rj_*P$ 是
$D(\mathcal{O}_X)$ 中的完美对象。另一方面，由引理
\ref{spaces-perfect-lemma-pushforward-with-support-in-open}，$Rj_*P = j_!P$。
例如由（\ref{spaces-perfect-equation-stalk-j-shriek}）可见，$j_!P$ 支撑在 $T$ 上。
因此得到逼近 $Rj_*P = j_!P \to j_!(E|_V) \to E$。
\end{proof}

\begin{lemma}
\label{spaces-perfect-lemma-approximation-affine}
设 $S$ 为概形，$X$ 为 $S$ 上可由仿射概形表示的代数空间。
若定义 \ref{spaces-perfect-definition-approximation-holds} 中的三元组
$(T, E, m)$ 满足：存在整数 $r \geq 0$ 使得
\begin{enumerate}
\item $E$ 是 $m$-拟凝聚的；
\item $H^i(E)$ 支撑在 $T$ 上（对 $i \geq m - r + 1$）；
\item $X \setminus T$ 是 $r$ 个仿射开集的并，
\end{enumerate}
则其逼近成立。特别地，仿射概形上的完美复形逼近成立。
\end{lemma}

\begin{proof}
设 $X_0$ 为表示 $X$ 的仿射概形。设 $T_0 \subset X_0$
为与 $T$ 对应的闭子集。% P09-ERR-0431
令 $\epsilon : X_\etale \to X_{0, Zar}$ 为（\ref{spaces-perfect-equation-epsilon}）
中的态射。由引理 \ref{spaces-perfect-lemma-derived-quasi-coherent-small-etale-site}，
可写成 $E = \epsilon^*E_0$，其中某个对象 $E_0$ 属于
$D_\QCoh(\mathcal{O}_{X_0})$。
由引理 \ref{spaces-perfect-lemma-descend-pseudo-coherent}，$E_0$ 是
$m$-拟凝聚的。比较上同调层的茎（见引理
\ref{spaces-perfect-lemma-epsilon-flat} 的证明），可见 $H^i(E_0)$ 支撑在 $T_0$ 上
（对 $i \geq m - r + 1$）。由《概形的导出范畴》引理
\ref{perfect-lemma-approximation-affine}，存在逼近 $P_0 \to E_0$，
它对应于三元组 $(T_0, E_0, m)$。
由引理 \ref{spaces-perfect-lemma-descend-perfect}，可见
$P = \epsilon^*P_0$ 是 $D(\mathcal{O}_X)$ 中的完美对象。
拉回后得到所需逼近 $P = \epsilon^*P_0 \to \epsilon^*E_0 = E$。
\end{proof}

\begin{lemma}
\label{spaces-perfect-lemma-induction-step}
设 $S$ 为概形，$(U \subset X, j : V \to X)$ 为 $S$ 上代数空间的
初等特异方形。假设 $U$ 拟紧、$V$ 仿射，且 $U \times_X V$ 拟紧。
若 $U$ 上用完美复形作逼近成立，则 $X$ 上用完美复形作逼近也成立。
\end{lemma}

\begin{proof}
设 $T \subset |X|$ 为闭子集，且 $X \setminus T \to X$ 拟紧。
设 $r_U$ 为定义 \ref{spaces-perfect-definition-approximation} 中适合于对子
$(U, T \cap |U|)$ 的整数。令 $T' = T \setminus |U|$，并赋予
$T'$ 诱导的既约子空间结构。由于 $|T'|$ 包含于
$|X| \setminus |U|$，可见 $j^{-1}(T') \to T'$ 是同构。此外，
$V \setminus j^{-1}(T')$ 拟紧，因为它是 $U \times_X V$ 与
$X \setminus T$ 在 $X$ 上的纤维积，并且已假设
$U \times_X V$ 拟紧且 $X \setminus T \to X$ 拟紧。
设 $r'$ 为覆盖 $V \setminus j^{-1}(T')$ 所需的仿射开集数。
我们断言 $r = \max(r_U, r')$ 对于对子 $(X, T)$ 可用。

\medskip\noindent
为证明此点，选取三元组 $(T, E, m)$，使得 $E$ 是
$(m - r)$-拟凝聚的，且 $H^i(E)$ 支撑在 $T$ 上（对
$i \geq m - r$）。设 $t$ 为使 $H^t(E)|_U$ 非零的最大整数。
（这样的整数存在，因为 $U$ 拟紧，且 $E|_U$ 是
$(m - r)$-拟凝聚的。）我们将证明 $E$ 可被逼近，并对 $t$ 作归纳。

\medskip\noindent
基本情形：$t \leq m - r'$。这意味着 $H^i(E)$ 支撑在
$T'$ 上（对 $i \geq m - r'$）。因此引理
\ref{spaces-perfect-lemma-approximation-affine} 保证存在逼近
$P \to E|_V$，它是三元组 $(T', E|_V, m)$ 在 $V$ 上的逼近。
应用引理 \ref{spaces-perfect-lemma-open} 可见 $(T', E, m)$ 可被逼近。
这样的逼近也是 $(T, E, m)$ 的逼近。

\medskip\noindent
归纳步骤。选取逼近 $P \to E|_U$，它对应于
$(T \cap |U|, E|_U, m)$。这尤其给出满射
$H^t(P) \to H^t(E|_U)$。
在证明的余下部分，我们将对可表的代数空间 $V$ 和
$U \times_X V$ 使用引理
\ref{spaces-perfect-lemma-derived-quasi-coherent-small-etale-site} 的等价
（以及评注 \ref{spaces-perfect-remark-match-total-direct-images} 的相容性）。
我们还将使用如下事实：$(m - r)$-拟凝聚性以及完美性在
Zariski 位点和平展位点上分别相同；见引理
\ref{spaces-perfect-lemma-descend-pseudo-coherent} 和
\ref{spaces-perfect-lemma-descend-perfect}。
因此，可对开浸入 $U \times_X V \subset V$ 使用
《概形的导出范畴》第 \ref{perfect-section-lift} 节的结果。
由此，《概形的导出范畴》引理
\ref{perfect-lemma-lift-perfect-complex-plus-shift-support} 表明，
存在完美对象 $Q$，它属于 $D(\mathcal{O}_V)$ 且支撑在 $j^{-1}(T)$ 上，
以及同构
$Q|_{U \times_X V} \to (P \oplus P[1])|_{U \times_X V}$。
由《概形的导出范畴》引理 \ref{perfect-lemma-lift-map}，
可将 $Q$ 替换为 $Q \otimes^\mathbf{L} I$，并且可以假设映射
$$
Q|_{U \times_X V} \longrightarrow
(P \oplus P[1])|_{U \times_X V} \longrightarrow
P|_{U \times_X V} \longrightarrow E|_{U \times_X V}
$$
可提升为 $Q \to E|_V$。由引理 \ref{spaces-perfect-lemma-glue}，得到
态射 $a : R \to E$，它是 $D(\mathcal{O}_X)$ 中的态射，并使得 $a|_U$
同构于 $P \oplus P[1] \to E|_U$，而 $a|_V$ 同构于
$Q \to E|_V$。% P09-ERR-0432
于是 $R$ 是支撑在 $T$ 上的完美对象，并且映射
$H^t(R) \to H^t(E)$ 限制到 $U$ 上是满射。选取特异三角
$$
R \to E \to E' \to R[1]
$$
则 $E'$ 是 $(m - r)$-拟凝聚的
（《位点上的上同调》引理
\ref{sites-cohomology-lemma-cone-pseudo-coherent}），
$H^i(E')|_U = 0$（对 $i \geq t$），且 $H^i(E')$ 支撑在
$T$ 上（对 $i \geq m - r$）。归纳地，得到逼近 $R' \to E'$，
它对应于三元组 $(T, E', m)$。将复合
$R' \to E' \to R[1]$ 嵌入特异三角
$R \to R'' \to R' \to R[1]$，并将态射 $R' \to E'$ 和
$R[1] \to R[1]$ 扩充成特异三角的态射
$$
\xymatrix{
R \ar[r] \ar[d] & R'' \ar[d] \ar[r] & R' \ar[d] \ar[r] & R[1] \ar[d] \\
R \ar[r] &  E \ar[r] & E' \ar[r] & R[1]
}
$$
这里使用 TR3。于是 $R''$ 是完美复形
（《位点上的上同调》引理
\ref{sites-cohomology-lemma-two-out-of-three-perfect}），且支撑在 $T$ 上。
简单的图追踪表明，$R'' \to E$ 就是所需的逼近。
\end{proof}

\begin{theorem}
\label{spaces-perfect-theorem-approximation}
设 $S$ 为概形。
设 $X$ 为 $S$ 上拟紧且拟分离的代数空间。
则 $X$ 上用完美复形作逼近成立。
\end{theorem}

\begin{proof}
这由引理 \ref{spaces-perfect-lemma-induction-principle} 的归纳原理以及引理
\ref{spaces-perfect-lemma-induction-step} 和 \ref{spaces-perfect-lemma-approximation-affine} 推出。
\end{proof}






\section{生成导出范畴}
\label{spaces-perfect-section-generating}

\noindent
本节是《概形的导出范畴》第 \ref{perfect-section-generating} 节的类似物。
不过，我们先证明下述引理；它是《概形的导出范畴》引理
\ref{perfect-lemma-lift-map-from-perfect-complex-with-support} 的类似物。

\begin{lemma}
\label{spaces-perfect-lemma-lift-map-from-perfect-complex-with-support}
设 $S$ 为概形，$X$ 为 $S$ 上拟紧且拟分离的代数空间。
设 $W \subset X$ 为拟紧开集。设 $T \subset |X|$ 为闭子集，使得
$X \setminus T \to X$ 是拟紧态射。设 $E$ 为
$D_\QCoh(\mathcal{O}_X)$ 中的对象。设 $\alpha : P \to E|_W$ 为映射，
其中 $P$ 是 $D(\mathcal{O}_W)$ 中支撑在 $T \cap W$ 上的完美对象。
则存在映射 $\beta : R \to E$，其中 $R$ 是
$D(\mathcal{O}_X)$ 中支撑在 $T$ 上的完美对象，使得 $P$ 是
$R|_W$ 的直和项（在 $D(\mathcal{O}_W)$ 中），并且与 $\alpha$ 和
$\beta|_W$ 相容。% P09-ERR-0433
\end{lemma}

\begin{proof}
我们将用引理 \ref{spaces-perfect-lemma-induction-principle-enlarge} 的归纳原理证明此事。
因此立即化归到如下情形：有初等特异方形
$(W \subset X, f : V \to X)$，其中 $V$ 仿射，且
$P \to E|_W$ 如引理陈述中所述。在证明的余下部分，我们将对可表代数空间
$V$ 和 $W \times_X V$ 使用引理
\ref{spaces-perfect-lemma-derived-quasi-coherent-small-etale-site}
（以及评注 \ref{spaces-perfect-remark-match-total-direct-images} 的相容性）。
我们还将使用如下事实：完美性在 Zariski 位点和平展位点上相同；
见引理 \ref{spaces-perfect-lemma-descend-perfect}。

\medskip\noindent
由《概形的导出范畴》引理
\ref{perfect-lemma-lift-perfect-complex-plus-shift-support}，
可选取完美对象 $Q$，它属于 $D(\mathcal{O}_V)$，支撑在 $f^{-1}T$ 上，
并且有同构
$Q|_{W \times_X V} \to (P \oplus P[1])|_{W \times_X V}$。
由《概形的导出范畴》引理 \ref{perfect-lemma-lift-map}，
可将 $Q$ 替换为 $Q \otimes^\mathbf{L} I$（仍支撑在 $f^{-1}T$ 上），
并且可以假设映射
$$
Q|_{W \times_X V} \to (P \oplus P[1])|_{W \times V} % P09-ERR-0434
\longrightarrow P|_{W \times_X V}
\longrightarrow
E|_{W \times_X V}
$$
可提升为 $Q \to E|_V$。由引理 \ref{spaces-perfect-lemma-glue}，得到态射
$a : R \to E$，它是 $D(\mathcal{O}_X)$ 中的态射，使得 $a|_W$
同构于 $P \oplus P[1] \to E|_W$，而 $a|_V$ 同构于
$Q \to E|_V$。% P09-ERR-0435
于是 $R$ 是支撑在 $T$ 上的完美对象，如所要求。
\end{proof}

\begin{remark}
\label{spaces-perfect-remark-addendum}
引理 \ref{spaces-perfect-lemma-lift-map-from-perfect-complex-with-support} 的证明表明
$$
R|_W = P \oplus P^{\oplus n_1}[1] \oplus \ldots \oplus P^{\oplus n_m}[m]
$$
其中某些 $m \geq 0$ 且 $n_j \geq 0$。因此，$R|_W$ 的最高次非零
上同调层等于 $P$ 的最高次非零上同调层。对映射
$P^{\oplus n_1}[1] \oplus \ldots \oplus P^{\oplus n_m}[m] \to R|_W$
重复该构造、取锥并使用归纳，即可使 $R|_W$ 与 $P$ 在任意给定次数
以上的上同调层相等。
\end{remark}

\begin{lemma}
\label{spaces-perfect-lemma-direct-summand-of-a-restriction}
设 $S$ 为概形。
设 $X$ 为 $S$ 上拟紧且拟分离的代数空间。
设 $W$ 为 $X$ 的拟紧开子空间。
设 $P$ 为 $D(\mathcal{O}_W)$ 中的完美对象。
则 $P$ 是 $D(\mathcal{O}_X)$ 中某个完美对象之限制的直和项。
\end{lemma}

\begin{proof}
这是引理 \ref{spaces-perfect-lemma-lift-map-from-perfect-complex-with-support} 的特殊情形。
\end{proof}

\begin{theorem}
\label{spaces-perfect-theorem-bondal-van-den-Bergh}
设 $S$ 为概形，$X$ 为 $S$ 上拟紧且拟分离的代数空间。
范畴 $D_\QCoh(\mathcal{O}_X)$ 可由单个完美对象生成。
更精确地，存在完美对象 $P$，它属于 $D(\mathcal{O}_X)$，使得对于
$E \in D_\QCoh(\mathcal{O}_X)$，下列条件等价：
\begin{enumerate}
\item $E = 0$；且
\item $\Hom_{D(\mathcal{O}_X)}(P[n], E) = 0$（对所有
$n \in \mathbf{Z}$）。
\end{enumerate}
\end{theorem}

\begin{proof}
我们将使用引理 \ref{spaces-perfect-lemma-induction-principle} 的归纳原理证明此事。

\medskip\noindent
若 $X$ 仿射，则 $\mathcal{O}_X$ 是完美生成元。这由引理
\ref{spaces-perfect-lemma-derived-quasi-coherent-small-etale-site} 和
《概形的导出范畴》引理
\ref{perfect-lemma-affine-compare-bounded} 推出。

\medskip\noindent
假设 $(U \subset X, f : V \to X)$ 是初等特异方形，其中 $U$ 拟紧，
定理对 $U$ 成立，且 $V$ 是仿射概形。
设 $P$ 为 $D(\mathcal{O}_U)$ 中的完美对象，它是
$D_\QCoh(\mathcal{O}_U)$ 的生成元。使用引理
\ref{spaces-perfect-lemma-direct-summand-of-a-restriction}，可选取完美对象
$Q$，它属于 $D(\mathcal{O}_X)$，其在 $U$ 上的限制是某个直和，
且 $P$ 是该直和的一个直和项。记 $V = \Spec(A)$。设 $Z \subset V$
为 $X \setminus U$ 的逆像所给的既约闭子概形，并同构地映到该逆像
（见定义 \ref{spaces-perfect-definition-elementary-distinguished-square}）。
这是 $V$ 的逆紧闭子集。选取 $f_1, \ldots, f_r \in A$，使得
$Z = V(f_1, \ldots, f_r)$。设 $K \in D(\mathcal{O}_V)$ 为与
$f_1, \ldots, f_r$ 在 $A$ 上的 Koszul 复形相对应的完美对象。
注意，因为 $K$ 支撑在 $Z$ 上，推前
$K' = Rf_*K$ 是 $D(\mathcal{O}_X)$ 中的完美对象，其在 $V$ 上的限制
是 $K$（见引理 \ref{spaces-perfect-lemma-pushforward-perfect} 和
\ref{spaces-perfect-lemma-pushforward-with-support-in-open}）。
我们断言 $Q \oplus K'$ 是 $D_\QCoh(\mathcal{O}_X)$ 的生成元。

\medskip\noindent
设 $E$ 为 $D_\QCoh(\mathcal{O}_X)$ 中的对象，使得从
$Q \oplus K'$ 的任意平移到 $E$ 都没有非零映射。
由引理 \ref{spaces-perfect-lemma-pushforward-with-support-in-open}，有
$K' =  f_! K$，因而
$$
\Hom_{D(\mathcal{O}_X)}(K'[n], E) = \Hom_{D(\mathcal{O}_V)}(K[n], E|_V)
$$
于是由《概形的导出范畴》引理
\ref{perfect-lemma-orthogonal-koszul-complex}
（还使用引理 \ref{spaces-perfect-lemma-derived-quasi-coherent-small-etale-site}），
这些群的消失蕴含 $E|_V$ 同构于
$R(U \times_X V \to V)_*E|_{U \times_X V}$。这又蕴含
$E = R(U \to X)_*E|_U$（略去一个小细节）。在此情形下
$$
\Hom_{D(\mathcal{O}_X)}(Q[n], E) = \Hom_{D(\mathcal{O}_U)}(Q|_U[n], E|_U)
$$
其中含有 $\Hom_{D(\mathcal{O}_U)}(P[n], E|_U)$ 作为直和项。
因此，根据对 $P$ 的选取，这些群的消失蕴含 $E|_U$ 为零，
从而 $E$ 为零。
\end{proof}












\begin{remark}
\label{spaces-perfect-remark-pullback-generator}
设 $S$ 为概形。
设 $f : X \to Y$ 为 $S$ 上拟紧且拟分离的代数空间之间的态射。
设 $E \in D_\QCoh(\mathcal{O}_Y)$ 为生成元
（见定理 \ref{spaces-perfect-theorem-bondal-van-den-Bergh}）。
则下列条件等价：
\begin{enumerate}
\item 对 $K \in D_\QCoh(\mathcal{O}_X)$，有
$Rf_*K = 0$ 当且仅当 $K = 0$；
\item $Rf_* : D_\QCoh(\mathcal{O}_X) \to D_\QCoh(\mathcal{O}_Y)$
反映同构；且
\item $Lf^*E$ 是 $D_\QCoh(\mathcal{O}_X)$ 的生成元。
\end{enumerate}
(1) 与 (2) 的等价性是如下事实的形式推论：
$Rf_* : D_\QCoh(\mathcal{O}_X) \to D_\QCoh(\mathcal{O}_Y)$ 是三角范畴的
正合函子。类似地，(1) 与 (3) 的等价性由 $Lf^*$ 是 $Rf_*$ 的左伴随
这一事实形式地推出。若 $f$ 仿射（引理 \ref{spaces-perfect-lemma-affine-morphism}），
或 $f$ 是开浸入，或 $f$ 是这类态射的复合，则这些条件成立。
\end{remark}

\noindent
下述结果是定理 \ref{spaces-perfect-theorem-bondal-van-den-Bergh} 的一个加强，
并且用完全相同的方法证明。% P09-ERR-0436
设 $T \subset |X|$ 为闭子集，其中 $X$ 是代数空间。
记 $D_T(\mathcal{O}_X)$ 为由上同调层支撑在 $T$ 上的复形所构成的
严格全、饱和三角子范畴。

\begin{lemma}
\label{spaces-perfect-lemma-generator-with-support}
设 $S$ 为概形，$X$ 为 $S$ 上拟紧且拟分离的代数空间。
设 $T \subset |X|$ 为闭子集，且 $|X| \setminus T$ 拟紧。
沿用上述记号，范畴 $D_{\QCoh, T}(\mathcal{O}_X)$ 可由单个完美对象生成。
\end{lemma}

\begin{proof}
我们将使用引理 \ref{spaces-perfect-lemma-induction-principle} 的归纳原理证明此事。
对于位点 $X_{spaces, \etale}$ 中可表、拟紧且拟分离的对象，该性质由
《概形的导出范畴》引理
\ref{perfect-lemma-generator-with-support} 成立。

\medskip\noindent
假设 $(U \subset X, f : V \to X)$ 是初等特异方形，且引理对 $U$
成立，而 $V$ 仿射。为完成证明，需说明结论对 $X$ 成立。
设 $P$ 为 $D(\mathcal{O}_U)$ 中支撑在 $T \cap U$ 上的完美对象，
它是 $D_{\QCoh, T \cap U}(\mathcal{O}_U)$ 的生成元。使用引理
\ref{spaces-perfect-lemma-lift-map-from-perfect-complex-with-support}，可选取
完美对象 $Q$，它属于 $D(\mathcal{O}_X)$ 且支撑在 $T$ 上，其在 $U$ 上的限制
是某个直和，且 $P$ 是该直和的一个直和项。
写 $V = \Spec(B)$。设 $Z = X \setminus U$。则 $f^{-1}Z$ 是 $V$ 的闭子集，
且 $V \setminus f^{-1}Z$ 拟紧。由于 $X$ 拟分离，可得
$f^{-1}Z \cap f^{-1}T = f^{-1}(Z \cap T)$ 是 $V$ 的闭子集，并且
$W = V \setminus f^{-1}(Z \cap T)$ 拟紧。因此可以选取
$g_1, \ldots, g_s \in B$，使得
$f^{-1}(Z \cap T) = V(g_1, \ldots, g_r)$。% P09-ERR-0437
设 $K \in D(\mathcal{O}_V)$ 为与 $g_1, \ldots, g_s$ 在 $B$ 上的
Koszul 复形相对应的完美对象。注意，因为 $K$ 支撑在闭子集
$f^{-1}(Z \cap T) \subset V$ 上，% P09-ERR-0438
推前
$K' = R(V \to X)_*K$ 是 $D(\mathcal{O}_X)$ 中的完美对象，
其在 $V$ 上的限制是 $K$（见引理 \ref{spaces-perfect-lemma-pushforward-perfect} 和
\ref{spaces-perfect-lemma-pushforward-with-support-in-open}）。
我们断言 $Q \oplus K'$ 是 $D_{\QCoh, T}(\mathcal{O}_X)$ 的生成元。

\medskip\noindent
设 $E$ 为 $D_{\QCoh, T}(\mathcal{O}_X)$ 中的对象，使得从
$Q \oplus K'$ 的任意平移到 $E$ 都没有非零映射。
由引理 \ref{spaces-perfect-lemma-pushforward-with-support-in-open}，有
$K' =  R(V \to X)_! K$，因而
$$
\Hom_{D(\mathcal{O}_X)}(K'[n], E) = \Hom_{D(\mathcal{O}_V)}(K[n], E|_V)
$$
于是由《概形的导出范畴》引理
\ref{perfect-lemma-orthogonal-koszul-complex}，有
$E|_V = Rj_*E|_W$，其中 $j : W \to V$ 是包含。图示为
$$
\xymatrix{
W \ar[r]_j & V & Z \cap T \ar[l] \ar[d] \\
V \setminus f^{-1}Z \ar[u]^{j'} \ar[ru]_{j''} & & Z \ar[lu]
}
$$
由于 $E$ 支撑在 $T$ 上，可见 $E|_W$ 支撑在
$f^{-1}T \cap W = f^{-1}T \cap (V \setminus f^{-1}Z)$ 上，
这是 $W$ 中的闭子集。由此得到
$$
E|_V = Rj_*(E|_W) = Rj_*(Rj'_*(E|_{U \cap V})) = Rj''_*(E|_{U \cap V})
$$
这里第二个等号是《上同调》引理
\ref{cohomology-lemma-pushforward-restriction} 的 (1)。该引理可以应用，
因为 $V$ 是概形且 $E$ 具有拟凝聚上同调层，所以沿拟紧开浸入 $j'$
的推前与底层概形上的推前一致；见评注
\ref{spaces-perfect-remark-match-total-direct-images}。
这蕴含 $E = R(U \to X)_*E|_U$（略去一个小细节）。在此情形下
$$
\Hom_{D(\mathcal{O}_X)}(Q[n], E) = \Hom_{D(\mathcal{O}_U)}(Q|_U[n], E|_U)
$$
其中含有 $\Hom_{D(\mathcal{O}_U)}(P[n], E|_U)$ 作为直和项。
因此，根据对 $P$ 的选取，这些群的消失蕴含 $E|_U$ 为零，
从而 $E$ 为零。
\end{proof}

\section{紧对象与完美对象}
\label{spaces-perfect-section-compact}

\noindent
本节是《概形的导出范畴》第 \ref{perfect-section-compact} 节的类似物。

\begin{proposition}
\label{spaces-perfect-proposition-compact-is-perfect}
设 $S$ 为概形。
设 $X$ 为 $S$ 上拟紧且拟分离的代数空间。
$D_\QCoh(\mathcal{O}_X)$ 中的对象是紧的，当且仅当它是完美的。
\end{proposition}

\begin{proof}
若 $K$ 是 $D(\mathcal{O}_X)$ 中的完美对象，其对偶为 $K^\vee$
（《位点上的上同调》引理
\ref{sites-cohomology-lemma-dual-perfect-complex}），则有
$$
\Hom_{D(\mathcal{O}_X)}(K, M) =
H^0(X, K^\vee \otimes_{\mathcal{O}_X}^\mathbf{L} M)
$$
且关于 $M$ 函子性成立。由于 $K^\vee \otimes_{\mathcal{O}_X}^\mathbf{L} -$
与直和可交换，并且由引理
\ref{spaces-perfect-lemma-quasi-coherence-pushforward-direct-sums}，$H^0(X, -)$ 在
$D_\QCoh(\mathcal{O}_X)$ 上与直和可交换，故 $K$ 在
$D_\QCoh(\mathcal{O}_X)$ 中是紧的。

\medskip\noindent
反之，设 $K$ 为 $D_\QCoh(\mathcal{O}_X)$ 中的紧对象。
为证明 $K$ 是完美的，只需证明 $K|_U$ 是完美的，其中 $U$ 是任意
仿射概形并且平展于 $X$；见《位点上的上同调》引理
\ref{sites-cohomology-lemma-perfect-independent-representative}。
注意 $j : U \to X$ 是拟紧且分离的态射。因此
$Rj_* : D_\QCoh(\mathcal{O}_U) \to D_\QCoh(\mathcal{O}_X)$
与直和可交换；见引理
\ref{spaces-perfect-lemma-quasi-coherence-pushforward-direct-sums}。
于是，限制到 $U$ 与 $Rj_*$ 的伴随性蕴含 $K|_U$ 是
$D_\QCoh(\mathcal{O}_U)$ 中的完美对象。因此化归到 $X$ 仿射的情形，
特别地，它是拟紧且拟分离的概形。通过引理
\ref{spaces-perfect-lemma-derived-quasi-coherent-small-etale-site} 和
\ref{spaces-perfect-lemma-descend-perfect}，化归到概形的情形，即《概形的导出范畴》命题
\ref{perfect-proposition-compact-is-perfect}。
\end{proof}

\begin{remark}
\label{spaces-perfect-remark-classical-generator}
设 $S$ 为概形，$X$ 为 $S$ 上拟紧且拟分离的代数空间。
设 $G$ 为 $D(\mathcal{O}_X)$ 中的完美对象，并且是
$D_\QCoh(\mathcal{O}_X)$ 的生成元。由定理
\ref{spaces-perfect-theorem-bondal-van-den-Bergh}，至少存在一个这样的对象。
结合引理 \ref{spaces-perfect-lemma-quasi-coherence-direct-sums}、命题
\ref{spaces-perfect-proposition-compact-is-perfect} 以及《导出范畴》命题
\ref{derived-proposition-generator-versus-classical-generator}，
可见 $G$ 是 $D_{perf}(\mathcal{O}_X)$ 的经典生成元。
\end{remark}

\noindent
下述结果是命题 \ref{spaces-perfect-proposition-compact-is-perfect} 的加强。
设 $T \subset |X|$ 为闭子集，其中 $X$ 是代数空间。
如前，$D_T(\mathcal{O}_X)$ 表示由上同调层支撑在 $T$ 上的复形所构成的
严格全、饱和三角子范畴。由于取直和与取上同调层可交换，
可得 $D_T(\mathcal{O}_X)$ 有直和，并且这些直和等于
$D(\mathcal{O}_X)$ 中的直和。

\begin{lemma}
\label{spaces-perfect-lemma-compact-is-perfect-with-support}
设 $S$ 为概形，$X$ 为 $S$ 上拟紧且拟分离的代数空间。
设 $T \subset |X|$ 为闭子集，且 $|X| \setminus T$ 拟紧。
$D_{\QCoh, T}(\mathcal{O}_X)$ 中的对象是紧的，当且仅当它作为
$D(\mathcal{O}_X)$ 中的对象是完美的。
\end{lemma}

\begin{proof}
由引理前的评注，$D_{\QCoh, T}(\mathcal{O}_X)$ 是有直和的三角范畴。
由命题 \ref{spaces-perfect-proposition-compact-is-perfect}，完美对象给出
$D(\mathcal{O}_X)$ 中的紧对象，% P09-ERR-0440
因而更是任意在取直和下封闭的子范畴中的紧对象。反之，我们将使用如下事实：
存在生成元 $E \in D_{\QCoh, T}(\mathcal{O}_X)$，它是
$\mathcal{O}_X$-模的完美复形；见引理
\ref{spaces-perfect-lemma-generator-with-support}。% P09-ERR-0439
于是由上文，$E$ 是紧的。因此由《导出范畴》命题
\ref{derived-proposition-generator-versus-classical-generator}，
$E$ 是 $D_{\QCoh, T}(\mathcal{O}_X)$ 的紧对象所成全子范畴的经典生成元。
因此，任意紧对象均可由 $E$ 通过有限次下列运算构造出来：
(a) 取平移，(b) 取有限直和，(c) 取锥，(d) 取直和项。
这些运算中的每一种都保持完美性，结论随之成立。
\end{proof}

\begin{remark}
\label{spaces-perfect-remark-classical-generator-with-support}
设 $S$ 为概形，$X$ 为 $S$ 上拟紧且拟分离的代数空间。
设 $T \subset |X|$ 为闭子集，且 $|X| \setminus T$ 拟紧。
设 $G$ 为 $D_{\QCoh, T}(\mathcal{O}_X)$ 中的完美对象，并且是
$D_{\QCoh, T}(\mathcal{O}_X)$ 的生成元。由引理
\ref{spaces-perfect-lemma-generator-with-support}，至少存在一个这样的对象。
结合 $D_{\QCoh, T}(\mathcal{O}_X)$ 有直和这一事实、引理
\ref{spaces-perfect-lemma-compact-is-perfect-with-support} 以及《导出范畴》命题
\ref{derived-proposition-generator-versus-classical-generator}，
可见 $G$ 是 $D_{perf, T}(\mathcal{O}_X)$ 的经典生成元。
\end{remark}

\noindent
下述引理应用了前一引理证明中所用的思想。

\begin{lemma}
\label{spaces-perfect-lemma-map-from-pseudo-coherent-to-complex-with-support}
设 $S$ 为概形，$X$ 为 $S$ 上拟紧且拟分离的代数空间。
设 $T \subset |X|$ 为闭子集，使得补集 $U \subset X$ 拟紧。
设 $\alpha : P \to E$ 为 $D_\QCoh(\mathcal{O}_X)$ 中的态射，且满足
下列条件之一：
\begin{enumerate}
\item $P$ 是完美的，且 $E$ 支撑在 $T$ 上；或
\item $P$ 是拟凝聚的，$E$ 支撑在 $T$ 上，且 $E$ 有下界。
\end{enumerate}
则存在 $\mathcal{O}_X$-模的完美复形 $I$ 以及映射
$I \to \mathcal{O}_X[0]$，使得 $I \otimes^\mathbf{L} P \to E$ 为零，
并且 $I|_U \to \mathcal{O}_U[0]$ 是同构。
\end{lemma}

\begin{proof}
置 $\mathcal{D} = D_{\QCoh, T}(\mathcal{O}_X)$。在两种情形中，复形
$K = R\SheafHom(P, E)$ 都是 $\mathcal{D}$ 中的对象。拟凝聚性见引理
\ref{spaces-perfect-lemma-quasi-coherence-internal-hom}。显然，$K$ 支撑在 $T$ 上，
因为 $R\SheafHom$ 的构造与限制到开集可交换。映射 $\alpha$ 定义了
$H^0(K) = \Hom_{D(\mathcal{O}_X)}(\mathcal{O}_X[0], K)$ 中的元素。
于是，只需对映射 $\alpha : \mathcal{O}_X[0] \to K$ 证明结论。

\medskip\noindent
设 $E \in \mathcal{D}$ 为完美生成元；见引理
\ref{spaces-perfect-lemma-generator-with-support}。如《导出范畴》引理
\ref{derived-lemma-write-as-colimit} 中那样，写成
$$
K = \text{hocolim} K_n
$$
这里使用生成元 $E$。
由于函子 $\mathcal{D} \to D(\mathcal{O}_X)$ 与直和可交换，
可见 $K = \text{hocolim} K_n$ 在 $D(\mathcal{O}_X)$ 中成立。
由于 $\mathcal{O}_X$ 是 $D(\mathcal{O}_X)$ 中的紧对象，得到某个 $n$
以及态射 $\alpha_n : \mathcal{O}_X \to K_n$，它给出 $\alpha$；
见《导出范畴》引理
\ref{derived-lemma-commutes-with-countable-sums}。
对态射 $\mathcal{O}_X[0] \to K_n$（在环境范畴
$D(\mathcal{O}_X)$ 中）应用《导出范畴》引理
\ref{derived-lemma-factor-through}，可见 $\alpha_n$ 分解为
$\mathcal{O}_X[0] \to Q \to K_n$，其中 $Q$ 是
$\langle E \rangle$ 中的对象。由此断定 $Q$ 是支撑在 $T$ 上的完美复形。

\medskip\noindent
选取特异三角
$$
I \to \mathcal{O}_X[0] \to Q \to I[1]
$$
由构造，$I$ 是完美的，映射 $I \to \mathcal{O}_X[0]$ 在 $U$ 上限制为
同构，并且复合 $I \to K$ 为零，因为 $\alpha$ 经 $Q$ 分解。
这证明了引理。
\end{proof}








\section{作为模范畴的导出范畴}
\label{spaces-perfect-section-derived-is-dga}

\noindent
本节是《概形的导出范畴》第
\ref{perfect-section-derived-is-dga} 节的类似结果。

\begin{lemma}
\label{spaces-perfect-lemma-tensor-with-QCoh-complex}
设 $S$ 为概形，$X$ 为 $S$ 上的代数空间。设 $K^\bullet$
为 $\mathcal{O}_X$-模复形，其上同调层拟凝聚。置
$(E, d) = \Hom_{\text{Comp}^{dg}(\mathcal{O}_X)}(K^\bullet, K^\bullet)$
为自同态微分分次代数。则《微分分次代数》引理
\ref{dga-lemma-tensor-with-complex-derived} 中的函子
$$
- \otimes_E^\mathbf{L} K^\bullet :
D(E, \text{d}) \longrightarrow D(\mathcal{O}_X)
$$
的像包含在 $D_\QCoh(\mathcal{O}_X)$ 中。
\end{lemma}

\begin{proof}
设 $P$ 为微分分次 $E$-模，且具有性质 $P$。% P09-ERR-0445
设 $F_\bullet$ 为《微分分次代数》第
\ref{dga-section-P-resolutions} 节中给出的 $P$ 上的滤过。
于是
$$
P \otimes_E K^\bullet = \text{hocolim}\ F_iP \otimes_E K^\bullet
$$
每个 $F_iP$ 都有有限滤过，其分次片是 $E[k]$ 的直和。
结论立即得出。
\end{proof}

\noindent
下述引理可以加强（对于所有仅在某个固定范围内具有非零上同调层的
$L$，其消失性具有一致性）。

\begin{lemma}
\label{spaces-perfect-lemma-ext-from-perfect-into-bounded-QCoh}
设 $S$ 为概形。
设 $X$ 为 $S$ 上拟紧且拟分离的代数空间。
设 $K$ 为 $D(\mathcal{O}_X)$ 中的完美对象。则
\begin{enumerate}
\item 存在整数 $a \leq b$，使得
$\Hom_{D(\mathcal{O}_X)}(K, L) = 0$，只要
$L \in D_\QCoh(\mathcal{O}_X)$ 且有 $H^i(L) = 0$，其中
$i \in [a, b]$；并且
\item 如果 $L$ 有界，则 $\Ext^n_{D(\mathcal{O}_X)}(K, L)$
除有限多个 $n$ 外均为零。
\end{enumerate}
\end{lemma}

\begin{proof}
由于 $\Ext^n_{D(\mathcal{O}_X)}(K, L) =
\Hom_{D(\mathcal{O}_X)}(K, L[n])$，第 (2) 项由第 (1) 项推出。
我们证明第 (1) 项。由于 $K$ 完美，有
$$
\Ext^i_{D(\mathcal{O}_X)}(K, L) =
H^i(X, K^\vee \otimes_{\mathcal{O}_X}^\mathbf{L} L)
$$
其中 $K^\vee$ 是 $K$ 的“对偶”完美复形；见《位点上的上同调》引理
\ref{sites-cohomology-lemma-dual-perfect-complex}。
注意，由引理 \ref{spaces-perfect-lemma-quasi-coherence-tensor-product} 和
\ref{spaces-perfect-lemma-pseudo-coherent}（后者用于说明完美复形具有拟凝聚上同调层），
$P = K^\vee \otimes_{\mathcal{O}_X}^\mathbf{L} L$
属于 $D_\QCoh(X)$。设 $K^\vee$ 的 Tor 振幅属于 $[a, b]$。
于是谱序列
$$
E_1^{p, q} = H^p(K^\vee \otimes_{\mathcal{O}_X}^\mathbf{L} H^q(L))
\Rightarrow
H^{p + q}(K^\vee \otimes_{\mathcal{O}_X}^\mathbf{L} L)
$$
表明 $H^j(K^\vee \otimes_{\mathcal{O}_X}^\mathbf{L} L)$ 为零，
只要 $H^q(L) = 0$ 对于 $q \in [j - b, j - a]$ 成立。
设 $N$ 为《空间的上同调》引理
\ref{spaces-cohomology-lemma-vanishing-quasi-separated}
中的整数 $\max(d_p + p)$。那么，
$H^0(X, K^\vee \otimes_{\mathcal{O}_X}^\mathbf{L} L)$ 为零，
只要上同调层
$$
H^{-N}(K^\vee \otimes_{\mathcal{O}_X}^\mathbf{L} L),
\ H^{-N + 1}(K^\vee \otimes_{\mathcal{O}_X}^\mathbf{L} L),
\ \ldots,
\ H^0(K^\vee \otimes_{\mathcal{O}_X}^\mathbf{L} L)
$$
都为零。事实上，由所引引理和引理
\ref{spaces-perfect-lemma-application-nice-K-injective}，有
$$
H^0(X, K^\vee \otimes_{\mathcal{O}_X}^\mathbf{L} L) =
H^0(X, \tau_{\geq -N}(K^\vee \otimes_{\mathcal{O}_X}^\mathbf{L} L))
$$
而由上同调层的消失性，它等于
$H^0(X, \tau_{\geq 1}(K^\vee \otimes_{\mathcal{O}_X}^\mathbf{L} L))$；
后者由《导出范畴》引理
\ref{derived-lemma-negative-vanishing} 为零。
因此，$\Hom_{D(\mathcal{O}_X)}(K, L)$ 为零，只要
$H^i(L) = 0$ 对于 $i \in [-b - N, -a]$ 成立。
\end{proof}

\noindent
下述定理是《概形的导出范畴》定理
\ref{perfect-theorem-DQCoh-is-Ddga} 的类似结果。

\begin{theorem}
\label{spaces-perfect-theorem-DQCoh-is-Ddga}
设 $S$ 为概形。
设 $X$ 为 $S$ 上拟紧且拟分离的代数空间。
则存在微分分次代数 $(E, \text{d})$，其非零上同调群 $H^i(E)$
只有有限多个，并且 $D_\QCoh(\mathcal{O}_X)$
等价于 $D(E, \text{d})$。
\end{theorem}

\begin{proof}
设 $K^\bullet$ 为 $\mathcal{O}$-模的 K-内射复形，% P09-ERR-0441
它是完美的并生成 $D_\QCoh(\mathcal{O}_X)$。
由定理 \ref{spaces-perfect-theorem-bondal-van-den-Bergh} 和 K-内射消解的存在性，
这样的对象存在。我们将证明，取
$$
(E, \text{d}) = \Hom_{\text{Comp}^{dg}(\mathcal{O}_X)}(K^\bullet, K^\bullet)
$$
时定理成立，其中 $\text{Comp}^{dg}(\mathcal{O}_X)$ 是
$\mathcal{O}$-模复形的微分分次范畴。参见《微分分次代数》第
\ref{dga-section-variant-base-change} 节。由于 $K^\bullet$ 是 K-内射的，
对于所有 $n \in \mathbf{Z}$，有
\begin{equation}
\label{spaces-perfect-equation-E-is-OK}
H^n(E) = \Ext^n_{D(\mathcal{O}_X)}(K^\bullet, K^\bullet)
\end{equation}
由引理 \ref{spaces-perfect-lemma-ext-from-perfect-into-bounded-QCoh}，
这些 Ext 中只有有限多个非零。考虑《微分分次代数》引理
\ref{dga-lemma-tensor-with-complex-derived} 中的函子
$$
- \otimes_E^\mathbf{L} K^\bullet :
D(E, \text{d}) \longrightarrow D(\mathcal{O}_X)
$$
由于 $K^\bullet$ 完美，由命题
\ref{spaces-perfect-proposition-compact-is-perfect}，它给出
$D(\mathcal{O}_X)$ 中的紧对象。% P09-ERR-0442
结合 (\ref{spaces-perfect-equation-E-is-OK})，由《微分分次代数》引理
\ref{dga-lemma-fully-faithful-in-compact-case}，上述函子全忠实。
由《微分分次代数》引理
\ref{dga-lemma-tensor-with-complex-hom-adjoint}，% P09-ERR-0443
它具有右伴随
$$
R\Hom(K^\bullet, - ) : D(\mathcal{O}_X) \longrightarrow D(E, \text{d})
$$
由伴随函子的一般理论，该右伴随是左拟逆函子。另一方面，由引理
\ref{spaces-perfect-lemma-tensor-with-QCoh-complex}，我们得到
$$
- \otimes_E^\mathbf{L} K^\bullet :
D(E, \text{d}) \longrightarrow D_\QCoh(\mathcal{O}_X)
$$
而由我们选取 $K^\bullet$ 为 $D_\QCoh(\mathcal{O}_X)$ 的生成元，
该伴随限制到 $D_\QCoh(\mathcal{O}_X)$ 后的核为零。
形式论证表明，这给出所需等价；见《导出范畴》引理
\ref{derived-lemma-fully-faithful-adjoint-kernel-zero}。
\end{proof}

\begin{remark}[带支撑的变式]
\label{spaces-perfect-remark-DQCoh-is-Ddga-with-support}
设 $S$ 为概形。% P09-ERR-0444
设 $X$ 为拟紧且拟分离的代数空间。
设 $T \subset |X|$ 为闭子集，且 $|X| \setminus T$ 拟紧。
定理 \ref{spaces-perfect-theorem-DQCoh-is-Ddga} 的类似结果对
$D_{\QCoh, T}(\mathcal{O}_X)$ 成立。
这由与该定理证明完全相同的论证得出，其中使用引理
\ref{spaces-perfect-lemma-generator-with-support}、\ref{spaces-perfect-lemma-compact-is-perfect-with-support}
以及引理 \ref{spaces-perfect-lemma-tensor-with-QCoh-complex} 的带支撑变式。
如果以后需要这一结果，我们将在此给出精确陈述和详细证明。
\end{remark}

\begin{remark}[微分分次代数的唯一性]
\label{spaces-perfect-remark-independence-choice}
设 $X$ 为环 $R$ 上拟紧且拟分离的代数空间。
由定理 \ref{spaces-perfect-theorem-DQCoh-is-Ddga} 证明中的构造，
存在微分分次代数 $(A, \text{d})$，它定义在 $R$ 上，使得三角范畴
$D_\QCoh(X)$ $R$-线性等价于 $D(A, \text{d})$。
可以问：$(A, \text{d})$ 在何种程度上唯一？
答案（仅仅）略好于只说 $(A, \text{d})$ 在导出等价意义下确定。
具体地，设 $(B, \text{d})$ 是第二个这样的偶。则有
$$
(A, \text{d}) = \Hom_{\text{Comp}^{dg}(\mathcal{O}_X)}(K^\bullet, K^\bullet)
$$
以及
$$
(B, \text{d}) = \Hom_{\text{Comp}^{dg}(\mathcal{O}_X)}(L^\bullet, L^\bullet)
$$
其中 $K^\bullet$ 和 $L^\bullet$ 是 $\mathcal{O}_X$-模的 K-内射复形，
分别对应于 $D_\QCoh(\mathcal{O}_X)$ 的完美生成元。置
$$
\Omega = \Hom_{\text{Comp}^{dg}(\mathcal{O}_X)}(K^\bullet, L^\bullet)
\quad
\Omega' = \Hom_{\text{Comp}^{dg}(\mathcal{O}_X)}(L^\bullet, K^\bullet)
$$
则 $\Omega$ 是微分分次 $B^{opp} \otimes_R A$-模，
而 $\Omega'$ 是微分分次 $A^{opp} \otimes_R B$-模。
此外，等价
$$
D(A, \text{d}) \to D_\QCoh(\mathcal{O}_X) \to
D(B, \text{d})
$$
由函子 $- \otimes_A^\mathbf{L} \Omega'$ 给出，拟逆也类似。
于是我们处于《微分分次代数》评注
\ref{dga-remark-hochschild-cohomology} 的情形。
如果以后需要这一评注，我们将在此给出精确陈述和详细证明。
\end{remark}











\section{刻画拟凝聚复形，I}
\label{spaces-perfect-section-pseudo-coherent-hocolim}

\noindent
这一材料将在《空间态射进阶》第
\ref{spaces-more-morphisms-section-characterize-pseudo-coherent} 节中继续讨论。
我们可以把拟凝聚对象刻画为完美对象的逼近所成的导出同伦极限。
% P09-ERR-0446

\begin{lemma}
\label{spaces-perfect-lemma-pseudo-coherent-hocolim}
设 $S$ 为概形。
设 $X$ 为 $S$ 上拟紧且拟分离的代数空间。
设 $K \in D(\mathcal{O}_X)$。下列条件等价：
\begin{enumerate}
\item $K$ 是拟凝聚的；并且
\item $K = \text{hocolim} K_n$，其中
$K_n$ 完美，并且映射
$\tau_{\geq -n}K_n \to \tau_{\geq -n}K$ 对所有 $n$ 都是同构。
\end{enumerate}
\end{lemma}

\begin{proof}
蕴含 (2) $\Rightarrow$ (1) 在任意带环位点上都成立。
事实上，假设 (2) 成立。回忆导出范畴中的完美对象是拟凝聚的；见
《位点上的上同调》引理 \ref{sites-cohomology-lemma-perfect}。
于是由定义，$\tau_{\geq -n}K_n$ 是 $(-n + 1)$-拟凝聚的，
从而 $\tau_{\geq -n}K$ 是 $(-n + 1)$-拟凝聚的，
因而 $K$ 是 $(-n + 1)$-拟凝聚的。这对所有 $n$ 都成立，
所以 $K$ 是拟凝聚的；见《位点上的上同调》定义
\ref{sites-cohomology-definition-pseudo-coherent}。

\medskip\noindent
假设 (1) 成立。先选取逼近 $K_1 \to K$；这是
$(X, K, -2)$ 的逼近，并由完美复形 $K_1$ 给出；见定义
\ref{spaces-perfect-definition-approximation-holds}、\ref{spaces-perfect-definition-approximation}
和定理 \ref{spaces-perfect-theorem-approximation}。
归纳地假设已有
$$
K_1 \to K_2 \to \ldots \to K_n \to K
$$
其中 $K_i$ 完美，并且
$\tau_{\geq -i}K_i \to \tau_{\geq -i}K$ 是同构，
对所有 $1 \leq i \leq n$ 都成立。% P09-ERR-0447
于是选取 $a \leq b$，即引理
\ref{spaces-perfect-lemma-ext-from-perfect-into-bounded-QCoh} 中用于完美对象 $K_n$
的整数。
选取逼近 $K_{n + 1} \to K$，即
$(X, K, \min(a - 1, -n - 1))$ 的逼近。选取特异三角
$$
K_{n + 1} \to K \to C \to K_{n + 1}[1]
$$
于是可见 $C \in D_\QCoh(\mathcal{O}_X)$ 满足
$H^i(C) = 0$，对所有 $i \geq a$ 都成立。
因此，由 $a, b$ 的选取可知
$\Hom_{D(\mathcal{O}_X)}(K_n, C) = 0$。
故复合 $K_n \to K \to C$ 为零。因此，由《导出范畴》引理
\ref{derived-lemma-representable-homological}，
我们可以将 $K_n \to K$ 经 $K_{n + 1}$ 分解，从而完成归纳步骤。

\medskip\noindent
还须证明 $K = \text{hocolim} K_n$。
这由把《导出范畴》引理
\ref{derived-lemma-cohomology-of-hocolim} 应用于函子
$H^i( - ) : D(\mathcal{O}_X) \to \textit{Mod}(\mathcal{O}_X)$
以及我们对 $K_n$ 的选取得出。
\end{proof}

\begin{lemma}
\label{spaces-perfect-lemma-pseudo-coherent-hocolim-with-support}
设 $X$ 为拟紧且拟分离的概形。
设 $T \subset X$ 为闭子集，且 $X \setminus T$ 拟紧。
设 $K \in D(\mathcal{O}_X)$ 支撑在 $T$ 上。下列条件等价：
\begin{enumerate}
\item $K$ 是拟凝聚的；并且
\item $K = \text{hocolim} K_n$，其中
$K_n$ 完美、支撑在 $T$ 上，并且
$\tau_{\geq -n}K_n \to \tau_{\geq -n}K$ 对所有 $n$ 都是同构。
\end{enumerate}
\end{lemma}

\begin{proof}
本引理的证明与引理
\ref{spaces-perfect-lemma-pseudo-coherent-hocolim} 的证明完全相同，
唯一差别是在选取逼近时使用三元组 $(T, K, m)$。
\end{proof}











\section{再论凝聚化函子}
\label{spaces-perfect-section-better-coherator}

\noindent
在第 \ref{spaces-perfect-section-coherator} 节中，我们构造并研究了右伴随 $RQ_X$；
它所伴随的典范函子是
$D(\QCoh(\mathcal{O}_X)) \to D(\mathcal{O}_X)$。
它被构造成凝聚化函子
$Q_X : \textit{Mod}(\mathcal{O}_X) \to \QCoh(\mathcal{O}_X)$
的右导出延拓。本节研究包含函子
$$
D_\QCoh(\mathcal{O}_X) \longrightarrow D(\mathcal{O}_X)
$$
何时具有右伴随。如果这个右伴随存在，我们将把它
记作\footnote{这大概不是标准记号。不过，我们已经用 $Q_X$
表示凝聚化函子，并用 $RQ_X$ 表示它的导出延拓。}
$$
DQ_X :
D(\mathcal{O}_X) \longrightarrow D_\QCoh(\mathcal{O}_X)
$$
事实证明，拟紧且拟分离的代数空间具有这样的右伴随。

\begin{lemma}
\label{spaces-perfect-lemma-better-coherator}
设 $S$ 为概形。
设 $X$ 为 $S$ 上拟紧且拟分离的代数空间。
包含函子 $D_\QCoh(\mathcal{O}_X) \to D(\mathcal{O}_X)$ 具有右伴随。
\end{lemma}

\begin{proof}[第一种证明]
我们将使用引理 \ref{spaces-perfect-lemma-induction-principle} 中的归纳原理证明此事。
如果 $D(\QCoh(\mathcal{O}_X)) \to D_\QCoh(\mathcal{O}_X)$ 是等价，
则本引理成立，因为第 \ref{spaces-perfect-section-coherator} 节的函子 $RQ_X$
是函子 $D(\QCoh(\mathcal{O}_X)) \to D(\mathcal{O}_X)$ 的右伴随。
特别地，本引理对仿射代数空间成立；见引理
\ref{spaces-perfect-lemma-affine-coherator}。因此，只须证明如下命题：
如果 $(U \subset X, f : V \to X)$ 是初等特异方形，其中 $U$ 拟紧、
$V$ 仿射，并且本引理对 $U$、$V$ 和 $U \times_X V$ 成立，
那么本引理对 $X$ 成立。

\medskip\noindent
右伴随存在，当且仅当对每个对象 $K$，只要它属于
$D(\mathcal{O}_X)$，
我们都能找到特异三角 % P09-ERR-0448
$$
E' \to E \to K \to E'[1]
$$
它位于 $D(\mathcal{O}_X)$ 中，并且 $E'$ 属于
$D_\QCoh(\mathcal{O}_X)$，同时 $\Hom(M, K) = 0$
对所有 $M$（属于 $D_\QCoh(\mathcal{O}_X)$）成立。见《导出范畴》引理
\ref{derived-lemma-right-adjoint}。考虑引理
\ref{spaces-perfect-lemma-exact-sequence-j-star} 的特异三角
$$
E \to Rj_{U, *}E|_U \oplus Rj_{V, *}E|_V \to
Rj_{U \times_X V, *}E|_{U \times_X V} \to E[1]
$$
它位于 $D(\mathcal{O}_X)$ 中。
由《导出范畴》引理 \ref{derived-lemma-prepare-adjoint}，只须为
$Rj_{U, *}E|_U$、$Rj_{V, *}E|_V$ 和
$Rj_{U \times_X V, *}E|_{U \times_X V}$
构造所需的特异三角。这把我们归约到下一段讨论的命题。

\medskip\noindent
设 $j : U \to X$ 为一个对应于如下情形的 \'etale 态射：
$U$ 拟紧且拟分离，并且本引理对 $U$ 成立。% P09-ERR-0449
设 $L$ 为 $D(\mathcal{O}_U)$ 的对象。
则存在特异三角
$$
E' \to Rj_*L \to K \to E'[1]
$$
它位于 $D(\mathcal{O}_X)$ 中，使得 $E'$ 属于
$D_\QCoh(\mathcal{O}_X)$，并且 $\Hom(M, K) = 0$ 对所有
$M$（属于 $D_\QCoh(\mathcal{O}_X)$）成立。
为说明这一点，选取特异三角
$$
L' \to L \to Q \to L'[1]
$$
它位于 $D(\mathcal{O}_U)$ 中，
使得 $L'$ 属于 $D_\QCoh(\mathcal{O}_U)$，并且
$\Hom(N, Q) = 0$ 对所有 $N$（属于
$D_\QCoh(\mathcal{O}_U)$）成立。
这是可能的，因为《导出范畴》引理
\ref{derived-lemma-right-adjoint} 中的陈述是一个当且仅当命题。
我们得到特异三角
$$
Rj_*L' \to Rj_*L \to Rj_*Q \to Rj_*L'[1]
$$
它位于 $D(\mathcal{O}_X)$ 中。
注意，由引理 \ref{spaces-perfect-lemma-quasi-coherence-direct-image}，
$Rj_*L'$ 属于 $D_\QCoh(\mathcal{O}_X)$。
另一方面，如果 $M$ 属于 $D_\QCoh(\mathcal{O}_X)$，% P09-ERR-0450
则
$$
\Hom(M, Rj_*Q) = \Hom(Lj^*M, Q) = 0
$$
因为由引理 \ref{spaces-perfect-lemma-quasi-coherence-pullback}，
$Lj^*M$ 属于 $D_\QCoh(\mathcal{O}_U)$。
这完成了证明。
\end{proof}

\begin{proof}[第二种证明]
由《导出范畴》命题 \ref{derived-proposition-brown}，该伴随存在。
其假设均满足：首先，注意 $D_\QCoh(\mathcal{O}_X)$ 有直和，
并且取直和与包含函子可交换（引理
\ref{spaces-perfect-lemma-quasi-coherence-direct-sums}）。
另一方面，$D_\QCoh(\mathcal{O}_X)$ 是紧生成的，因为它有完美生成元；
见定理 \ref{spaces-perfect-theorem-bondal-van-den-Bergh}，% P09-ERR-0451
并且由命题 \ref{spaces-perfect-proposition-compact-is-perfect}，完美对象是紧的。
\end{proof}

\begin{lemma}
\label{spaces-perfect-lemma-pushforward-better-coherator}
设 $S$ 为概形。
设 $f : X \to Y$ 为 $S$ 上代数空间的拟紧且拟分离态射。
如果右伴随 $DQ_X$ 和 $DQ_Y$——即包含函子
$D_\QCoh \to D$ 对 $X$ 和 $Y$ 的右伴随——存在，则
$$
Rf_* \circ DQ_X = DQ_Y \circ Rf_*
$$
\end{lemma}

\begin{proof}
该陈述有意义，因为由引理
\ref{spaces-perfect-lemma-quasi-coherence-direct-image}，$Rf_*$
把 $D_\QCoh(\mathcal{O}_X)$ 映入 $D_\QCoh(\mathcal{O}_Y)$。
该陈述成立，是因为 $Lf^*$ 同样把
$D_\QCoh(\mathcal{O}_Y)$ 映入 $D_\QCoh(\mathcal{O}_X)$
（引理 \ref{spaces-perfect-lemma-quasi-coherence-pullback}），
因而 $Rf_* \circ DQ_X$ 和 $DQ_Y \circ Rf_*$
都是 $Lf^* : D_\QCoh(\mathcal{O}_Y) \to D(\mathcal{O}_X)$ 的右伴随。
\end{proof}

\begin{remark}
\label{spaces-perfect-remark-explain-consequence}
设 $S$ 为概形。设 $(U \subset X, f : V \to X)$ 为
$S$ 上代数空间的初等特异方形。
假设 $X$、$U$、$V$ 拟紧且拟分离。
由引理 \ref{spaces-perfect-lemma-better-coherator}，函子
$DQ_X$、$DQ_U$、$DQ_V$、$DQ_{U \times_X V}$ 存在。
此外，有典范特异三角
$$
DQ_X(K) \to Rj_{U, *}DQ_U(K|_U) \oplus Rj_{V, *}DQ_V(K|_V)
\to Rj_{U \times_X V, *}DQ_{U \times_X V}(K|_{U \times_X V}) \to
$$
对任意 $K \in D(\mathcal{O}_X)$ 都成立。
% P09-ERR-0452
这是把正合函子 $DQ_X$ 应用于引理
\ref{spaces-perfect-lemma-exact-sequence-j-star} 的特异三角，
并使用三次引理 \ref{spaces-perfect-lemma-pushforward-better-coherator} 得出的。
\end{remark}

\begin{lemma}
\label{spaces-perfect-lemma-boundedness-better-coherator}
设 $S$ 为概形。
设 $X$ 为 $S$ 上拟紧且拟分离的代数空间。
引理 \ref{spaces-perfect-lemma-better-coherator} 的函子 $DQ_X$
具有如下有界性：存在整数 $N = N(X)$，使得，如果
$K$ 属于 $D(\mathcal{O}_X)$，并且 $H^i(U, K) = 0$，
其中 $U$ 是在 $X$ 上仿射且 \'etale 的，而
$i \not \in [a, b]$，那么上同调层 $H^i(DQ_X(K))$ 为零，
只要 $i \not \in [a, b + N]$。
\end{lemma}

\begin{proof}
我们将使用引理 \ref{spaces-perfect-lemma-induction-principle} 的归纳原理证明此事。

\medskip\noindent
如果 $X$ 仿射，则本引理对 $N = 0$ 成立，因为此时
$RQ_X = DQ_X$ 由取与 $R\Gamma(X, K)$ 相伴的拟凝聚层复形给出。
见引理 \ref{spaces-perfect-lemma-affine-coherator}。

\medskip\noindent
设 $(U \subset W, f : V \to W)$ 为初等特异方形，其中
$W$ 拟紧且拟分离，$U \subset W$ 为拟紧开子空间，$V$ 仿射，
并且本引理对 $U$、$V$ 和 $U \times_W V$ 成立。
设相应整数为 $N(U)$、$N(V)$ 和 $N(U \times_W V)$。
现在假设 $K$ 属于 $D(\mathcal{O}_X)$，并且
$H^i(W, K) = 0$，其中 $W$ 是在 $X$ 上仿射且 \'etale 的，
而 $i \not \in [a, b]$。% P09-ERR-0453
则 $K|_U$、$K|_V$、$K|_{U \times_W V}$ 具有相同性质。
因此可见 $RQ_U(K|_U)$、$RQ_V(K|_V)$ 和 % P09-ERR-0454
$RQ_{U \cap V}(K|_{U \times_W V})$ % P09-ERR-0455
在区间
$[a, b + \max(N(U), N(V), N(U \times_W V))$ % P09-ERR-0456
之外具有消失的上同调层。
由于函子 $Rj_{U, *}$、$Rj_{V, *}$、$Rj_{U \times_W V, *}$
由引理 \ref{spaces-perfect-lemma-quasi-coherence-direct-image} 在
$D_\QCoh$ 上具有有限上同调维数，可见存在 $N$，使得
$Rj_{U, *}DQ_U(K|_U)$、$Rj_{V, *}DQ_V(K|_V)$ 和
$Rj_{U \cap V, *}DQ_{U \times_W V}(K|_{U \times_W V})$
在区间 $[a, b + N]$ 之外具有消失的上同调层。
最后，由评注 \ref{spaces-perfect-remark-explain-consequence} 的特异三角得出结论。
\end{proof}

\begin{example}
\label{spaces-perfect-example-inverse-limit-quasi-coherent}
设 $S$ 为概形。
设 $X$ 为 $S$ 上拟紧且拟分离的代数空间。
设 $(\mathcal{F}_n)$ 为 $X$ 上拟凝聚层的逆系统。
由于 $DQ_X$ 是右伴随，它与积可交换，因而与导出极限可交换。
因此
$$
DQ_X(R\lim \mathcal{F}_n) =
(R\lim\text{ 于 }D_\QCoh(\mathcal{O}_X))(\mathcal{F}_n)
$$
其中第一个 $R\lim$ 取于 $D(\mathcal{O}_X)$ 中。
事实上，为此记 $K = R\lim \mathcal{F}_n$。
对任意 $U$，只要它在 $X$ 上仿射且 \'etale，就有
$$
H^i(U, K) =
H^i(R\Gamma(U, R\lim \mathcal{F}_n)) =
H^i(R\lim R\Gamma(U, \mathcal{F}_n)) =
H^i(R\lim \Gamma(U, \mathcal{F}_n))
$$
因为上同调与导出极限可交换，并且拟凝聚层
$\mathcal{F}_n$ 在仿射对象上没有高阶上同调。
由阿贝尔群范畴中对 $R\lim$ 的计算可见，
$H^i(U, K) = 0$，除非 $i \in [0, 1]$。
最后断定，$R\lim$（取于 $D_\QCoh(\mathcal{O}_X)$ 中），
即上文的 $DQ_X(K)$，属于 $D^b_\QCoh(\mathcal{O}_X)$，
并且由引理 \ref{spaces-perfect-lemma-boundedness-better-coherator}，
其负次数上同调层为零。
\end{example}










\section{上同调与基变换，IV}
\label{spaces-perfect-section-cohomology-and-base-change-perfect}

\noindent
本节是《概形的导出范畴》第
\ref{perfect-section-cohomology-and-base-change-perfect} 节的类似结果。

\begin{lemma}
\label{spaces-perfect-lemma-cohomology-base-change}
设 $S$ 为概形。设 $f : X \to Y$ 为 $S$ 上代数空间的
拟紧且拟分离态射。对 $E$ 属于
$D_\QCoh(\mathcal{O}_X)$ 且 $K$ 属于
$D_\QCoh(\mathcal{O}_Y)$，有
$$
Rf_*(E) \otimes_{\mathcal{O}_Y}^\mathbf{L} K =
Rf_*(E \otimes_{\mathcal{O}_X}^\mathbf{L} Lf^*K)
$$
\end{lemma}

\begin{proof}
不加任何假设时，存在映射
$Rf_*(E) \otimes_{\mathcal{O}_Y}^\mathbf{L} K \to
Rf_*(E \otimes_{\mathcal{O}_X}^\mathbf{L} Lf^*K)$。
具体地，它是典范映射
$$
Lf^*(Rf_*(E) \otimes_{\mathcal{O}_Y}^\mathbf{L} K) =
Lf^*(Rf_*(E)) \otimes_{\mathcal{O}_X}^\mathbf{L} Lf^*K
\longrightarrow
E \otimes_{\mathcal{O}_X}^\mathbf{L} Lf^*K
$$
的伴随映射，而该映射来自 $Lf^*Rf_*E \to E$。见《位点上的上同调》引理
\ref{sites-cohomology-lemma-pullback-tensor-product} 和
\ref{sites-cohomology-lemma-adjoint}。
为检验它是同构，可以在 $Y$ 上 \'etale 局部地工作。
因此归约到 $Y$ 为仿射概形的情形。

\medskip\noindent
假设 $K = \bigoplus K_i$ 是若干复形
$K_i \in D_\QCoh(\mathcal{O}_Y)$ 的直和。
如果该陈述对每个 $K_i$ 成立，那么它对 $K$ 成立。
事实上，函子 $Lf^*$ 和 $\otimes^\mathbf{L}$ 由构造保持直和，
而由引理 \ref{spaces-perfect-lemma-quasi-coherence-pushforward-direct-sums}，
$Rf_*$ 与直和可交换（对具有拟凝聚上同调层的复形）。
此外，假设 $K \to L \to M \to K[1]$ 是
$D_\QCoh(Y)$ 中的特异三角。那么，如果本引理的陈述对
$K, L, M$ 中两个成立，它就对第三个成立
（因为所涉函子都是三角范畴的正合函子）。

\medskip\noindent
假设 $Y$ 仿射，记 $Y = \Spec(A)$。函子
$\widetilde{\ } : D(A) \to D_\QCoh(\mathcal{O}_Y)$ 是等价；见引理
\ref{spaces-perfect-lemma-derived-quasi-coherent-small-etale-site} 和
《概形的导出范畴》引理
\ref{perfect-lemma-affine-compare-bounded}。
令 $T$ 表示关于 $K \in D(A)$ 的如下性质：
本引理的陈述对 $\widetilde{K}$ 成立。
上面的讨论和《代数进阶》评注
\ref{more-algebra-remark-P-resolution} 表明，% P09-ERR-0457
只须证明 $T$ 对 $A[k]$ 成立。这完成了证明，因为本引理的陈述
对结构层的平移显然成立。
\end{proof}

\begin{definition}
\label{spaces-perfect-definition-tor-independent}
设 $S$ 为概形，$B$ 为 $S$ 上的代数空间。
设 $X$、$Y$ 为 $B$ 上的代数空间。称 $X$ 和
$Y$ 在 $B$ 上 {\it Tor 独立}，当且仅当对每个交换图
$$
\xymatrix{
\Spec(k) \ar[d]_{\overline{y}} \ar[dr]_{\overline{b}} \ar[r]_-{\overline{x}} &
X \ar[d] \\
Y \ar[r] & B
}
$$
的几何点，环 $\mathcal{O}_{X, \overline{x}}$ 和
$\mathcal{O}_{Y, \overline{y}}$ 在
$\mathcal{O}_{B, \overline{b}}$ 上 Tor 独立（见《代数进阶》定义
\ref{more-algebra-definition-tor-independent}）。
\end{definition}

\noindent
下述引理特别表明：在代数空间可表的情形中，本定义与我们已有的定义一致。

\begin{lemma}
\label{spaces-perfect-lemma-tor-independent}
设 $S$ 为概形，$B$ 为 $S$ 上的代数空间。
设 $X$、$Y$ 为 $B$ 上的代数空间。下列条件等价：
\begin{enumerate}
\item $X$ 和 $Y$ 在 $B$ 上 Tor 独立；
\item 对每个交换图
$$
\xymatrix{
U \ar[d] \ar[r] & W \ar[d] & V \ar[d] \ar[l] \\
X \ar[r] & B & Y \ar[l]
}
$$
若竖直箭头为 \'etale，则 $U$ 和 $V$ 在 $W$ 上 Tor 独立；% P09-ERR-0458
\item 对某个如 (2) 中的交换图，满足 (a) $W \to B$ 为
\'etale 满射，(b) $U \to X \times_B W$ 为 \'etale 满射，
(c) $V \to Y \times_B W$ 为 \'etale 满射，空间 $U$ 和 $V$
在 $W$ 上 Tor 独立；并且
\item 对某个如 (3) 中的交换图，其中 $U$、$V$、$W$ 是概形，
概形 $U$ 和 $V$ 在 $W$ 上 Tor 独立，其意义见《概形的导出范畴》定义
\ref{perfect-definition-tor-independent}。
\end{enumerate}
\end{lemma}

\begin{proof}
对代数空间的 \'etale 态射 $\varphi : U \to X$ 以及几何点
$\overline{u}$，局部环映射
$\mathcal{O}_{X, \varphi(\overline{u})} \to
\mathcal{O}_{U, \overline{u}}$ 是同构。
因此 (1) 与 (2) 的等价性成立，蕴含 (1) $\Rightarrow$ (3)
也成立。假设 (3)，并选取一个如定义
\ref{spaces-perfect-definition-tor-independent} 中的几何点图。
这些假设表明，我们可以先把 $\overline{b}$ 提升到几何点
$\overline{w}$（属于 $W$），再把几何点
$(\overline{x}, \overline{b})$ 提升到几何点
$\overline{u}$（属于 $U$），最后把几何点
$(\overline{y}, \overline{b})$ 提升到几何点
$\overline{v}$（属于 $V$）。使用《空间的性质》引理
\ref{spaces-properties-lemma-geometric-lift-to-usual} 找到这些提升。
利用上面关于局部环的评注，可断定给定的图满足定义中的条件。

\medskip\noindent
作出这些初步选择后，显然，(4) 归结为如下陈述：
当 $X$、$Y$、$B$ 是概形时，定义
\ref{spaces-perfect-definition-tor-independent} 与《概形的导出范畴》定义
\ref{perfect-definition-tor-independent} 一致。

\medskip\noindent
设 $\overline{x}, \overline{b}, \overline{y}$ 如定义
\ref{spaces-perfect-definition-tor-independent} 中那样，位于点
$x, y, b$ 之上。回忆
$\mathcal{O}_{X, \overline{x}} = \mathcal{O}_{X, x}^{sh}$
（《空间的性质》引理
\ref{spaces-properties-lemma-describe-etale-local-ring}），其余两个也类似。
由《代数》引理
\ref{algebra-lemma-strictly-henselian-functorial-improve} 可见，
$\mathcal{O}_{X, \overline{x}}$ 是
$\mathcal{O}_{X, x} \otimes_{\mathcal{O}_{B, b}}
\mathcal{O}_{B, \overline{b}}$ 的严格亨泽尔化。
特别地，环映射
$$
\mathcal{O}_{X, x} \otimes_{\mathcal{O}_{B, b}} \mathcal{O}_{B, \overline{b}}
\longrightarrow
\mathcal{O}_{X, \overline{x}}
$$
是平坦的（《代数进阶》引理
\ref{more-algebra-lemma-dumb-properties-henselization}）。
由《代数进阶》引理
\ref{more-algebra-lemma-tor-independent-flat} 可见
$$
\text{Tor}_i^{\mathcal{O}_{B, b}}(\mathcal{O}_{X, x}, \mathcal{O}_{Y, y})
\otimes_{\mathcal{O}_{X, x} \otimes_{\mathcal{O}_{B, b}} \mathcal{O}_{Y, y}}
(\mathcal{O}_{X, \overline{x}} \otimes_{\mathcal{O}_{B, \overline{b}}}
\mathcal{O}_{Y, \overline y})
=
\text{Tor}_i^{\mathcal{O}_{B, \overline{b}}}(
\mathcal{O}_{X, \overline{x}}, \mathcal{O}_{Y, \overline{y}})
$$
因此，如果作为概形的 $X$ 和 $Y$ 在 $B$ 上 Tor 独立，
那么作为代数空间的 $X$ 和 $Y$ 在 $B$ 上 Tor 独立。

\medskip\noindent
反之，可以假设 $X$、$Y$、$B$ 仿射。注意，环映射
$$
\mathcal{O}_{X, x} \otimes_{\mathcal{O}_{B, b}} \mathcal{O}_{Y, y}
\longrightarrow
\mathcal{O}_{X, \overline{x}} \otimes_{\mathcal{O}_{B, \overline{b}}}
\mathcal{O}_{Y, \overline y}
$$
由上面的观察是平坦的。此外，该映射在谱上的像包含所有素理想
$$
\mathfrak s \subset
\mathcal{O}_{X, x} \otimes_{\mathcal{O}_{B, b}} \mathcal{O}_{Y, y}
$$
只要它们位于 $\mathfrak m_x$ 和 $\mathfrak m_y$ 之上。
因此，由此及上面显示的 Tor 公式可见，如果作为代数空间的
$X$ 和 $Y$ 在 $B$ 上 Tor 独立，那么
$$
\text{Tor}_i^{\mathcal{O}_{B, b}}
(\mathcal{O}_{X, x}, \mathcal{O}_{Y, y})_\mathfrak s = 0
$$
对所有 $i > 0$ 以及上面所有的 $\mathfrak s$ 成立。
把《代数进阶》引理
\ref{more-algebra-lemma-tor-independent} 应用于环映射
$\Gamma(B, \mathcal{O}_B) \to \Gamma(X, \mathcal{O}_X)$
以及
$\Gamma(B, \mathcal{O}_B) \to \Gamma(X, \mathcal{O}_X)$，% P09-ERR-0459
可知这蕴含 $X$ 和 $Y$ 在 $B$ 上 Tor 独立。
\end{proof}

\begin{lemma}
\label{spaces-perfect-lemma-compare-base-change}
设 $S$ 为概形。设 $g : Y' \to Y$ 为 $S$ 上代数空间的态射。
设 $f : X \to Y$ 为 $S$ 上代数空间的拟紧且拟分离态射。
考虑基变换图
$$
\xymatrix{
X' \ar[r]_{g'} \ar[d]_{f'} &
X \ar[d]^f \\
Y' \ar[r]^g &
Y
}
$$
如果 $X$ 和 $Y'$ 在 $Y$ 上 Tor 独立，那么对所有
$E \in D_\QCoh(\mathcal{O}_X)$，有
$Rf'_*L(g')^*E = Lg^*Rf_*E$。
\end{lemma}

\begin{proof}
对任意对象 $E$（属于 $D(\mathcal{O}_X)$），可以使用
《位点上的上同调》评注
\ref{sites-cohomology-remark-base-change} 得到典范基变换映射
$Lg^*Rf_*E \to Rf'_*L(g')^*E$。为检验它是同构，可以在
$Y'$ 上 \'etale 局部地工作。因此可以假设
$g : Y' \to Y$ 是仿射概形的态射。特别地，$g$ 是仿射的，
并且只须证明
$$
Rg_*Lg^*Rf_*E \to Rg_*Rf'_*L(g')^*E = Rf_*(Rg'_* L(g')^* E)
$$
是同构；见引理 \ref{spaces-perfect-lemma-affine-morphism}
（并使用引理 \ref{spaces-perfect-lemma-quasi-coherence-pullback}、
\ref{spaces-perfect-lemma-quasi-coherence-tensor-product} 和
\ref{spaces-perfect-lemma-quasi-coherence-direct-image} 来说明对象
$Rf'_*L(g')^*E$ 和 $Lg^*Rf_*E$ 具有拟凝聚上同调层）。
注意，$g'$ 也是仿射的（《空间的态射》引理
\ref{spaces-morphisms-lemma-base-change-affine}）。
由引理 \ref{spaces-perfect-lemma-affine-morphism-pull-push}，该映射变为映射
$$
Rf_*E \otimes_{\mathcal{O}_Y}^\mathbf{L} g_*\mathcal{O}_{Y'}
\longrightarrow
Rf_*(E \otimes_{\mathcal{O}_X}^\mathbf{L} g'_*\mathcal{O}_{X'})
$$
注意 $g'_*\mathcal{O}_{X'} = f^*g_*\mathcal{O}_{Y'}$。
因此由引理 \ref{spaces-perfect-lemma-cohomology-base-change}，只须证明
$Lf^*g_*\mathcal{O}_{Y'} = f^*g_*\mathcal{O}_{Y'}$。
这由假设 $X$ 和 $Y'$ 在 $Y$ 上 Tor 独立得出。
事实上，为检验这一点，可以在 $X$ 上 \'etale 局部地工作，
因此还可以假设 $X$ 仿射。记
$X = \Spec(A)$、$Y = \Spec(R)$ 且 $Y' = \Spec(R')$。
我们的假设蕴含 $A$ 和 $R'$ 在 $R$ 上 Tor 独立（见引理
\ref{spaces-perfect-lemma-tor-independent} 和《代数进阶》引理
\ref{more-algebra-lemma-tor-independent}），即
$\text{Tor}_i^R(A, R') = 0$ 对 $i > 0$ 成立。换言之，
$A \otimes_R^\mathbf{L} R' = A \otimes_R R'$，
这恰好意味着
$Lf^*g_*\mathcal{O}_{Y'} = f^*g_*\mathcal{O}_{Y'}$。
\end{proof}

\noindent
下述引理将在对偶化复形一章中使用。

\begin{lemma}
\label{spaces-perfect-lemma-affine-morphism-and-hom-out-of-perfect}
设 $g : S' \to S$ 为仿射概形的态射。考虑笛卡尔方形
$$
\xymatrix{
X' \ar[r]_{g'} \ar[d]_{f'} & X \ar[d]^f \\
S' \ar[r]^g & S
}
$$
其中各代数空间拟紧且拟分离。假设 $g$ 和 $f$ Tor 独立。
写作 $S = \Spec(R)$ 且 $S' = \Spec(R')$。对
$M, K \in D(\mathcal{O}_X)$，典范映射
$$
R\Hom_X(M, K) \otimes^\mathbf{L}_R R'
\longrightarrow
R\Hom_{X'}(L(g')^*M, L(g')^*K)
$$
在 $D(R')$ 中，于下列两种情形为同构：
\begin{enumerate}
\item $M \in D(\mathcal{O}_X)$ 完美，且 $K \in D_\QCoh(X)$；或
\item $M \in D(\mathcal{O}_X)$ 拟凝聚，
$K \in D_\QCoh^+(X)$，且 $R'$ 在 $R$ 上具有有限 Tor 维数。
\end{enumerate}
\end{lemma}

\begin{proof}
存在典范映射
$R\Hom_X(M, K) \to R\Hom_{X'}(L(g')^*M, L(g')^*K)$，
它位于全局 Hom 复形的 $D(\Gamma(X, \mathcal{O}_X))$ 中；见
《位点上的上同调》第
\ref{sites-cohomology-section-global-RHom} 节。
限制标量后，可以把它看作 $D(R)$ 中的映射。
然后使用限制标量与 $- \otimes_R^\mathbf{L} R'$ 的伴随性，
得到本引理中显示的映射。定义了该映射后，只须证明它在
阿贝尔群的导出范畴中是同构。

\medskip\noindent
右端等于
$$
R\Hom_X(M, R(g')_*L(g')^*K) =
R\Hom_X(M, K \otimes_{\mathcal{O}_X}^\mathbf{L} g'_*\mathcal{O}_{X'})
$$
这是引理 \ref{spaces-perfect-lemma-affine-morphism-pull-push} 的结论。
在两种情形中，由引理 \ref{spaces-perfect-lemma-quasi-coherence-internal-hom}，
复形 $R\SheafHom(M, K)$ 都是
$D_\QCoh(\mathcal{O}_X)$ 的对象。存在自然映射
$$
R\SheafHom(M, K) \otimes_{\mathcal{O}_X}^\mathbf{L} g'_*\mathcal{O}_{X'}
\longrightarrow
R\SheafHom(M, K \otimes_{\mathcal{O}_X}^\mathbf{L} g'_*\mathcal{O}_{X'})
$$
它在两种情形中都是同构；引理
\ref{spaces-perfect-lemma-internal-hom-evaluate-tensor-isomorphism}。% P09-ERR-0461
为说明该引理适用于情形 (2)，注意
$g'_*\mathcal{O}_{X'} = Rg'_*\mathcal{O}_{X'} =
Lf^*g_*\mathcal{O}_X$，其中第二个等号由引理
\ref{spaces-perfect-lemma-compare-base-change} 给出。% P09-ERR-0462
使用《概形的导出范畴》引理
\ref{perfect-lemma-tor-dimension-affine}、引理
\ref{spaces-perfect-lemma-descend-tor-amplitude} 和《位点上的上同调》引理
\ref{sites-cohomology-lemma-tor-amplitude-pullback}，可断定
$g'_*\mathcal{O}_{X'}$ 具有有限 Tor 维数。
因此，在两种情形中，把 $K$ 替换为 $R\SheafHom(M, K)$，
我们都归约到证明
$$
R\Gamma(X, K) \otimes^\mathbf{L}_A A' \longrightarrow
R\Gamma(X, K \otimes^\mathbf{L}_{\mathcal{O}_X} g'_*\mathcal{O}_{X'})
$$
是同构。% P09-ERR-0463
注意，由引理 \ref{spaces-perfect-lemma-affine-morphism-pull-push}，左端等于
$R\Gamma(X', L(g')^*K)$。因此结论由引理
\ref{spaces-perfect-lemma-compare-base-change} 得出。
\end{proof}

\begin{remark}
\label{spaces-perfect-remark-multiplication-map}
沿用引理 \ref{spaces-perfect-lemma-affine-morphism-and-hom-out-of-perfect} 的记号。% P09-ERR-0460
下图交换：
$$
\xymatrix{
R\Hom_X(M, Rg'_*L) \otimes_R^\mathbf{L} R' \ar[r] \ar[d]_\mu &
R\Hom_{X'}(L(g')^*M, L(g')^*Rg'_*L) \ar[d]^a \\
R\Hom_X(M, R(g')_*L) \ar@{=}[r] &
R\Hom_{X'}(L(g')^*M, L)
}
$$
其中上方水平箭头是该引理中的映射，$\mu$ 是乘法映射，
而 $a$ 来自伴随映射 $L(g')^*Rg'_*L \to L$。
乘法映射是伴随映射
$K' \otimes_R^\mathbf{L} R' \to K'$，其中
$K' \in D(R')$ 任意。
\end{remark}

\begin{lemma}
\label{spaces-perfect-lemma-tor-independence-and-tor-amplitude}
设 $S$ 为概形。考虑代数空间的笛卡尔方形
$$
\xymatrix{
X' \ar[r]_{g'} \ar[d]_{f'} & X \ar[d]^f \\
Y' \ar[r]^g & Y
}
$$
它位于 $S$ 上。
假设 $g$ 和 $f$ Tor 独立。
\begin{enumerate}
\item 如果 $E \in D(\mathcal{O}_X)$ 具有 $[a, b]$ 中的
Tor 振幅（把它视为 $f^{-1}\mathcal{O}_Y$-模复形），
那么 $L(g')^*E$ 具有 $[a, b]$ 中的 Tor 振幅
（把它视为 $f^{-1}\mathcal{O}_{Y'}$-模复形）。% P09-ERR-0464
\item 如果 $\mathcal{G}$ 是 $\mathcal{O}_X$-模且在
$Y$ 上平坦，那么
$L(g')^*\mathcal{G} = (g')^*\mathcal{G}$。
\end{enumerate}
\end{lemma}

\begin{proof}
可以在茎上计算 Tor 维数；见《位点上的上同调》引理
\ref{sites-cohomology-lemma-tor-amplitude-stalk} 和《空间的性质》定理
\ref{spaces-properties-theorem-exactness-stalks}。
如果 $\overline{x}'$ 是 $X'$ 的几何点，且 $\overline{x}$
是它在 $X$ 中的像，那么
$$
(L(g')^*E)_{\overline{x}'} =
E_{\overline{x}}
\otimes_{\mathcal{O}_{X, \overline{x}}}^\mathbf{L}
\mathcal{O}_{X', \overline{x}'}
$$
设 $\overline{y}'$（属于 $Y'$）和 $\overline{y}$（属于 $Y$）
分别是 $\overline{x}'$ 和 $\overline{x}$ 的像。
由于 $X$ 和 $Y'$ 在 $Y$ 上 Tor 独立，可以应用《代数进阶》引理
\ref{more-algebra-lemma-base-change-comparison}，看出所显示公式的右端等于
$E_{\overline{x}}
\otimes_{\mathcal{O}_{Y, \overline{y}}}^\mathbf{L}
\mathcal{O}_{Y', \overline{y}'}$
（在 $D(\mathcal{O}_{Y', \overline{y}'})$ 中）。
因此 (1) 由《代数进阶》引理
\ref{more-algebra-lemma-pull-tor-amplitude} 得出。
为证明 (2)，注意 $\mathcal{G}$ 的平坦性等价于
$\mathcal{G}[0]$ 具有 $[0, 0]$ 中的 Tor 振幅。
应用 (1) 即得结论。
\end{proof}










\section{上同调与基变换，V}
\label{spaces-perfect-section-cohomology-base-change}

\noindent
本节是《概形的导出范畴》第
\ref{perfect-section-cohomology-base-change} 节的类似结果。
在第 \ref{spaces-perfect-section-cohomology-and-base-change-perfect} 节中，
我们看到，当态射 Tor 独立时，基变换定理成立。
即使在仿射情形中，若无这种条件，也不可能有基变换定理；见
《代数进阶》第 \ref{more-algebra-section-tor-independence} 节。
本节分析何时可以“逐个复形”得到基变换结果。

\medskip\noindent
为此，设 $S$ 为基概形，并假设我们有交换图
$$
\xymatrix{
X' \ar[r]_{g'} \ar[d]_{f'} &
X \ar[d]^f \\
Y' \ar[r]^g &
Y
}
$$
它由 $S$ 上的代数空间组成（通常我们会假设它是笛卡尔图）。
设 $K \in D_\QCoh(\mathcal{O}_X)$，并设
$L(g')^*K \to K'$ 是 $D_\QCoh(\mathcal{O}_{X'})$ 中的映射。
对几何点 $\overline{x}'$（属于 $X'$），考虑几何点
$\overline{x} = g'(\overline{x}')$、
$\overline{y}' = f'(\overline{x}')$、
$\overline{y} = f(\overline{x}) = g(\overline{y}')$，
它们分别属于 $X$、$Y'$、$Y$。
于是可以考虑映射
$$
K_{\overline{x}}
\otimes_{\mathcal{O}_{Y, \overline{y}}}^\mathbf{L}
\mathcal{O}_{Y', \overline{y}'}
\to
K_{\overline{x}}
\otimes_{\mathcal{O}_{X, \overline{x}}}^\mathbf{L}
\mathcal{O}_{X', \overline{x}'} \to
K'_{\overline{x}'}
$$
其中第一个箭头是《代数进阶》方程
(\ref{more-algebra-equation-comparison-map})，第二个箭头来自
$$
(L(g')^*K)_{\overline{x}'} =
K_{\overline{x}} \otimes_{\mathcal{O}_{X, \overline{x}}}^\mathbf{L}
\mathcal{O}_{X', \overline{x}'}
$$
以及给定映射 $L(g')^*K \to K'$。
对每个 $i \in \mathbf{Z}$，得到一个
$\mathcal{O}_{X, \overline{x}} \otimes_{\mathcal{O}_{Y, \overline{y}}}
\mathcal{O}_{Y', \overline{y}'}$-模结构，其承载模为
$H^i(K_{\overline{x}}
\otimes_{\mathcal{O}_{Y, \overline{y}}}^\mathbf{L}
\mathcal{O}_{Y', \overline{y}'})$。
把所有这些放在一起，得到典范映射
\begin{equation}
\label{spaces-perfect-equation-bc}
H^i(K_{\overline{x}}
\otimes_{\mathcal{O}_{Y, \overline{y}}}^\mathbf{L}
\mathcal{O}_{Y', \overline{y}'})
\otimes_{(\mathcal{O}_{X, \overline{x}}
\otimes_{\mathcal{O}_{Y, \overline{y}}}
\mathcal{O}_{Y', \overline{y}'})}
\mathcal{O}_{X', \overline{x}'}
\longrightarrow
H^i(K'_{\overline{x}'})
\end{equation}
它是 $\mathcal{O}_{X', \overline{x}'}$-模的映射。

\begin{lemma}
\label{spaces-perfect-lemma-single-complex-base-change-condition}
设 $S$ 为概形。设
$$
\xymatrix{
X' \ar[r]_{g'} \ar[d]_{f'} &
X \ar[d]^f \\
Y' \ar[r]^g &
Y
}
$$
为 $S$ 上代数空间的笛卡尔图。
设 $K \in D_\QCoh(\mathcal{O}_X)$，并设 $L(g')^*K \to K'$
是 $D_\QCoh(\mathcal{O}_{X'})$ 中的映射。下列条件等价：
\begin{enumerate}
\item 对任意 $x' \in X'$ 和 $i \in \mathbf{Z}$，% P09-ERR-0465
映射 (\ref{spaces-perfect-equation-bc}) 是同构；
\item 对任意交换图
$$
\xymatrix{
& U \ar[d] \ar[rd]^a \\
V' \ar[r] \ar[rd]^c & V \ar[rd]^b & X \ar[d]^f \\
& Y' \ar[r]^g & Y
}
$$
若 $a, b, c$ 为 \'etale，$U, V, V'$ 为概形，并且
$U' = V' \times_V U$，则《概形的导出范畴》引理
\ref{perfect-lemma-single-complex-base-change-condition} 的等价条件
对 $(U \to X)^*K$ 和 $(U' \to X')^*K'$ 成立；并且
\item 存在一个如 (2) 中的图，使得 $U' \to X'$ 满射。
\end{enumerate}
\end{lemma}

\begin{proof}
注意，(1) 在 $X'$ 上是 \'etale 局部的。
逐一考察已知条件之间的形式蕴含，可见只须证明：当
$X, Y, Y', X'$ 分别由概形 $X_0, Y_0, Y'_0, X'_0$ 表示时，
本引理的条件 (1) 等价于《概形的导出范畴》引理
\ref{perfect-lemma-single-complex-base-change-condition} 的条件 (1)。
用 $f_0, g_0, g'_0, f'_0$ 表示这些概形之间分别对应于
$f, g, g', f'$ 的态射。% P09-ERR-0466
可以假设 $K = \epsilon^*K_0$ 且 $K' = \epsilon^*K'_0$，
其中 $K_0 \in D_\QCoh(\mathcal{O}_{X_0})$ 且
$K'_0 \in D_\QCoh(\mathcal{O}_{X'_0})$；见引理
\ref{spaces-perfect-lemma-derived-quasi-coherent-small-etale-site}。
此外，映射 $Lg^*K \to K'$ 是映射 % P09-ERR-0467
$L(g_0)^*K_0 \to K'_0$ 的拉回；记号见评注
\ref{spaces-perfect-remark-match-total-direct-images}。
回忆 $\mathcal{O}_{X, \overline{x}}$ 是
$\mathcal{O}_{X, x}$ 的严格亨泽尔化
（《空间的性质》引理
\ref{spaces-properties-lemma-describe-etale-local-ring}），并且
$$
K_{\overline{x}} =
K_{0, x} \otimes_{\mathcal{O}_{X, x}}^\mathbf{L} \mathcal{O}_{X, \overline{x}}
\quad\text{且}\quad
K'_{\overline{x}'} =
K'_{0, x'} \otimes_{\mathcal{O}_{X', x'}}^\mathbf{L}
\mathcal{O}_{X', \overline{x}'}
$$
（类似于《空间的性质》引理
\ref{spaces-properties-lemma-stalk-quasi-coherent}）。
考虑交换图
$$
\xymatrix{
H^i(K_{\overline{x}}
\otimes_{\mathcal{O}_{Y, \overline{y}}}^\mathbf{L}
\mathcal{O}_{Y', \overline{y}'})
\otimes_{(\mathcal{O}_{X, \overline{x}}
\otimes_{\mathcal{O}_{Y, \overline{y}}}
\mathcal{O}_{Y', \overline{y}'})}
\mathcal{O}_{X', \overline{x}'}
\ar[r] &
H^i(K'_{\overline{x}'}) \\
H^i(K_{0, x} \otimes_{\mathcal{O}_{Y, y}}^\mathbf{L} \mathcal{O}_{Y', y'})
\otimes_{(\mathcal{O}_{X, x} \otimes_{\mathcal{O}_{Y, y}} \mathcal{O}_{Y', y'})}
\mathcal{O}_{X', x'}
\ar[u] \ar[r] &
H^i(K'_{0, x'}) \ar[u]
}
$$
必须证明，下方水平箭头是同构当且仅当上方水平箭头是同构。
由于 $\mathcal{O}_{X', x'} \to \mathcal{O}_{X', \overline{x}'}$
是忠实平坦的（《代数进阶》引理
\ref{more-algebra-lemma-dumb-properties-henselization}），
只须证明上方箭头是下方箭头沿此映射的基变换。
对于右侧对象，这立即由上面给出的茎之间的关系得出。
对于左侧对象，只须证明
\begin{align*}
&
H^i\left(
(K_{0, x}
\otimes_{\mathcal{O}_{X, x}}^\mathbf{L} \mathcal{O}_{X, \overline{x}})
\otimes_{\mathcal{O}_{Y, \overline{y}}}^\mathbf{L}
\mathcal{O}_{Y', \overline{y}'}\right) \\
& =
H^i(K_{0, x} \otimes_{\mathcal{O}_{Y, y}}^\mathbf{L} \mathcal{O}_{Y', y'})
\otimes_{(\mathcal{O}_{X, x} \otimes_{\mathcal{O}_{Y, y}} \mathcal{O}_{Y', y'})}
(\mathcal{O}_{X, \overline{x}}
\otimes_{\mathcal{O}_{Y, \overline{y}}}
\mathcal{O}_{Y', \overline{y}'})
\end{align*}
这由《代数进阶》引理
\ref{more-algebra-lemma-lemma-tor-independent-flat-compare} 得出。
该引理的平坦性假设由上文以及《代数》引理
\ref{algebra-lemma-strictly-henselian-functorial-improve} 成立；
后一个引理说明 $\mathcal{O}_{X, \overline{x}}$ 是
$\mathcal{O}_{X, x} \otimes_{\mathcal{O}_{Y, y}}
\mathcal{O}_{Y, \overline{y}}$ 的严格亨泽尔化，并且
$\mathcal{O}_{Y', \overline{y}'}$ 是
$\mathcal{O}_{Y', y'} \otimes_{\mathcal{O}_{Y, y}}
\mathcal{O}_{Y, \overline{y}}$ 的严格亨泽尔化。
\end{proof}

\begin{lemma}
\label{spaces-perfect-lemma-single-complex-base-change}
设 $S$ 为概形。设
$$
\xymatrix{
X' \ar[r]_{g'} \ar[d]_{f'} &
X \ar[d]^f \\
Y' \ar[r]^g &
Y
}
$$
为 $S$ 上代数空间的笛卡尔图。
设 $K \in D_\QCoh(\mathcal{O}_X)$，并设 $L(g')^*K \to K'$
是 $D_\QCoh(\mathcal{O}_{X'})$ 中的映射。如果
\begin{enumerate}
\item 引理 \ref{spaces-perfect-lemma-single-complex-base-change-condition} 的等价条件成立；并且
\item $f$ 拟紧且拟分离，
\end{enumerate}
那么复合映射 $Lg^*Rf_*K \to Rf'_*L(g')^*K \to Rf'_*K'$
是同构。
\end{lemma}

\begin{proof}
为检验该映射是同构，可以在 $Y'$ 上 \'etale 局部地工作。
因此可以假设 $g : Y' \to Y$ 是仿射概形的态射。
在这种情形中，将使用引理 \ref{spaces-perfect-lemma-induction-principle} 的归纳原理，
证明：对于拟紧且拟分离代数空间 $U$（它在 $X$ 上 \'etale），
类似构造的映射
$Lg^*R(U \to Y)_*K|_U \to R(U' \to Y')_*K'|_{U'}$
是同构。这里
$U' = X' \times_{g', X} U = Y' \times_{g, Y} U$。

\medskip\noindent
如果 $U$ 是概形（例如仿射概形），那么该结果成立。
事实上，此时 $Y, Y', U, U'$ 是概形；由引理
\ref{spaces-perfect-lemma-derived-quasi-coherent-small-etale-site}，$K$ 和 $K'$
来自底层概形导出范畴中的对象；而《概形的导出范畴》引理
\ref{perfect-lemma-single-complex-base-change-condition} 的条件
对这些复形成立，这是引理
\ref{spaces-perfect-lemma-single-complex-base-change-condition} 的结论。
因此（利用评注
\ref{spaces-perfect-remark-match-total-direct-images} 中说明的相容性），
可以应用概形情形下的结果，即《概形的导出范畴》引理
\ref{perfect-lemma-single-complex-base-change}。

\medskip\noindent
下面是归纳步骤。设 $(U \subset W, V \to W)$ 为初等特异方形，
其中 $W$ 是在 $X$ 上 \'etale 的拟紧且拟分离代数空间，
$U$ 拟紧，$V$ 仿射，并且结果对 $U$、$V$ 和
$U \times_W V$ 成立。为简化记号，在此把 $W$ 替换为 $X$
（这一步在此处是允许的）。% P09-ERR-0468
用 $a : U \to Y$、$b : V \to Y$ 和
$c : U \times_X V \to Y$ 表示显然的态射。% P09-ERR-0469
设 $a' : U' \to Y'$、$b' : V' \to Y'$ 和
$c' : U' \times_{X'} V' \to Y'$ 分别为 $a$、$b$ 和 $c$ 的基变换。
使用相对 Mayer--Vietoris 的特异三角
（引理 \ref{spaces-perfect-lemma-unbounded-relative-mayer-vietoris}），
得到交换图
$$
\xymatrix{
Lg^*Rf_*K \ar[r] \ar[d] & Rf'_*K' \ar[d] \\
Lg^*Ra_*K|_U \oplus Lg^*Rb_*K|_V \ar[r] \ar[d] &
Ra'_* K'|_{U'} \oplus Rb'_* K'|_{V'} \ar[d] \\
Lg^*Rc_*K|_{U \times_X V} \ar[r] \ar[d] &
Rc'_*K'|_{U' \times_{X'} V'} \ar[d] \\
Lg^*Rf_* K[1] \ar[r] &
Rf'_* K'[1]
}
$$
由于第二和第三个水平箭头是同构，第一个也是同构
（《导出范畴》引理 \ref{derived-lemma-third-isomorphism-triangle}），
从而本引理的证明完成。
\end{proof}

\begin{lemma}
\label{spaces-perfect-lemma-single-complex-base-change-condition-inherited}
设 $S$ 为概形。设
$$
\xymatrix{
X' \ar[r]_{g'} \ar[d]_{f'} &
X \ar[d]^f \\
S' \ar[r]^g &
S
}
$$
为 $S$ 上代数空间的笛卡尔图。
设 $K \in D_\QCoh(\mathcal{O}_X)$，并设 $L(g')^*K \to K'$
是 $D_\QCoh(\mathcal{O}_{X'})$ 中的映射。
如果引理 \ref{spaces-perfect-lemma-single-complex-base-change-condition} 的等价条件成立，
那么
\begin{enumerate}
\item 对 $E \in D_\QCoh(\mathcal{O}_X)$，
引理 \ref{spaces-perfect-lemma-single-complex-base-change-condition} 的等价条件
对映射
$L(g')^*(E \otimes^\mathbf{L} K) \to L(g')^*E \otimes^\mathbf{L} K'$
成立；
\item 如果 $E$ 属于 $D(\mathcal{O}_X)$ 且完美，% P09-ERR-0470
那么引理 \ref{spaces-perfect-lemma-single-complex-base-change-condition} 的等价条件
对映射
$L(g')^*R\SheafHom(E, K) \to R\SheafHom(L(g')^*E, K')$ 成立；并且
\item 如果 $K$ 有下界，且 $E$ 属于 $D(\mathcal{O}_X)$
并且拟凝聚，那么引理 % P09-ERR-0471
\ref{spaces-perfect-lemma-single-complex-base-change-condition} 的等价条件对映射
$L(g')^*R\SheafHom(E, K) \to R\SheafHom(L(g')^*E, K')$ 成立。
\end{enumerate}
\end{lemma}

\begin{proof}
由引理
\ref{spaces-perfect-lemma-quasi-coherence-pullback}、
\ref{spaces-perfect-lemma-quasi-coherence-tensor-product} 和
\ref{spaces-perfect-lemma-quasi-coherence-internal-hom}，以及《位点上的上同调》引理
\ref{sites-cohomology-lemma-pseudo-coherent-pullback} 和
\ref{sites-cohomology-lemma-perfect-pullback}，
所涉复形具有拟凝聚上同调层，因而该陈述有意义。
说明这一点后，在情形 (1) 中，利用张量积的传递性计算
(\ref{spaces-perfect-equation-bc}) 的源和靶，即可检验映射
(\ref{spaces-perfect-equation-bc}) 是同构；见《代数进阶》引理
\ref{more-algebra-lemma-triple-tensor-product}。
情形 (2) 和 (3) 中的映射是复合映射
$$
L(g')^*R\SheafHom(E, K) \to R\SheafHom(L(g')^*E, L(g')^*K)
\to R\SheafHom(L(g')^*E, K')
$$
其中第一个箭头是《位点上的上同调》评注
\ref{sites-cohomology-remark-prepare-fancy-base-change} 中的箭头，
第二个箭头来自给定映射 $L(g')^*K \to K'$。
为证明映射 (\ref{spaces-perfect-equation-bc}) 是同构，在情形 (2) 中，
把 $E_x$ 表示为有限生成投射 $\mathcal{O}_{X. x}$-模的有界复形；% P09-ERR-0472
在情形 (3) 中，则用有限秩自由模的有上界复形表示它；
然后计算该箭头的源和靶。
略去若干细节。
\end{proof}

\begin{lemma}
\label{spaces-perfect-lemma-base-change-tensor}
设 $S$ 为概形。设 $f : X \to Y$ 为 $S$ 上代数空间的
拟紧且拟分离态射。设 $E \in D_\QCoh(\mathcal{O}_X)$。
设 $\mathcal{G}^\bullet$ 是由拟凝聚 $\mathcal{O}_X$-模组成的
有上界复形，并且每一项在 $Y$ 上平坦。那么
$$
Rf_*(E \otimes^\mathbf{L}_{\mathcal{O}_X} \mathcal{G}^\bullet)
$$
的形成与任意基变换可交换（精确陈述见证明）。
\end{lemma}

\begin{proof}
该陈述的意义如下。设 $g : Y' \to Y$ 为代数空间的态射，
并考虑基变换图
$$
\xymatrix{
X' \ar[r]_{g'} \ar[d]_{f'} &
X \ar[d]^f \\
Y' \ar[r]^g &
Y
}
$$
换言之，$X' = Y' \times_Y X$。本引理断言
$$
Lg^*Rf_*(E \otimes^\mathbf{L}_{\mathcal{O}_X} \mathcal{G}^\bullet)
\longrightarrow
Rf'_*(L(g')^*E \otimes^\mathbf{L}_{\mathcal{O}_{X'}} (g')^*\mathcal{G}^\bullet)
$$
是同构。注意，在右端对 $\mathcal{G}^\bullet$ {\bf 不}使用导出拉回。
为证明这一点，应用引理 \ref{spaces-perfect-lemma-single-complex-base-change} 和
\ref{spaces-perfect-lemma-single-complex-base-change-condition-inherited}，可见只须证明典范映射
$$
L(g')^*\mathcal{G}^\bullet \to (g')^*\mathcal{G}^\bullet
$$
满足引理 \ref{spaces-perfect-lemma-single-complex-base-change-condition} 的等价条件。
这可通过在茎上检验条件得出；在茎上，它立即源于如下事实：
$$
\mathcal{G}^\bullet_{\overline{x}}
\otimes_{\mathcal{O}_{Y, \overline{y}}}
\mathcal{O}_{Y', \overline{y}'}
$$
由我们对复形 $\mathcal{G}^\bullet$ 的假设计算导出张量积。
\end{proof}

\begin{lemma}
\label{spaces-perfect-lemma-base-change-RHom}
设 $S$ 为概形。设 $f : X \to Y$ 为 $S$ 上代数空间的
拟紧且拟分离态射。设 $E$ 是 $D(\mathcal{O}_X)$ 的对象。
设 $\mathcal{G}^\bullet$ 为拟凝聚 $\mathcal{O}_X$-模复形。如果
\begin{enumerate}
\item $E$ 完美，$\mathcal{G}^\bullet$ 有上界，% P09-ERR-0473
并且 $\mathcal{G}^n$ 在 $Y$ 上平坦；或
\item $E$ 拟凝聚，$\mathcal{G}^\bullet$ 有界，
并且 $\mathcal{G}^n$ 在 $Y$ 上平坦，
\end{enumerate}
那么
$$
Rf_*R\SheafHom(E, \mathcal{G}^\bullet)
$$
的形成与任意基变换可交换（精确陈述见证明）。
\end{lemma}

\begin{proof}
该陈述的意义如下。设 $g : Y' \to Y$ 为代数空间的态射，
并考虑基变换图
$$
\xymatrix{
X' \ar[r]_h \ar[d]_{f'} & % P09-ERR-0474
X \ar[d]^f \\
Y' \ar[r]^g &
Y
}
$$
换言之，$X' = Y' \times_Y X$。本引理断言
$$
Lg^*Rf_*R\SheafHom(E, \mathcal{G}^\bullet)
\longrightarrow
R(f')_*R\SheafHom(L(g')^*E, (g')^*\mathcal{G}^\bullet)
$$
是同构。注意，在右端对 $\mathcal{G}^\bullet$ {\bf 不}使用导出拉回。
为证明这一点，应用引理 \ref{spaces-perfect-lemma-single-complex-base-change} 和
\ref{spaces-perfect-lemma-single-complex-base-change-condition-inherited}，可见只须证明典范映射
$$
L(g')^*\mathcal{G}^\bullet \to (g')^*\mathcal{G}^\bullet
$$
满足引理 \ref{spaces-perfect-lemma-single-complex-base-change-condition} 的等价条件。
这一点已在引理 \ref{spaces-perfect-lemma-base-change-tensor} 的证明中得到证明。
\end{proof}















\section{构造完美复形}
\label{spaces-perfect-section-producing-perfect}

\noindent
下述引理是我们构造完美复形的主要技术工具。
这个结果后面的版本将通过诺特逼近归约到它。

\begin{lemma}
\label{spaces-perfect-lemma-perfect-direct-image}
设 $S$ 为概形。设 $Y$ 为 $S$ 上的诺特代数空间。
设 $f : X \to Y$ 为局部有限型且拟分离的代数空间态射。
设 $E \in D(\mathcal{O}_X)$，满足
\begin{enumerate}
\item $E \in D^b_{\textit{Coh}}(\mathcal{O}_X)$；
\item $H^i(E)$ 的支撑在 $Y$ 上对所有 $i$ 都适当；
\item $E$ 作为 $D(f^{-1}\mathcal{O}_Y)$ 的对象具有有限 Tor 维数。
\end{enumerate}
那么 $Rf_*E$ 是 $D(\mathcal{O}_Y)$ 的完美对象。
\end{lemma}

\begin{proof}
由引理 \ref{spaces-perfect-lemma-direct-image-coherent}，$Rf_*E$ 是
$D^b_{\textit{Coh}}(\mathcal{O}_Y)$ 的对象。因此 $Rf_*E$ 拟凝聚
（引理 \ref{spaces-perfect-lemma-identify-pseudo-coherent-noetherian}）。
所以只须证明 $Rf_*E$ 具有有限 Tor 维数；见《位点上的上同调》引理
\ref{sites-cohomology-lemma-perfect}。
由引理 \ref{spaces-perfect-lemma-tor-qc-qs}，只须检验
$Rf_*(E) \otimes_{\mathcal{O}_Y}^\mathbf{L} \mathcal{F}$
对所有拟凝聚层 $\mathcal{F}$（位于 $Y$ 上）
都具有一致有界的上同调。它有上界是显然的，因为 $Rf_*(E)$ 有上界。
设 $T \subset |X|$ 为 $H^i(E)$ 的支撑之并，其中 $i$ 取遍所有整数。
由假设 (1)、(2) 和引理 \ref{spaces-perfect-lemma-union-closed-proper-over-base}，
$T$ 在 $Y$ 上适当。特别地，存在拟紧开子空间
$X' \subset X$，它包含 $T$。置 $f' = f|_{X'}$。我们有
$Rf_*(E) = Rf'_*(E|_{X'})$，因为 $E$ 在
$X \setminus T$ 上限制为零。
因此可以把 $X$ 换成 $X'$，并假设 $f$ 拟紧。
已经假设 $f$ 拟分离。因此
$$
Rf_*(E) \otimes_{\mathcal{O}_Y}^\mathbf{L} \mathcal{F} =
Rf_*\left(E \otimes_{\mathcal{O}_X}^\mathbf{L} Lf^*\mathcal{F}\right) =
Rf_*\left(E \otimes_{f^{-1}\mathcal{O}_Y}^\mathbf{L} f^{-1}\mathcal{F}\right)
$$
这由引理 \ref{spaces-perfect-lemma-cohomology-base-change} 和《位点上的上同调》引理
\ref{sites-cohomology-lemma-variant-derived-pullback} 得出。
由假设 (3)，复形
$E \otimes_{f^{-1}\mathcal{O}_Y}^\mathbf{L} f^{-1}\mathcal{F}$
的上同调层位于某个给定有限范围中，记为 $[a, b]$。
于是它的 $Rf_*$ 的上同调位于 $[a, \infty)$ 中，结论得证。
\end{proof}

\begin{lemma}
\label{spaces-perfect-lemma-tensor-perfect}
设 $S$ 为概形。设 $B$ 为 $S$ 上的诺特代数空间。
设 $f : X \to B$ 为局部有限型且拟分离的代数空间态射。
设 $E \in D(\mathcal{O}_X)$ 完美。
设 $\mathcal{G}^\bullet$ 为由拟凝聚 $\mathcal{O}_X$-模组成的有界复形，
每一项在 $B$ 上平坦，且支撑在 $B$ 上适当。那么
$K = Rf_*(E \otimes^\mathbf{L}_{\mathcal{O}_X} \mathcal{G}^\bullet)$
是 $D(\mathcal{O}_B)$ 的完美对象。
\end{lemma}

\begin{proof}
由引理 \ref{spaces-perfect-lemma-perfect-direct-image}，对象 $K$ 完美。
下面检验该引理适用：局部地，$E$ 同构于有限秩自由
$\mathcal{O}_X$-模的有限复形。因此局部地，
$E \otimes^\mathbf{L}_{\mathcal{O}_X} \mathcal{G}^\bullet$
同构于一个有限复形，其各项具有形式
$$
\bigoplus\nolimits_{i = a, \ldots, b} (\mathcal{G}^i)^{\oplus r_i}
$$
其中 $a, b, r_a, \ldots, r_b$ 是整数。这立即蕴含上同调层
$H^i(E \otimes^\mathbf{L}_{\mathcal{O}_X} \mathcal{G})$ 是凝聚的。% P09-ERR-0475
Tor 维数的假设也成立，因为 $\mathcal{G}^i$ 在
$f^{-1}\mathcal{O}_Y$ 上平坦。% P09-ERR-0476
\end{proof}

\begin{lemma}
\label{spaces-perfect-lemma-ext-perfect}
设 $S$ 为概形。设 $B$ 为 $S$ 上的诺特代数空间。
设 $f : X \to B$ 为局部有限型且拟分离的代数空间态射。
设 $E \in D(\mathcal{O}_X)$ 完美。
设 $\mathcal{G}^\bullet$ 为由拟凝聚 $\mathcal{O}_X$-模组成的有界复形，
每一项在 $B$ 上平坦，且支撑在 $B$ 上适当。那么
$K = Rf_*R\SheafHom(E, \mathcal{G})$ 是 % P09-ERR-0477
$D(\mathcal{O}_B)$ 的完美对象。
\end{lemma}

\begin{proof}
由于 $E$ 是完美复形，存在对偶完美复形 $E^\vee$；见《位点上的上同调》引理
\ref{sites-cohomology-lemma-dual-perfect-complex}。
注意
$R\SheafHom(E, \mathcal{G}^\bullet) =
E^\vee \otimes^\mathbf{L}_{\mathcal{O}_X} \mathcal{G}^\bullet$。
因此，由引理 \ref{spaces-perfect-lemma-tensor-perfect}，$K$ 是完美的。
\end{proof}






\section{Ext 的投影公式}
\label{spaces-perfect-section-ext}

\noindent
引理 \ref{spaces-perfect-lemma-compute-ext}（或文献中的类似结果）
有时可用于验证障碍理论的性质，而这些性质是验证 Quot 函子、
Hilbert 概形及其他模问题的某个 Artin 判据所需的。
假设 $f : X \to Y$ 是代数空间的适当、平坦、有限表现态射，
并且 $E \in D(\mathcal{O}_X)$ 完美。这里该引理说明
$$
\Ext^i_X(E, f^*\mathcal{F}) =
\Ext^i_Y((Rf_*E^\vee)^\vee, \mathcal{F})
$$
其中 $\mathcal{F}$ 是 $Y$ 上的拟凝聚层。
这样写使它看起来像 Ext 的投影公式；事实上，该结果很容易由引理
\ref{spaces-perfect-lemma-cohomology-base-change} 得出。

\begin{lemma}
\label{spaces-perfect-lemma-compute-tensor-perfect}
假设和记号如引理 \ref{spaces-perfect-lemma-tensor-perfect}。那么存在函子性同构
$$
H^i(B, K \otimes^\mathbf{L}_{\mathcal{O}_B} \mathcal{F})
\longrightarrow
H^i(X, E \otimes^\mathbf{L}_{\mathcal{O}_X}
(\mathcal{G}^\bullet \otimes_{\mathcal{O}_X} f^*\mathcal{F}))
$$
其中 $\mathcal{F}$ 是 $B$ 上的拟凝聚层；这些同构与边界映射相容
（见证明）。
\end{lemma}

\begin{proof}
我们有
$$
\mathcal{G}^\bullet \otimes_{\mathcal{O}_X}^\mathbf{L} Lf^*\mathcal{F} =
\mathcal{G}^\bullet \otimes_{f^{-1}\mathcal{O}_B}^\mathbf{L} f^{-1}\mathcal{F} =
\mathcal{G}^\bullet \otimes_{f^{-1}\mathcal{O}_B} f^{-1}\mathcal{F} =
\mathcal{G}^\bullet \otimes_{\mathcal{O}_X} f^*\mathcal{F}
$$
第一个等号由《位点上的上同调》引理
\ref{sites-cohomology-lemma-variant-derived-pullback} 得出，% P09-ERR-0478
第二个等号源于 $\mathcal{G}^n$ 是平坦
$f^{-1}\mathcal{O}_B$-模，第三个等号源于拉回的定义。
因此得到
\begin{align*}
H^i(X, E \otimes^\mathbf{L}_{\mathcal{O}_X}
(\mathcal{G}^\bullet \otimes_{\mathcal{O}_X} f^*\mathcal{F}))
& =
H^i(X, E \otimes^\mathbf{L}_{\mathcal{O}_X} \mathcal{G}^\bullet
\otimes_{\mathcal{O}_X}^\mathbf{L} Lf^*\mathcal{F}) \\
& =
H^i(B,
Rf_*(E \otimes^\mathbf{L}_{\mathcal{O}_X} \mathcal{G}^\bullet
\otimes^\mathbf{L}_{\mathcal{O}_X} Lf^*\mathcal{F})) \\
& =
H^i(B,
Rf_*(E \otimes^\mathbf{L}_{\mathcal{O}_X} \mathcal{G}^\bullet)
\otimes^\mathbf{L}_{\mathcal{O}_B} \mathcal{F}) \\
& =
H^i(B, K \otimes^\mathbf{L}_{\mathcal{O}_B} \mathcal{F})
\end{align*}
第一个等号由上文得出，第二个等号由 Leray 得出
（《位点上的上同调》评注 \ref{sites-cohomology-remark-before-Leray}），
第三个等号由引理 \ref{spaces-perfect-lemma-cohomology-base-change} 得出。
关于边界映射的陈述含义如下：给定短正合列
$0 \to \mathcal{F}_1 \to \mathcal{F}_2 \to \mathcal{F}_3 \to 0$，
这些同构嵌入交换图
$$
\xymatrix{
H^i(B, K \otimes^\mathbf{L}_{\mathcal{O}_B} \mathcal{F}_3)
\ar[r] \ar[d]_\delta &
H^i(X, E \otimes^\mathbf{L}_{\mathcal{O}_X}
(\mathcal{G}^\bullet \otimes_{\mathcal{O}_X} f^*\mathcal{F}_3)) \ar[d]^\delta \\
H^{i + 1}(B, K \otimes^\mathbf{L}_{\mathcal{O}_B} \mathcal{F}_1)
\ar[r] &
H^{i + 1}(X, E \otimes^\mathbf{L}_{\mathcal{O}_X}
(\mathcal{G}^\bullet \otimes_{\mathcal{O}_X} f^*\mathcal{F}_1))
}
$$
其中边界映射来自特异三角
$$
K \otimes^\mathbf{L}_{\mathcal{O}_B} \mathcal{F}_1 \to
K \otimes^\mathbf{L}_{\mathcal{O}_B} \mathcal{F}_2 \to
K \otimes^\mathbf{L}_{\mathcal{O}_B} \mathcal{F}_3 \to
K \otimes^\mathbf{L}_{\mathcal{O}_B} \mathcal{F}_1[1]
$$
以及 $D(\mathcal{O}_X)$ 中与复形的短正合列
$$
0 \to
\mathcal{G}^\bullet \otimes_{\mathcal{O}_X} f^*\mathcal{F}_1 \to
\mathcal{G}^\bullet \otimes_{\mathcal{O}_X} f^*\mathcal{F}_2 \to
\mathcal{G}^\bullet \otimes_{\mathcal{O}_X} f^*\mathcal{F}_3 \to 0
$$
相伴的特异三角。这个序列正合，因为 $\mathcal{G}^n$ 在 $B$ 上平坦。
略去对所显示图之交换性的验证。
\end{proof}

\begin{lemma}
\label{spaces-perfect-lemma-compute-ext-perfect}
假设和记号如引理 \ref{spaces-perfect-lemma-ext-perfect}。% P09-ERR-0479
那么存在函子性同构
$$
H^i(B, K \otimes^\mathbf{L}_{\mathcal{O}_B} \mathcal{F})
\longrightarrow
\Ext^i_{\mathcal{O}_X}(E,
\mathcal{G}^\bullet \otimes_{\mathcal{O}_X} f^*\mathcal{F})
$$
其中 $\mathcal{F}$ 是 $B$ 上的拟凝聚层；这些同构与边界映射相容
（见证明）。
\end{lemma}

\begin{proof}
如引理 \ref{spaces-perfect-lemma-ext-perfect} 的证明中那样，令
$E^\vee$ 为对偶完美复形，并回忆
$K = Rf_*(E^\vee \otimes_{\mathcal{O}_X}^\mathbf{L} \mathcal{G}^\bullet)$。
此外，我们还有
$$
\Ext^i_{\mathcal{O}_X}(E,
\mathcal{G}^\bullet \otimes_{\mathcal{O}_X} f^*\mathcal{F})
=
H^i(X, E^\vee \otimes^\mathbf{L}_{\mathcal{O}_X}
(\mathcal{G}^\bullet \otimes_{\mathcal{O}_X} f^*\mathcal{F}))
$$
这是 $E^\vee$ 的构造所给出的。所以这些同构的存在性由引理
\ref{spaces-perfect-lemma-compute-tensor-perfect} 应用于 $E^\vee$ 和
$\mathcal{G}^\bullet$ 得出。
关于边界映射的陈述含义如下：给定短正合列
$0 \to \mathcal{F}_1 \to \mathcal{F}_2 \to \mathcal{F}_3 \to 0$，
这些同构组成交换图
$$
\xymatrix{
H^i(B, K \otimes^\mathbf{L}_{\mathcal{O}_B} \mathcal{F}_3)
\ar[r] \ar[d]_\delta &
\Ext^i_{\mathcal{O}_X}(E,
\mathcal{G}^\bullet \otimes_{\mathcal{O}_X} f^*\mathcal{F}_3) \ar[d]^\delta \\
H^{i + 1}(B, K \otimes^\mathbf{L}_{\mathcal{O}_B} \mathcal{F}_1)
\ar[r] &
\Ext^{i + 1}_{\mathcal{O}_X}(E,
\mathcal{G}^\bullet \otimes_{\mathcal{O}_X} f^*\mathcal{F}_1)
}
$$
其中边界映射来自特异三角
$$
K \otimes^\mathbf{L}_{\mathcal{O}_B} \mathcal{F}_1 \to
K \otimes^\mathbf{L}_{\mathcal{O}_B} \mathcal{F}_2 \to
K \otimes^\mathbf{L}_{\mathcal{O}_B} \mathcal{F}_3 \to
K \otimes^\mathbf{L}_{\mathcal{O}_B} \mathcal{F}_1[1]
$$
以及 $D(\mathcal{O}_X)$ 中与复形的短正合列
$$
0 \to
\mathcal{G}^\bullet \otimes_{\mathcal{O}_X} f^*\mathcal{F}_1 \to
\mathcal{G}^\bullet \otimes_{\mathcal{O}_X} f^*\mathcal{F}_2 \to
\mathcal{G}^\bullet \otimes_{\mathcal{O}_X} f^*\mathcal{F}_3 \to 0
$$
相伴的特异三角。这个序列正合，因为 $\mathcal{G}^n$ 在 $B$ 上平坦。
略去对所显示图之交换性的验证。
\end{proof}

\begin{lemma}
\label{spaces-perfect-lemma-compute-ext}
设 $S$ 为概形。设 $f : X \to B$ 为 $S$ 上代数空间的态射，
$E \in D(\mathcal{O}_X)$，并设 $\mathcal{F}^\bullet$ 为 % P09-ERR-0480
$\mathcal{O}_X$-模复形。假设
\begin{enumerate}
\item $B$ 是诺特的；
\item $f$ 局部有限型且拟分离；
\item $E \in D^-_{\textit{Coh}}(\mathcal{O}_X)$；
\item $\mathcal{G}^\bullet$ 是由凝聚 $\mathcal{O}_X$-模组成的
有界复形，每一项在 $B$ 上平坦，且支撑在 $B$ 上适当。% P09-ERR-0481
\end{enumerate}
那么下列两个陈述成立：
\begin{enumerate}
\item[(A)] 对每个 $m \in \mathbf{Z}$，存在完美对象 $K$
（属于 $D(\mathcal{O}_B)$）和函子性映射
$$
\alpha^i_\mathcal{F} :
\Ext^i_{\mathcal{O}_X}(E,
\mathcal{G}^\bullet \otimes_{\mathcal{O}_X} f^*\mathcal{F})
\longrightarrow
H^i(B, K \otimes^\mathbf{L}_{\mathcal{O}_B} \mathcal{F})
$$
其中 $\mathcal{F}$ 是 $B$ 上的拟凝聚层；这些映射与边界映射相容
（见证明），并且 $\alpha^i_\mathcal{F}$ 在 $i \leq m$ 时是同构；
\item[(B)] 存在拟凝聚对象 $L \in D(\mathcal{O}_B)$ 和函子性同构
$$
\Ext^i_{\mathcal{O}_B}(L, \mathcal{F}) \longrightarrow
\Ext^i_{\mathcal{O}_X}(E,
\mathcal{G}^\bullet \otimes_{\mathcal{O}_X} f^*\mathcal{F})
$$
其中 $\mathcal{F}$ 是 $B$ 上的拟凝聚层；这些同构与边界映射相容。
\end{enumerate}
\end{lemma}

\begin{proof}
(A) 的证明。假设 $\mathcal{G}^i$ 仅当 $i \in [a, b]$ 时非零。
可以把 $X$ 换成 $\mathcal{G}^i$ 的支撑之并的拟紧开邻域。
因此可以假设 $X$ 是诺特的。在这种情形中，$X$ 和 $f$
都是拟紧且拟分离的。选取逼近 $P \to E$，其中 $P$ 是
$(X, E, -m - 1 + a)$ 的完美复形逼近
（这由定理 \ref{spaces-perfect-theorem-approximation} 保证存在）。
于是诱导映射
$$
\Ext^i_{\mathcal{O}_X}(E,
\mathcal{G}^\bullet \otimes_{\mathcal{O}_X} f^*\mathcal{F})
\longrightarrow
\Ext^i_{\mathcal{O}_X}(P,
\mathcal{G}^\bullet \otimes_{\mathcal{O}_X} f^*\mathcal{F})
$$
当 $i \leq m$ 时是同构。事实上，该映射的核和余核分别是
下列两个群的商模和子模：
$$
\Ext^i_{\mathcal{O}_X}(C,
\mathcal{G}^\bullet \otimes_{\mathcal{O}_X} f^*\mathcal{F})
\quad\text{分别}\quad
\Ext^{i + 1}_{\mathcal{O}_X}(C,
\mathcal{G}^\bullet \otimes_{\mathcal{O}_X} f^*\mathcal{F})
$$
其中 $C$ 是 $P \to E$ 的锥。由于 $C$ 在次数
$\geq -m - 1 + a$ 的上同调层消失，由《导出范畴》引理
\ref{derived-lemma-negative-exts}，这些 $\Ext$ 群对
$i \leq m + 1$ 为零。
这把问题归约到 $E$ 为完美复形的情形，也就是引理
\ref{spaces-perfect-lemma-compute-ext-perfect}。% P09-ERR-0482
关于边界的陈述已在引理
\ref{spaces-perfect-lemma-compute-ext-perfect} 的证明中说明。

\medskip\noindent
(B) 的证明。如 (A) 的证明中那样，可以假设 $X$ 是诺特的。
由引理 \ref{spaces-perfect-lemma-identify-pseudo-coherent-noetherian}，$E$ 拟凝聚。
由引理 \ref{spaces-perfect-lemma-pseudo-coherent-hocolim}，可以写成
$E = \text{hocolim} E_n$，其中 $E_n$ 完美，并且 $E_n \to E$
在截断 $\tau_{\geq -n}$ 上诱导同构。设 $E_n^\vee$ 为对偶完美复形
（《位点上的上同调》引理
\ref{sites-cohomology-lemma-dual-perfect-complex}）。
得到完美对象的逆系统
$\ldots \to E_3^\vee \to E_2^\vee \to E_1^\vee$。
这又给出逆系统
$$
\ldots \to K_3 \to K_2 \to K_1\quad\text{其中}\quad
K_n = Rf_*(E_n^\vee \otimes_{\mathcal{O}_X}^\mathbf{L} \mathcal{G}^\bullet)
$$
它在 $Y$ 上完美；见引理 \ref{spaces-perfect-lemma-tensor-perfect}。% P09-ERR-0483
由引理 \ref{spaces-perfect-lemma-compute-ext-perfect} 及其证明，以及上一段的论证
（取 $P = E_n$），对任意拟凝聚层 $\mathcal{F}$（位于 $Y$ 上），
存在函子性典范映射
$$
\xymatrix{
& \Ext^i_{\mathcal{O}_X}(E,
\mathcal{G}^\bullet \otimes_{\mathcal{O}_X} f^*\mathcal{F})
\ar[ld] \ar[rd] \\
H^i(Y, K_{n + 1} \otimes_{\mathcal{O}_Y}^\mathbf{L} \mathcal{F})
\ar[rr] & &
H^i(Y, K_n \otimes_{\mathcal{O}_Y}^\mathbf{L} \mathcal{F})
}
$$
它们对 $i \leq n + a$ 是同构。
设 $L_n = K_n^\vee$ 为对偶完美复形。
于是 $L_1 \to L_2 \to L_3 \to \ldots$ 是
$D(\mathcal{O}_Y)$ 中的完美对象系统，使得对于任意拟凝聚层
$\mathcal{F}$（位于 $Y$ 上），映射
$$
\Ext^i_{\mathcal{O}_Y}(L_{n + 1}, \mathcal{F})
\longrightarrow
\Ext^i_{\mathcal{O}_Y}(L_n, \mathcal{F})
$$
对 $i \leq n + a - 1$ 是同构。这蕴含 $L_n \to L_{n + 1}$
在截断 $\tau_{\geq -n - a + 2}$ 上诱导同构
（提示：取 $L_n \to L_{n + 1}$ 的锥，并考察它最后一个非零上同调层）。
因此 $L = \text{hocolim} L_n$ 拟凝聚；见引理
\ref{spaces-perfect-lemma-pseudo-coherent-hocolim}。同伦余极限的映射性质给出
$\Ext^i_{\mathcal{O}_Y}(L, \mathcal{F}) =
\Ext^i_{\mathcal{O}_Y}(L_n, \mathcal{F})$
对 $i \leq n + a - 3$ 成立，从而证明完成。
\end{proof}

\begin{remark}
\label{spaces-perfect-remark-base-change-of-L}
拟凝聚复形 $L$（即引理 \ref{spaces-perfect-lemma-compute-ext} 的 (B) 部分中的复形）
典范地由这一情形确定。例如，$L$ 按 (B) 中方式的构造与基变换相容。
换言之，给定概形的笛卡儿图 % P09-ERR-0484
$$
\xymatrix{
X' \ar[r]_{g'} \ar[d]_{f'} &
X \ar[d]^f \\
Y' \ar[r]^g &
Y
}
$$
我们有函子性的典范同构
$$
\Ext^i_{\mathcal{O}_{Y'}}(Lg^*L, \mathcal{F}') \longrightarrow
\Ext^i_{\mathcal{O}_X}(L(g')^*E,
(g')^*\mathcal{G}^\bullet \otimes_{\mathcal{O}_{X'}} (f')^*\mathcal{F}')
$$
% P09-ERR-0485
其中 $\mathcal{F}'$ 是 $Y'$ 上的拟凝聚层。注意，在右端我们\textbf{没有}对 % P09-ERR-0486
$\mathcal{G}^\bullet$ 使用导出拉回。若将来需要这种结果，我们会在此
表述一个精确结果并给出详细证明。
\end{remark}

\section{极限与导出范畴}
\label{spaces-perfect-section-limits}

\noindent
本节汇集关于如下代数空间之导出范畴的一些结果：该代数空间是
一个代数空间逆系统的极限。更确切地说，我们将在下列情形中工作。

\begin{situation}
\label{spaces-perfect-situation-descent}
设 $S$ 为概形。设 $X = \lim_{i \in I} X_i$ 是 $S$ 上代数空间的
有向逆系统之极限，其转移态射
$f_{i'i} : X_{i'} \to X_i$ 是仿射的。把投影记为 % P09-ERR-0487
$f_i : X \to X_i$。假设 $X_i$ 对所有 $i \in I$ 都拟紧且拟分离。
还选定一个元素 $0 \in I$。
\end{situation}

\begin{lemma}
\label{spaces-perfect-lemma-descend-homomorphisms}
在情形 \ref{spaces-perfect-situation-descent} 中，令 $E_0$ 和 $K_0$ 为
$D(\mathcal{O}_{X_0})$ 的对象。置
$E_i = Lf_{i0}^*E_0$ 和 $K_i = Lf_{i0}^*K_0$（$i \geq 0$），
并置 $E = Lf_0^*E_0$ 和 $K = Lf_0^*K_0$。那么映射
$$
\colim_{i \geq 0} \Hom_{D(\mathcal{O}_{X_i})}(E_i, K_i)
\longrightarrow
\Hom_{D(\mathcal{O}_X)}(E, K)
$$
在下列任一情形中为同构：
\begin{enumerate}
\item $E_0$ 完美且 $K_0 \in D_\QCoh(\mathcal{O}_{X_0})$；或
\item $E_0$ 拟凝聚，且
$K_0 \in D_\QCoh(\mathcal{O}_{X_0})$ 具有有限 Tor 维数。
\end{enumerate}
\end{lemma}

\begin{proof}
对每个拟紧拟分离对象 $U_0$（属于 $(X_0)_{spaces, \etale}$），
考虑如下条件 $P$：典范映射
$$
\colim_{i \geq 0} \Hom_{D(\mathcal{O}_{U_i})}(E_i|_{U_i}, K_i|_{U_i})
\longrightarrow
\Hom_{D(\mathcal{O}_U)}(E|_U, K|_U)
$$
是同构，其中 $U = X \times_{X_0} U_0$ 且
$U_i = X_i \times_{X_0} U_0$。我们将用引理
\ref{spaces-perfect-lemma-induction-principle} 的归纳原理证明 $P$ 对每个 $U_0$ 成立。
该引理的条件 (2) 立即由导出范畴中 Hom 的 Mayer--Vietoris 性质得出；
见引理 \ref{spaces-perfect-lemma-mayer-vietoris-hom}。
因此只需在 $X_0$ 仿射时证明本引理。

\medskip\noindent
若 $X_0$ 仿射，则结果由概形情形得出；见《概形的导出范畴》引理
\ref{perfect-lemma-descend-homomorphisms}。
为此，使用引理 \ref{spaces-perfect-lemma-derived-quasi-coherent-small-etale-site}
的等价，并使用引理
\ref{spaces-perfect-lemma-descend-pseudo-coherent}、
\ref{spaces-perfect-lemma-descend-tor-amplitude} 和
\ref{spaces-perfect-lemma-descend-perfect} 中说明的性质对应。
\end{proof}

\begin{lemma}
\label{spaces-perfect-lemma-perfect-on-limit}
在情形 \ref{spaces-perfect-situation-descent} 中，$D(\mathcal{O}_X)$ 的完美对象范畴
是 $D(\mathcal{O}_{X_i})$ 的完美对象范畴之余极限。
\end{lemma}

\begin{proof}
对每个拟紧拟分离对象 $U_0$（属于 $(X_0)_{spaces, \etale}$），
考虑条件 $P$：函子
$$
\colim_{i \geq 0} D_{perf}(\mathcal{O}_{U_i})
\longrightarrow
D_{perf}(\mathcal{O}_U)
$$
是等价；其中 ${}_{perf}$ 表示由完美对象构成的充满子范畴，而
$U = X \times_{X_0} U_0$ 且
$U_i = X_i \times_{X_0} U_0$。我们将用引理
\ref{spaces-perfect-lemma-induction-principle} 的归纳原理证明 $P$ 对每个 $U_0$ 成立。
首先，由引理 \ref{spaces-perfect-lemma-descend-homomorphisms}，已经知道该函子全忠实。
因此只需证明本质满性。

\medskip\noindent
先验证归纳原理的条件 (2)。因而，设有初等特异方形
$(U_0 \subset X_0, V_0 \to X_0)$，并且 $P$ 对
$U_0$、$V_0$ 和 $U_0 \times_{X_0} V_0$ 成立。设 $E$ 是
$D(\mathcal{O}_X)$ 的完美对象。可以找到 $i \geq 0$、完美对象
$E_{U, i}$（位于 $U_i$ 上）和完美对象 $E_{V, i}$（位于 $V_i$ 上），
它们到 $U$ 和 $V$ 的拉回分别同构于
$E|_U$ 和 $E|_V$。记
$$
a : E_{U, i} \to (R(X \to X_i)_*E)|_{U_i}
\quad\text{且}\quad
b : E_{V, i} \to (R(X \to X_i)_*E)|_{V_i}
$$
为与同构 $L(U \to U_i)^*E_{U, i} \to E|_U$
和 $L(V \to V_i)^*E_{V, i} \to E|_V$ 伴随的映射。
由全忠实性，% P09-ERR-0488
在增大 $i$ 后，可以找到同构
$c : E_{U, i}|_{U_i \times_{X_i} V_i} \to E_{V, i}|_{U_i \times_{X_i} V_i}$，
其拉回给出下列识别：
$$
L(U \to U_i)^*E_{U, i}|_{U \times_X V} \to E|_{U \times_X V} \to
L(V \to V_i)^*E_{V, i}|_{U \times_X V}.
$$
应用引理 \ref{spaces-perfect-lemma-glue}，得到对象 $E_i$（位于 $X_i$ 上）和映射
$d : E_i \to R(X \to X_i)_*E$；其限制是映射 $a$ 和 $b$，
分别位于 $U_i$ 和 $V_i$ 上。于是显然 $E_i$ 完美，并且
$d$ 与同构 $L(X \to X_i)^*E_i \to E$ 伴随。

\medskip\noindent
最后验证归纳原理的条件 (1)，换言之，验证本引理在 $X_0$ 仿射时成立。
这由概形情形得出；见《概形的导出范畴》引理
\ref{perfect-lemma-descend-perfect}。
为此，使用引理
\ref{spaces-perfect-lemma-derived-quasi-coherent-small-etale-site} 的等价，
并使用引理 \ref{spaces-perfect-lemma-descend-perfect} 中说明的性质对应。
\end{proof}

\section{上同调与基变换，VI}
\label{spaces-perfect-section-cohomology-and-base-change-final}

\noindent
这是关于上同调与基变换的最后一节，继续第
\ref{spaces-perfect-section-cohomology-and-base-change-perfect}、
\ref{spaces-perfect-section-cohomology-base-change} 和
\ref{spaces-perfect-section-producing-perfect} 节的讨论。
备注 \ref{spaces-perfect-remark-explain-perfect-direct-image} 给出了一个易于把握的特例。

\begin{lemma}
\label{spaces-perfect-lemma-base-change-tensor-perfect}
设 $S$ 为概形。设 $f : X \to Y$ 是 $S$ 上代数空间之间的
有限表现态射。设 $E \in D(\mathcal{O}_X)$ 为完美对象。
设 $\mathcal{G}^\bullet$ 是由有限表现 $\mathcal{O}_X$-模组成的
有界复形，它在 $Y$ 上平坦，且支撑在 $Y$ 上适当。那么
$$
K = Rf_*(E \otimes_{\mathcal{O}_X}^\mathbf{L} \mathcal{G}^\bullet)
$$
是 $D(\mathcal{O}_Y)$ 的完美对象，并且它的构造与任意基变换可交换。
\end{lemma}

\begin{proof}
关于基变换的陈述是引理 \ref{spaces-perfect-lemma-base-change-tensor}。
因此只需证明 $K$ 是完美对象。若 $Y$ 诺特，则这由引理
\ref{spaces-perfect-lemma-tensor-perfect} 得出。
我们将用诺特逼近把问题归约到这种情形。
建议读者跳过本证明的其余部分。

\medskip\noindent
问题在 $Y$ 上是局部的，故可假设 $Y$ 仿射。
设 $Y = \Spec(R)$。把 $R = \colim R_i$ 写成诺特环
$R_i$ 的滤过余极限。由《空间的极限》引理
\ref{spaces-limits-lemma-descend-finite-presentation}，
存在 $i$ 和代数空间 $X_i$，它在 $R_i$ 上有限表现，且到
$R$ 的基变换是 $X$。由《空间的极限》引理
\ref{spaces-limits-lemma-descend-modules-finite-presentation}，
增大 $i$ 后，可以假设存在由有限表现
$\mathcal{O}_{X_i}$-模组成的有界复形 $\mathcal{G}_i^\bullet$，
它到 $X$ 的拉回是 $\mathcal{G}^\bullet$。增大 $i$ 后，
可以假设 $\mathcal{G}_i^n$ 在 $R_i$ 上平坦；见《空间的极限》引理
\ref{spaces-limits-lemma-descend-flat}。
增大 $i$ 后，可以假设 $\mathcal{G}_i^n$ 的支撑在 $R_i$ 上适当；
见《空间的极限》引理
\ref{spaces-limits-lemma-eventually-proper-support}。
最后，由引理 \ref{spaces-perfect-lemma-perfect-on-limit}，增大 $i$ 后，
可以假设存在完美对象 $E_i$（属于 $D(\mathcal{O}_{X_i})$），
它到 $X$ 的拉回是 $E$。由引理 \ref{spaces-perfect-lemma-tensor-perfect}，
$$
K_i =
Rf_{i, *}(E_i \otimes_{\mathcal{O}_{X_i}}^\mathbf{L} \mathcal{G}_i^\bullet)
$$
在 $\Spec(R_i)$ 上完美，其中 $f_i : X_i \to \Spec(R_i)$ 是结构态射。
由基变换结果（引理 \ref{spaces-perfect-lemma-base-change-tensor}），$K_i$ 到
$Y = \Spec(R)$ 的拉回是 $K$，故结论成立。
\end{proof}

\begin{remark}
\label{spaces-perfect-remark-explain-perfect-direct-image}
设 $R$ 为环。设 $X$ 为 $R$ 上有限表现的代数空间。
设 $\mathcal{G}$ 为有限表现 $\mathcal{O}_X$-模，它在 $R$ 上平坦，
且支撑在 $R$ 上适当。由引理
\ref{spaces-perfect-lemma-base-change-tensor-perfect}，存在由有限投射
$R$-模组成的有限复形 $M^\bullet$，使得
$$
R\Gamma(X_{R'}, \mathcal{G}_{R'}) = M^\bullet \otimes_R R'
$$
对 $R$-代数 $R'$ 函子性成立。
\end{remark}

\begin{lemma}
\label{spaces-perfect-lemma-base-change-tensor-pseudo-coherent}
设 $S$ 为概形。
设 $f : X \to Y$ 是 $S$ 上代数空间之间的有限表现态射。
设 $E \in D(\mathcal{O}_X)$ 为拟凝聚对象。
设 $\mathcal{G}^\bullet$ 是由有限表现 $\mathcal{O}_X$-模组成的
上有界复形，它在 $Y$ 上平坦，且支撑在 $Y$ 上适当。那么
$$
K = Rf_*(E \otimes_{\mathcal{O}_X}^\mathbf{L} \mathcal{G}^\bullet)
$$
是 $D(\mathcal{O}_Y)$ 的拟凝聚对象，并且它的构造与任意基变换可交换。
\end{lemma}

\begin{proof}
关于基变换的陈述是引理 \ref{spaces-perfect-lemma-base-change-tensor}。
因此只需证明 $K$ 是拟凝聚对象。
这将由引理 \ref{spaces-perfect-lemma-base-change-tensor-perfect}
通过完美复形逼近得出。建议读者跳过本证明的其余部分。

\medskip\noindent
问题在 $Y$ 上是平展局部的，故可假设 $Y$ 仿射。
于是 $X$ 拟紧且拟分离。此外，存在整数 $N$，使得全正像
$Rf_* : D_\QCoh(\mathcal{O}_X) \to D_\QCoh(\mathcal{O}_Y)$
具有上同调维数 $N$；见引理
\ref{spaces-perfect-lemma-quasi-coherence-direct-image} 中的说明。
选取整数 $b$，使得 $\mathcal{G}^i = 0$ 对 $i > b$ 成立。
只需证明 $K$ 对每个 $m$ 都是 $m$-拟凝聚的。
选取逼近 $P \to E$，其中 $P$ 是完美复形，且该逼近属于
$(X, E, m - N - 1 - b)$。这由定理
\ref{spaces-perfect-theorem-approximation} 保证存在。
选取特异三角
$$
P \to E \to C \to P[1]
$$
，它位于 $D_\QCoh(\mathcal{O}_X)$ 中。$C$ 的上同调层在次数
$\geq m - N - 1 - b$ 为零。
因此 $C \otimes^\mathbf{L} \mathcal{G}^\bullet$ 的上同调层
在次数 $\geq m - N - 1$ 为零。
从而 $Rf_*(C \otimes^\mathbf{L} \mathcal{G})$ 的上同调层 % P09-ERR-0489
在次数 $\geq m - 1$ 为零。因此
$$
Rf_*(P \otimes^\mathbf{L} \mathcal{G}) \to
Rf_*(E \otimes^\mathbf{L} \mathcal{G})
$$
在次数 $\geq m$ 的上同调层上是同构。
其次，假设 $H^i(P) = 0$ 对 $i > a$ 成立。那么
$
P \otimes^\mathbf{L} \sigma_{\geq m - N - 1 - a}\mathcal{G}^\bullet
\longrightarrow
P \otimes^\mathbf{L} \mathcal{G}^\bullet
$
在次数 $\geq m - N - 1$ 的上同调层上是同构。
因而再次得到
$$
Rf_*(P \otimes^\mathbf{L} \sigma_{\geq m - N - 1 - a}\mathcal{G}^\bullet) \to
Rf_*(P \otimes^\mathbf{L} \mathcal{G}^\bullet)
$$
在次数 $\geq m$ 的上同调层上是同构。
由引理 \ref{spaces-perfect-lemma-base-change-tensor-perfect}，源是完美复形。
因此 $K$ 如所需地为 $m$-拟凝聚。
\end{proof}

\begin{lemma}
\label{spaces-perfect-lemma-flat-proper-perfect-direct-image-general}
设 $S$ 为概形。设 $f : X \to Y$ 是 $S$ 上代数空间之间的
适当有限表现态射。
\begin{enumerate}
\item 设 $E \in D(\mathcal{O}_X)$ 完美且 $f$ 平坦。那么
$Rf_*E$ 是 $D(\mathcal{O}_Y)$ 的完美对象，并且它的构造
与任意基变换可交换。
\item 设 $\mathcal{G}$ 为有限表现 $\mathcal{O}_X$-模，
且在 $S$ 上平坦。那么 $Rf_*\mathcal{G}$ 是 % P09-ERR-0490
$D(\mathcal{O}_Y)$ 的完美对象，并且它的构造与任意基变换可交换。
\end{enumerate}
\end{lemma}

\begin{proof}
这是引理 \ref{spaces-perfect-lemma-base-change-tensor-perfect} 的特例，分别取
(1) $\mathcal{G}^\bullet$ 等于 $\mathcal{O}_X$（置于次数 $0$），以及
(2) $E = \mathcal{O}_X$，并令 $\mathcal{G}^\bullet$
由 $\mathcal{G}$ 组成（置于次数 $0$）。
\end{proof}

\begin{lemma}
\label{spaces-perfect-lemma-flat-proper-pseudo-coherent-direct-image-general}
设 $S$ 为概形。设 $f : X \to Y$ 是 $S$ 上代数空间之间的
平坦适当有限表现态射。设 $E \in D(\mathcal{O}_X)$ 拟凝聚。
那么 $Rf_*E$ 是 $D(\mathcal{O}_Y)$ 的拟凝聚对象，
并且它的构造与任意基变换可交换。
\end{lemma}

\noindent
更一般地，若 $f : X \to Y$ 适当，而 $E$（位于 $X$ 上）相对于
$Y$ 拟凝聚（《空间态射进阶》定义
\ref{spaces-more-morphisms-definition-relative-pseudo-coherence}），
则 $Rf_*E$ 拟凝聚
（但在这种一般性下，其构造不与基变换可交换）。
概形的这种情形在 \cite{Kiehl} 中得到证明。

\begin{proof}
这是引理 \ref{spaces-perfect-lemma-base-change-tensor-pseudo-coherent} 的特例，
取 $\mathcal{G} = \mathcal{O}_X$。
\end{proof}

\begin{lemma}
\label{spaces-perfect-lemma-pullback-and-limits}
设 $R$ 为环。设 $X$ 为代数空间，并设
$f : X \to \Spec(R)$ 适当、平坦且有限表现。
设 $(M_n)$ 为具有满射转移映射的 $R$-模逆系统。
那么典范映射
$$
\mathcal{O}_X \otimes_R (\lim M_n)
\longrightarrow
\lim \mathcal{O}_X \otimes_R M_n
$$
诱导从源到对目标应用 $DQ_X$ 所得对象的同构。
\end{lemma}

\begin{proof}
该陈述意指：对任意对象 $E$（属于
$D_\QCoh(\mathcal{O}_X)$），诱导映射
$$
\Hom(E, \mathcal{O}_X \otimes_R (\lim M_n))
\longrightarrow
\Hom(E, \lim \mathcal{O}_X \otimes_R M_n)
$$
是同构。由于 $D_\QCoh(\mathcal{O}_X)$ 有完美生成元
（定理 \ref{spaces-perfect-theorem-bondal-van-den-Bergh}），只需对完美 $E$ 验证。
由引理 \ref{spaces-perfect-lemma-Rlim-quasi-coherent}，
$\lim \mathcal{O}_X \otimes_R M_n = R\lim \mathcal{O}_X \otimes_R M_n$。
《位点上的上同调》第 \ref{sites-cohomology-section-global-RHom} 节的
正合函子
$R\Hom_X(E, -) : D_\QCoh(\mathcal{O}_X) \to D(R)$
与积可交换，因而与导出极限可交换，所以
$$
R\Hom_X(E, \lim \mathcal{O}_X \otimes_R M_n) =
R\lim R\Hom_X(E, \mathcal{O}_X \otimes_R M_n)
$$
令 $E^\vee$ 为对偶完美复形；见《位点上的上同调》引理
\ref{sites-cohomology-lemma-dual-perfect-complex}。我们有
$$
R\Hom_X(E, \mathcal{O}_X \otimes_R M_n) =
R\Gamma(X, E^\vee \otimes_{\mathcal{O}_X}^\mathbf{L} Lf^*M_n) =
R\Gamma(X, E^\vee) \otimes_R^\mathbf{L} M_n
$$
，这由引理 \ref{spaces-perfect-lemma-cohomology-base-change} 得出。
由引理 \ref{spaces-perfect-lemma-flat-proper-perfect-direct-image-general} 可知
$R\Gamma(X, E^\vee)$ 是 $R$-模的完美复形。
特别地，它是拟凝聚复形；由《代数进阶》引理
\ref{more-algebra-lemma-pseudo-coherent-tensor-limit}，得到
$$
R\lim R\Gamma(X, E^\vee) \otimes_R^\mathbf{L} M_n =
R\Gamma(X, E^\vee) \otimes_R^\mathbf{L} \lim M_n
$$
，正如所需。
\end{proof}

\begin{lemma}
\label{spaces-perfect-lemma-perfect-enough}
设 $A$ 为环。设 $X$ 为 $A$ 上拟紧拟分离的代数空间。
设 $K \in D^-_\QCoh(\mathcal{O}_X)$。若
$R\Gamma(X, E \otimes^\mathbf{L} K)$ 在 $D(A)$ 中拟凝聚，
对每个完美 $E$（属于 $D(\mathcal{O}_X)$）都如此，那么
$R\Gamma(X, E \otimes^\mathbf{L} K)$ 在 $D(A)$ 中拟凝聚，
对每个拟凝聚 $E$（属于 $D(\mathcal{O}_X)$）都如此。
\end{lemma}

\begin{proof}
存在整数 $N$，使得
$R\Gamma(X, -) : D_\QCoh(\mathcal{O}_X) \to D(A)$
具有上同调维数 $N$；见引理
\ref{spaces-perfect-lemma-quasi-coherence-direct-image} 中的说明。
取 $b \in \mathbf{Z}$，使得 $H^i(K) = 0$ 对 $i > b$ 成立。
设 $E$ 在 $X$ 上拟凝聚。
只需证明 $R\Gamma(X, E \otimes^\mathbf{L} K)$
对每个 $m$ 都是 $m$-拟凝聚的。
选取逼近 $P \to E$，其中 $P$ 是完美复形，且该逼近属于
$(X, E, m - N - 1 - b)$。这由定理
\ref{spaces-perfect-theorem-approximation} 保证存在。
选取特异三角
$$
P \to E \to C \to P[1]
$$
，它位于 $D_\QCoh(\mathcal{O}_X)$ 中。$C$ 的上同调层
在次数 $\geq m - N - 1 - b$ 为零。因此
$C \otimes^\mathbf{L} K$ 的上同调层在次数
$\geq m - N - 1$ 为零。
从而 $R\Gamma(X, C \otimes^\mathbf{L} K)$ 的上同调 % P09-ERR-0491
在次数 $\geq m - 1$ 为零。因此
$$
R\Gamma(X, P \otimes^\mathbf{L} K) \to R\Gamma(X, E \otimes^\mathbf{L} K)
$$
在次数 $\geq m$ 的上同调上是同构。
根据假设，源拟凝聚。
因此 $R\Gamma(X, E \otimes^\mathbf{L} K)$
如所需地为 $m$-拟凝聚。
\end{proof}

\begin{lemma}
\label{spaces-perfect-lemma-base-change-RHom-perfect}
设 $S$ 为概形。
设 $f : X \to Y$ 是 $S$ 上代数空间之间的有限表现态射。
设 $E \in D(\mathcal{O}_X)$ 为完美对象。设 $\mathcal{G}^\bullet$
为由有限表现 $\mathcal{O}_X$-模组成的有界复形，
它在 $Y$ 上平坦，且支撑在 $Y$ 上适当。那么
$$
K = Rf_*R\SheafHom(E, \mathcal{G}^\bullet)
$$
是 $D(\mathcal{O}_Y)$ 的完美对象，并且它的构造与任意基变换可交换。
\end{lemma}

\begin{proof}
关于基变换的陈述是引理 \ref{spaces-perfect-lemma-base-change-RHom}。
因此只需证明 $K$ 是完美对象。若 $Y$ 诺特，则这由引理
\ref{spaces-perfect-lemma-ext-perfect} 得出。
我们将用诺特逼近把问题归约到这种情形。
建议读者跳过本证明的其余部分。

\medskip\noindent
问题在 $Y$ 上是局部的，故可假设 $Y$ 仿射。
设 $Y = \Spec(R)$。把 $R = \colim R_i$ 写成诺特环
$R_i$ 的滤过余极限。由《空间的极限》引理
\ref{spaces-limits-lemma-descend-finite-presentation}，
存在 $i$ 和代数空间 $X_i$，它在 $R_i$ 上有限表现，且到
$R$ 的基变换是 $X$。由《空间的极限》引理
\ref{spaces-limits-lemma-descend-modules-finite-presentation}，
增大 $i$ 后，可以假设存在由有限表现
$\mathcal{O}_{X_i}$-模组成的有界复形 $\mathcal{G}_i^\bullet$，% P09-ERR-0492
它到 $X$ 的拉回是 $\mathcal{G}$。增大 $i$ 后，% P09-ERR-0493
可以假设 $\mathcal{G}_i^n$ 在 $R_i$ 上平坦；见《空间的极限》引理
\ref{spaces-limits-lemma-descend-flat}。
增大 $i$ 后，可以假设 $\mathcal{G}_i^n$ 的支撑在 $R_i$ 上适当；
见《空间的极限》引理
\ref{spaces-limits-lemma-eventually-proper-support}。
最后，由引理 \ref{spaces-perfect-lemma-descend-perfect}，增大 $i$ 后，
可以假设存在完美对象 $E_i$（属于 $D(\mathcal{O}_{X_i})$），
它到 $X$ 的拉回是 $E$。把引理
\ref{spaces-perfect-lemma-compute-ext-perfect} 应用于
$X_i \to \Spec(R_i)$、$E_i$、$\mathcal{G}_i^\bullet$，
并使用已经证明的基变换性质，即得结果。
\end{proof}

\section{完美复形}
\label{spaces-perfect-section-perfect-complexes}

\noindent
先讨论完美复形的 Betti 数之跳跃轨迹。
首先必须定义 Betti 数。

\medskip\noindent
设 $S$ 为概形。设 $X$ 为 $S$ 上的代数空间。
设 $E$ 为 $D(\mathcal{O}_X)$ 的对象。
设 $x \in |X|$。我们要定义
$\beta_i(x) \in \{0, 1, 2, \ldots \} \cup \{\infty\}$。
为此，选取态射 $f : \Spec(k) \to X$，它属于 $x$ 的等价类。
于是 $Lf^*E$ 是 $D(\Spec(k)_\etale, \mathcal{O})$ 的对象。
由《平展上同调》引理
\ref{etale-cohomology-lemma-all-modules-quasi-coherent} 和定理
\ref{etale-cohomology-theorem-quasi-coherent}，得到
$D(\Spec(k)_\etale, \mathcal{O}) = D(k)$，即 $k$-向量空间的导出范畴。
因此 $Lf^*E$ 是 $k$-向量空间的复形，可以取
$\beta_i(x) = \dim_k H^i(Lf^*E)$。容易看出，这不依赖于
$x$ 中代表元的选取。此外，若 $X$ 是概形，这与《概形的导出范畴》第
\ref{perfect-section-perfect-complexes} 节中使用的概念相同。

\begin{lemma}
\label{spaces-perfect-lemma-jump-loci}
设 $S$ 为概形。设 $X$ 为 $S$ 上的代数空间。
设 $E \in D(\mathcal{O}_X)$ 拟凝聚（例如完美）。
对任意 $i \in \mathbf{Z}$，考虑上面定义的函数
$$
\beta_i : |X| \longrightarrow \{0, 1, 2, \ldots\}
$$
。那么
\begin{enumerate}
\item $\beta_i$ 的构造与任意基变换可交换；
\item 函数 $\beta_i$ 上半连续；且
\item $\beta_i$ 的水平集在平展局部可构造。
\end{enumerate}
\end{lemma}

\begin{proof}
选取概形 $U$ 和满射平展态射 $\varphi : U \to X$。
那么 $L\varphi^*E$ 是概形 $U$ 上的拟凝聚复形（使用引理
\ref{spaces-perfect-lemma-descend-pseudo-coherent}），从而可以应用概形情形的结果；
见《概形的导出范畴》引理 \ref{perfect-lemma-jump-loci}。
第 (3) 部分意指这些水平集到 $U$ 的逆像局部可构造；见
《空间的性质》定义
\ref{spaces-properties-definition-locally-constructible}。
\end{proof}

\begin{lemma}
\label{spaces-perfect-lemma-jump-loci-geometric}
设 $Y$ 为概形，并设 $X$ 为 $Y$ 上的代数空间，使得结构态射
$f : X \to Y$ 平坦、适当且有限表现。
设 $\mathcal{F}$ 为有限表现 $\mathcal{O}_X$-模，且在 $Y$ 上平坦。
固定 $i \in \mathbf{Z}$，考虑函数
$$
\beta_i : |Y| \to \{0, 1, 2, \ldots\},\quad
y \longmapsto \dim_{\kappa(y)} H^i(X_y, \mathcal{F}_y)
$$
。那么
\begin{enumerate}
\item $\beta_i$ 的构造与任意基变换可交换；
\item 函数 $\beta_i$ 上半连续；且
\item $\beta_i$ 的水平集在 $Y$ 中局部可构造。
\end{enumerate}
\end{lemma}

\begin{proof}
由上同调与基变换（更确切地说，由引理
\ref{spaces-perfect-lemma-flat-proper-perfect-direct-image-general}），对象
$K = Rf_*\mathcal{F}$ 是 $Y$ 的导出范畴中的完美对象，
并且它的构造与任意基变换可交换。特别地，
$$
H^i(X_y, \mathcal{F}_y) = H^i(K \otimes_{\mathcal{O}_Y}^\mathbf{L} \kappa(y))
$$
。因此本引理由引理 \ref{spaces-perfect-lemma-jump-loci} 得出。
\end{proof}

\begin{lemma}
\label{spaces-perfect-lemma-chi-locally-constant}
设 $S$ 为概形。设 $X$ 为 $S$ 上的代数空间。
设 $E \in D(\mathcal{O}_X)$ 完美。函数
$$
\chi_E : |X| \longrightarrow \mathbf{Z},\quad
x \longmapsto \sum (-1)^i \beta_i(x)
$$
在 $X$ 上局部常值。
\end{lemma}

\begin{proof}
略。提示：
由概形情形通过平展局部化得出。见《概形的导出范畴》引理
\ref{perfect-lemma-chi-locally-constant}。
\end{proof}

\begin{lemma}
\label{spaces-perfect-lemma-open-where-cohomology-in-degree-i-rank-r}
设 $S$ 为概形。设 $X$ 为 $S$ 上的代数空间。
设 $E \in D(\mathcal{O}_X)$ 完美。
给定 $i, r \in \mathbf{Z}$，存在由下列性质刻画的开子空间
$U \subset X$：
\begin{enumerate}
\item $E|_U \cong H^i(E|_U)[-i]$，且 $H^i(E|_U)$ 是局部自由
$\mathcal{O}_U$-模，秩为 $r$；
\item 态射 $f : Y \to X$ 经由 $U$ 分解，当且仅当
$Lf^*E$ 同构于局部自由模，秩为 $r$，并置于次数 $i$。
\end{enumerate}
\end{lemma}

\begin{proof}
略。提示：
由概形情形通过平展局部化得出。见《概形的导出范畴》引理
\ref{perfect-lemma-open-where-cohomology-in-degree-i-rank-r}。
\end{proof}

\begin{lemma}
\label{spaces-perfect-lemma-open-where-cohomology-in-degree-i-rank-r-geometric}
设 $S$ 为概形。
设 $f : X \to Y$ 是 $S$ 上代数空间的态射，且适当、平坦、有限表现。
设 $\mathcal{F}$ 为有限表现 $\mathcal{O}_X$-模，且在 $Y$ 上平坦。
固定 $i, r \in \mathbf{Z}$。那么存在开子空间
$V \subset Y$，具有如下性质：
态射 $T \to Y$ 经由 $V$ 分解，当且仅当
$Rf_{T, *}\mathcal{F}_T$ 同构于有限局部自由模，秩为 $r$，
并置于次数 $i$。
\end{lemma}

\begin{proof}
由上同调与基变换（引理
\ref{spaces-perfect-lemma-flat-proper-perfect-direct-image-general}），对象
$K = Rf_*\mathcal{F}$ 是 $Y$ 的导出范畴中的完美对象，
并且它的构造与任意基变换可交换。因此本引理由引理
\ref{spaces-perfect-lemma-open-where-cohomology-in-degree-i-rank-r} 立即得出。
\end{proof}

\begin{lemma}
\label{spaces-perfect-lemma-locally-closed-where-H0-locally-free}
设 $S$ 为概形。设 $X$ 为 $S$ 上的代数空间。
设 $E \in D(\mathcal{O}_X)$ 完美，且 Tor 振幅在
$[a, b]$ 中，其中 $a, b \in \mathbf{Z}$。
设 $r \geq 0$。
那么存在由下列性质刻画的局部闭子空间 $j : Z \to X$：
\begin{enumerate}
\item $H^a(Lj^*E)$ 是局部自由 $\mathcal{O}_Z$-模，秩为 $r$；且
\item 态射 $f : Y \to X$ 经由 $Z$ 分解，当且仅当对所有态射
$g : Y' \to Y$，$\mathcal{O}_{Y'}$-模
$H^a(L(f \circ g)^*E)$ 都局部自由且秩为 $r$。
\end{enumerate}
此外，$j : Z \to X$ 是有限表现的，并且
\begin{enumerate}
\item[(3)] 若 $f : Y \to X$ 分解为
$Y \xrightarrow{g} Z \to X$，则
$H^a(Lf^*E) = g^*H^a(Lj^*E)$；
\item[(4)] 若 $\beta_a(x) \leq r$ 对所有 $x \in |X|$ 成立，
则 $j$ 是闭浸入；并且给定 $f : Y \to X$，下列条件等价：
\begin{enumerate}
\item $f : Y \to X$ 经由 $Z$ 分解；
\item $H^0(Lf^*E)$ 是局部自由 $\mathcal{O}_Y$-模，秩为 $r$；% P09-ERR-0494
\end{enumerate}
而当 $r = 1$ 时，它们还等价于
\begin{enumerate}
\item[(c)] $\mathcal{O}_Y \to
\SheafHom_{\mathcal{O}_Y}(H^0(Lf^*E), H^0(Lf^*E))$
为单射。
\end{enumerate}
\end{enumerate}
\end{lemma}

\begin{proof}
略。提示：
由概形情形通过平展局部化得出。见《概形的导出范畴》引理
\ref{perfect-lemma-locally-closed-where-H0-locally-free}。
\end{proof}

\begin{lemma}
\label{spaces-perfect-lemma-proper-flat-h0}
设 $S$ 为概形。
设 $f : X \to Y$ 是 $S$ 上代数空间的态射。假设
\begin{enumerate}
\item $f$ 适当、平坦且有限表现；并且
\item 对于一个态射 $\Spec(k) \to Y$，其中 $k$ 是域，有 % P09-ERR-0495
$k = H^0(X_k, \mathcal{O}_{X_k})$。
\end{enumerate}
那么
\begin{enumerate}
\item[(a)] $f_*\mathcal{O}_X = \mathcal{O}_Y$，且这在任意基变换后仍成立；
\item[(b)] 在 $Y$ 的平展局部有
$$
Rf_*\mathcal{O}_X = \mathcal{O}_Y \oplus P
$$
（在 $D(\mathcal{O}_Y)$ 中），其中 $P$ 完美且 Tor 振幅在
$[1, \infty)$ 中。
\end{enumerate}
\end{lemma}

\begin{proof}
只需在 $Y$ 的平展局部证明 (a) 和 (b)，故可以并且确实假设
$Y$ 是仿射概形。
由上同调与基变换
（引理 \ref{spaces-perfect-lemma-flat-proper-perfect-direct-image-general}），复形
$E = Rf_*\mathcal{O}_X$
完美，并且它的构造与任意基变换可交换。
特别地，对 $y \in Y$ 有
$H^0(E \otimes^\mathbf{L} \kappa(y)) =
H^0(X_y, \mathcal{O}_{X_y}) = \kappa(y)$。
因此，用引理 \ref{spaces-perfect-lemma-jump-loci} 的记号，
$\beta_0(y) \leq 1$ 对所有 $y \in Y$ 成立。应用引理
\ref{spaces-perfect-lemma-locally-closed-where-H0-locally-free}，取
$a = 0$ 和 $r = 1$。得到泛闭子概形
$j : Z \to Y$，其中 $H^0(Lj^*E)$ 可逆，并由该引理
(4)(a)、(b)、(c) 的等价性刻画。
由于 $E$ 的构造与基变换可交换，有
$$
Lf^*E = R\text{pr}_{1, *}\mathcal{O}_{X \times_Y X}
$$
。态射 $\text{pr}_1 : X \times_Y X$ 有截面，
即对角态射 $\Delta$（$X$ 在 $Y$ 上）。得到映射
$$
\mathcal{O}_X \longrightarrow R\text{pr}_{1, *}\mathcal{O}_{X \times_Y X}
\longrightarrow \mathcal{O}_X
$$
（在 $D(\mathcal{O}_X)$ 中），其复合是恒等态射。因此
$R\text{pr}_{1, *}\mathcal{O}_{X \times_Y X} = \mathcal{O}_X \oplus E'$
（在 $D(\mathcal{O}_X)$ 中）。从而 $\mathcal{O}_X$ 是
$H^0(Lf^*E)$ 的直和因子；由上述引理 (4)(c) 与 (4)(a) 的等价性，
得出 $X \to Y$ 经由 $Z$ 分解。
由于 $\{X \to Y\}$ 是 fppf 覆盖，有 $Z = Y$。
于是 $f_*\mathcal{O}_X$ 是可逆 $\mathcal{O}_Y$-模。
我们断定 $\mathcal{O}_Y \to f_*\mathcal{O}_X$ 是同构，因为若环同态
$A \to B$ 使 $B$ 作为 $A$-模可逆，则该同态是同构。
由于假设在基变换下保持，(a) 得证。

\medskip\noindent
(b) 的证明。上面已经看到，对每个 $y \in Y$，映射
$\mathcal{O}_Y \to H^0(E \otimes^\mathbf{L} \kappa(y))$ 满射。
因此可以应用《代数进阶》引理
\ref{more-algebra-lemma-better-cut-complex-in-two}，得出在 $y$ 的
一个开邻域上有分解
$Rf_*\mathcal{O}_X = \mathcal{O}_Y \oplus P$。
\end{proof}

\begin{lemma}
\label{spaces-perfect-lemma-proper-flat-geom-red-connected}
设 $S$ 为概形。
设 $f : X \to Y$ 是 $S$ 上代数空间的态射。假设
\begin{enumerate}
\item $f$ 适当、平坦且有限表现；并且
\item $f$ 的几何纤维是既约且连通的。
\end{enumerate}
那么 $f_*\mathcal{O}_X = \mathcal{O}_Y$，且这在任意基变换后仍成立。
\end{lemma}

\begin{proof}
由引理 \ref{spaces-perfect-lemma-proper-flat-h0}，只需证明
$k = H^0(X_k, \mathcal{O}_{X_k})$ 对所有态射
$\Spec(k) \to Y$（其中 $k$ 是域）成立。这由《域上的空间》引理
\ref{spaces-over-fields-lemma-proper-geometrically-reduced-global-sections}
以及 $X_k$ 几何连通且几何既约这一事实得出。
\end{proof}

\section{其他应用}
\label{spaces-perfect-section-other-applications}

\noindent
本节陈述并证明一些可由上述理论推出的结果。

\begin{lemma}
\label{spaces-perfect-lemma-countable-cohomology}
设 $S$ 为概形。设 $X$ 为 $S$ 上拟紧拟分离的代数空间。
设 $K$ 为 $D_\QCoh(\mathcal{O}_X)$ 的对象，并且上同调层
$H^i(K)$ 在平展于 $X$ 的仿射概形上的截面集合可数。
那么对任意拟紧拟分离平展态射 $U \to X$ 和任意完美对象
$E$（属于 $D(\mathcal{O}_X)$），集合
$$
H^i(U, K \otimes^\mathbf{L} E),\quad \Ext^i(E|_U, K|_U)
$$
都可数。
\end{lemma}

\begin{proof}
由《位点上的上同调》引理
\ref{sites-cohomology-lemma-dual-perfect-complex}，
只需对群 $H^i(U, K \otimes^\mathbf{L} E)$ 证明结果。
我们将使用归纳原理证明本引理；见引理
\ref{spaces-perfect-lemma-induction-principle}。

\medskip\noindent
当 $U = \Spec(A)$ 仿射时，结果由概形情形得出；
见《概形的导出范畴》引理
\ref{perfect-lemma-countable-cohomology}。

\medskip\noindent
为完成证明，只需说明：若 $(U \subset W, V \to W)$
是初等特异三角，% P09-ERR-0496; same source-defect class as P09-ERR-0408
并且结果对 $U$、$V$ 和 $U \times_W V$ 成立，则结果对 $W$ 成立。
这立即由 Mayer--Vietoris 序列得出；见引理
\ref{spaces-perfect-lemma-unbounded-mayer-vietoris}。
\end{proof}

\begin{lemma}
\label{spaces-perfect-lemma-countable}
设 $S$ 为概形。
设 $X$ 为 $S$ 上拟紧拟分离的代数空间。
假设 $\mathcal{O}_X$ 在平展于 $X$ 的仿射概形上的截面集合可数。
设 $K$ 为 $D_\QCoh(\mathcal{O}_X)$ 的对象。下列条件等价：
\begin{enumerate}
\item $K = \text{hocolim} E_n$，其中 $E_n$ 是
$D(\mathcal{O}_X)$ 的完美对象；且
\item 上同调层 $H^i(K)$ 在平展于 $X$ 的仿射概形上的截面集合可数。
\end{enumerate}
\end{lemma}

\begin{proof}
若 (1) 成立，则 (2) 成立，因为同伦余极限与取上同调层可交换，% P09-ERR-0497
（由《导出范畴》引理
\ref{derived-lemma-cohomology-of-hocolim}），并且完美复形局部同构于
有限自由 $\mathcal{O}_X$-模的有限复形，所以根据关于 $X$ 的假设满足 (2)。

\medskip\noindent
假设 (2)。
选取 K-内射复形 $\mathcal{K}^\bullet$ 来代表 $K$。
选取完美生成元 $E$（属于 $D_\QCoh(\mathcal{O}_X)$），并用 K-内射复形
$\mathcal{I}^\bullet$ 代表它。根据定理
\ref{spaces-perfect-theorem-DQCoh-is-Ddga} 及其证明，存在三角范畴等价
$F : D_\QCoh(\mathcal{O}_X) \to D(A, \text{d})$，
其中 $(A, \text{d})$ 是微分分次代数
$$
(A, \text{d}) =
\Hom_{\text{Comp}^{dg}(\mathcal{O}_X)}
(\mathcal{I}^\bullet, \mathcal{I}^\bullet)
$$
，且它把 $K$ 映到微分分次模
$$
M = \Hom_{\text{Comp}^{dg}(\mathcal{O}_X)}
(\mathcal{I}^\bullet, \mathcal{K}^\bullet)
$$
。注意 $H^i(A) = \Ext^i(E, E)$ 且
$H^i(M) = \Ext^i(E, K)$。
此外，由于 $F$ 是等价，它及其拟逆都与同伦余极限可交换。
因此，只需把 $M$ 写成 $D(A, \text{d})$ 的紧对象之同伦余极限。
由《微分分次代数》引理 \ref{dga-lemma-countable}，只需证明 % P09-ERR-0498
$\Ext^i(E, E)$ 和 $\Ext^i(E, K)$ 对每个 $i$ 都可数。
这由引理 \ref{spaces-perfect-lemma-countable-cohomology} 得出。
\end{proof}

\begin{lemma}
\label{spaces-perfect-lemma-computing-sections-as-colim}
设 $A$ 为环。设 $f : U \to X$ 是 $A$ 上代数空间之间的
平坦有限表现态射。那么
\begin{enumerate}
\item 存在由完美对象 $L_n$（属于 $D(\mathcal{O}_X)$）构成的逆系统，
使得
$$
R\Gamma(U, Lf^*K) = \text{hocolim}\ R\Hom_X(L_n, K)
$$
在 $D(A)$ 中对 $K$（属于 $D_\QCoh(\mathcal{O}_X)$）函子性成立；且
\item 存在由完美对象 $E_n$（属于 $D(\mathcal{O}_X)$）构成的系统，
使得
$$
R\Gamma(U, Lf^*K) = \text{hocolim}\ R\Gamma(X, E_n \otimes^\mathbf{L} K)
$$
在 $D(A)$ 中对 $K$（属于 $D_\QCoh(\mathcal{O}_X)$）函子性成立。
\end{enumerate}
\end{lemma}

\begin{proof}
由引理 \ref{spaces-perfect-lemma-cohomology-base-change}，
$$
R\Gamma(U, Lf^*K) = R\Gamma(X, Rf_*\mathcal{O}_U \otimes^\mathbf{L} K)
$$
对 $K$ 函子性成立。注意，$R\Gamma(X, -)$ 与同伦余极限可交换，
因为由引理 \ref{spaces-perfect-lemma-quasi-coherence-pushforward-direct-sums}，
它与直和可交换。类似地，$- \otimes^\mathbf{L} K$
与导出余极限可交换，因为 $- \otimes^\mathbf{L} K$ 与直和可交换
（因为 $D(\mathcal{O}_X)$ 中的直和由代表复形的直和给出）。
因此，为证明 (2)，只需把
$Rf_*\mathcal{O}_U = \text{hocolim} E_n$ 写成由完美对象
$E_n$（属于 $D(\mathcal{O}_X)$）组成的系统。
做到这一点后，置 $L_n = E_n^\vee$ 即得 (1)；见《位点上的上同调》引理
\ref{sites-cohomology-lemma-dual-perfect-complex}。

\medskip\noindent
把 $A = \colim A_i$ 写成环 $A_i$ 的余极限，它们
在 $\mathbf{Z}$ 上有限型。由《空间的极限》引理
\ref{spaces-limits-lemma-descend-finite-presentation}，
可以找到 $i$ 和有限表现态射
$U_i \to X_i \to \Spec(A_i)$，其到 $\Spec(A)$ 的基变换恢复
$U \to X \to \Spec(A)$。
增大 $i$ 后，可以假设 $f_i : U_i \to X_i$ 平坦；见《空间的极限》引理
\ref{spaces-limits-lemma-descend-flat}。
由引理 \ref{spaces-perfect-lemma-compare-base-change}，对象
$Rf_{i, *}\mathcal{O}_{U_i}$ 沿 $g : X \to X_i$ 的导出拉回等于
$Rf_*\mathcal{O}_U$。
由于 $Lg^*$ 与导出余极限可交换，只需对 $f_i$ 证明所需结论。
因此可以假设 $U$ 和 $X$ 在 $\mathbf{Z}$ 上有限型。

\medskip\noindent
假设 $f : U \to X$ 是 $\mathbf{Z}$ 上有限型代数空间的态射。
为完成证明，我们将说明 $Rf_*\mathcal{O}_U$ 是完美复形的同伦余极限。
为此应用引理 \ref{spaces-perfect-lemma-countable}。所以只需证明
$R^if_*\mathcal{O}_U$ 在平展于 $X$ 的仿射概形上的截面集合可数。
这由引理 \ref{spaces-perfect-lemma-countable-cohomology} 应用于结构层得出。
\end{proof}









\section{分解性质}
\label{spaces-perfect-section-resolution-property}

\noindent
本节是《概形的导出范畴》第
\ref{perfect-section-resolution-property} 节在代数空间情形的对应版本；
请先阅读该节。目前尚不清楚域上的光滑固有代数空间是否总具有分解性质，
抑或并非如此。如知此问题的答案，请发送邮件至
\href{mailto:stacks.project@gmail.com}{stacks.project@gmail.com}。

\medskip\noindent
尽管对一般代数空间考虑如下概念几乎没有意义，我们仍可作出这个定义。

\begin{definition}
\label{spaces-perfect-definition-resolution-property}
设 $S$ 为概形，$X$ 为 $S$ 上的代数空间。称 $X$ 具有
{\it 分解性质}，如果每个有限型拟凝聚 $\mathcal{O}_X$-模
都是某个有限局部自由 $\mathcal{O}_X$-模的商。
\end{definition}

\noindent
若 $X$ 是拟紧拟分离代数空间，则只需检验每个有限表示的
$\mathcal{O}_X$-模（自动拟凝聚）都是某个有限局部自由
$\mathcal{O}_X$-模的商；见《空间的极限》引理 % P09-ERR-0499
\ref{spaces-limits-lemma-finite-directed-colimit-surjective-maps}。
若 $X$ 是诺特代数空间，则有限型拟凝聚模恰为凝聚
$\mathcal{O}_X$-模；见《空间的上同调》引理
\ref{spaces-cohomology-lemma-coherent-Noetherian}。

\begin{lemma}
\label{spaces-perfect-lemma-resolution-property-ample-relative}
设 $S$ 为概形。设 $f : X \to Y$ 是 $S$ 上代数空间之间的态射。假设
\begin{enumerate}
\item $Y$ 拟紧、拟分离且具有分解性质；
\item 存在一个 $f$-丰沛可逆模，其定义在 $X$ 上
（《空间上的除子》定义
\ref{spaces-divisors-definition-relatively-ample}）。
\end{enumerate}
则 $X$ 具有分解性质。
\end{lemma}

\begin{proof}
令 $\mathcal{F}$ 为有限型拟凝聚 $\mathcal{O}_X$-模，
$\mathcal{L}$ 为 $f$-丰沛可逆模。取仿射概形 $V$ 以及满平展态射
$V \to Y$。置 $U = V \times_Y X$，则 $\mathcal{L}|_U$ 在 $U$ 上丰沛。
由《性质》命题
\ref{properties-proposition-characterize-ample}，存在有限多个映射 % P09-ERR-0500
$s_i : \mathcal{L}^{\otimes n_i}|_U \to \mathcal{F}|_U$，
它们联合为满。考虑拟凝聚 $\mathcal{O}_Y$-模
$$
\mathcal{H}_n =
f_*(\mathcal{F} \otimes_{\mathcal{O}_X} \mathcal{L}^{\otimes n})
$$
可将 $s_i$ 视为 $V$ 上的一个截面；这里所取的层是
$\mathcal{H}_{-n_i}$。
假设可以找到有限局部自由 $\mathcal{O}_Y$-模 $\mathcal{E}_i$ 和映射
$\mathcal{E}_i \to \mathcal{H}_{-n_i}$，使 $s_i$ 属于其像。
那么相应的映射
$$
f^*\mathcal{E}_i
\otimes_{\mathcal{O}_X}
\mathcal{L}^{\otimes n_i} \longrightarrow \mathcal{F}
$$
将联合为满，于是引理得证。由《空间的极限》引理
\ref{spaces-limits-lemma-directed-colimit-finite-type}，
对每个 $i$ 都能找到有限型拟凝聚子模
$\mathcal{H}'_i \subset  \mathcal{H}_{-n_i}$，
它包含截面 $s_i$（定义在 $V$ 上）。
因此，$Y$ 的分解性质给出满射
$\mathcal{E}_i \to \mathcal{H}'_i$，结论随即成立。
\end{proof}

\begin{lemma}
\label{spaces-perfect-lemma-resolution-property-goes-up-affine}
设 $S$ 为概形。设 $f : X \to Y$ 是 $S$ 上代数空间之间的仿射或拟仿射
态射，并且 $Y$ 拟紧拟分离。若 $Y$ 具有分解性质，则 $X$ 亦然。
\end{lemma}

\begin{proof}
由《空间上的除子》引理
\ref{spaces-divisors-lemma-quasi-affine-O-ample}，
这是引理 \ref{spaces-perfect-lemma-resolution-property-ample-relative} 的特殊情形。
\end{proof}

\noindent
下面给出一个可以证明分解性质向下传递的情形。

\begin{lemma}
\label{spaces-perfect-lemma-resolution-property-goes-down-finite-flat}
设 $S$ 为概形。设 $f : X \to Y$ 是 $S$ 上代数空间之间的满有限局部自由
态射。若 $X$ 具有分解性质，则 $Y$ 亦然。
\end{lemma}

\begin{proof}
该条件意味着 $f$ 仿射，且 $f_*\mathcal{O}_X$ 是秩为正的有限局部自由
$\mathcal{O}_Y$-模。令 $\mathcal{G}$ 为有限型拟凝聚
$\mathcal{O}_Y$-模。由假设，存在满射
$\mathcal{E} \to f^*\mathcal{G}$，其源是有限局部自由
$\mathcal{O}_X$-模 $\mathcal{E}$。由于 $f_*$ 正合
（《空间的上同调》第 \ref{spaces-cohomology-section-finite-morphisms} 节），
我们得到满射
$$
f_*\mathcal{E}
\longrightarrow
f_*f^*\mathcal{G} = \mathcal{G} \otimes_{\mathcal{O}_Y} f_*\mathcal{O}_X
$$
取对偶后，得到满射 % P09-ERR-0501
$$
f_*\mathcal{E}
\otimes_{\mathcal{O}_Y}
\SheafHom_{\mathcal{O}_Y}(f_*\mathcal{O}_X, \mathcal{O}_Y)
\longrightarrow
\mathcal{G}
$$
由于 $f_*\mathcal{E}$ 有限局部自由，结论成立。
\end{proof}

\noindent
关于代数空间分解性质的更多内容，见《空间态射进阶》第
\ref{spaces-more-morphisms-section-resolution-property} 节。













\section{判定有界性}
\label{spaces-perfect-section-detecting-boundedness}

\noindent
本节将说明，第 \ref{spaces-perfect-section-generating} 节所构造的拟紧拟分离概形之
$D_\QCoh$ 的紧生成元具有一种特殊性质。% P09-ERR-0502
建议先阅读该节，因为它与本节非常相似。

\begin{lemma}
\label{spaces-perfect-lemma-bounded-truncation}
设 $S$ 为概形。设 $X$ 为 $S$ 上的拟紧拟分离代数空间。
设 $P \in D_{perf}(\mathcal{O}_X)$ 且
$E \in D_{\QCoh}(\mathcal{O}_X)$。设 $a \in \mathbf{Z}$。
以下条件等价：
\begin{enumerate}
\item $\Hom_{D(\mathcal{O}_X)}(P[-i], E) = 0$
对 $i \gg 0$ 成立；且
\item $\Hom_{D(\mathcal{O}_X)}(P[-i], \tau_{\geq a} E) = 0$
对 $i \gg 0$ 成立。
\end{enumerate}
\end{lemma}

\begin{proof}
利用三角 $ \tau_{< a} E \to E \to \tau_{\geq a} E \to$，
可知只要能证明
$$
\Hom_{D(\mathcal{O}_X)}(P[-i], \tau_{< a} E) =
\Hom_{D(\mathcal{O}_X)}(P, (\tau_{< a} E)[i]) = 0
$$
对 $i \gg 0$ 成立，所述等价性就成立。由于 $P$ 完美，
这由引理 \ref{spaces-perfect-lemma-ext-from-perfect-into-bounded-QCoh} 得出。
\end{proof}

\begin{lemma}
\label{spaces-perfect-lemma-bounded-below-truncation}
设 $S$ 为概形。设 $X$ 为 $S$ 上的拟紧拟分离代数空间。
设 $P \in D_{perf}(\mathcal{O}_X)$ 且
$E \in D_{\QCoh}(\mathcal{O}_X)$。设 $a \in \mathbf{Z}$。
以下条件等价：
\begin{enumerate}
\item $\Hom_{D(\mathcal{O}_X)}(P[-i], E) = 0$
对 $i \ll 0$ 成立；且
\item $\Hom_{D(\mathcal{O}_X)}(P[-i], \tau_{\leq a} E) = 0$
对 $i \ll 0$ 成立。
\end{enumerate}
\end{lemma}

\begin{proof}
利用三角 $ \tau_{\leq a} E \to E \to \tau_{> a} E \to$，
可知只要能证明
$$
\Hom_{D(\mathcal{O}_X)}(P[-i], \tau_{> a} E) =
\Hom_{D(\mathcal{O}_X)}(P, (\tau_{> a} E)[i]) = 0
$$
对 $i \ll 0$ 成立，所述等价性就成立。由于 $P$ 完美，
这由引理 \ref{spaces-perfect-lemma-ext-from-perfect-into-bounded-QCoh} 得出。
\end{proof}

\begin{proposition}
\label{spaces-perfect-proposition-detecting-bounded-above}
设 $S$ 为概形。设 $X$ 为 $S$ 上的拟紧拟分离代数空间。
设 $G \in D_{perf}(\mathcal{O}_X)$ 是生成
$D_\QCoh (\mathcal{O}_X)$ 的完美复形。设
$E \in D_\QCoh (\mathcal{O}_X)$。以下条件等价：
\begin{enumerate}
\item $E \in D^-_\QCoh (\mathcal{O}_X)$；
\item $\Hom_{D(\mathcal{O}_X)}(G[-i], E) = 0$
对 $i \gg 0$ 成立；
\item $\Ext^i_X(G, E) = 0$ 对 $i \gg 0$ 成立；
\item $R\Hom_X(G, E)$ 属于 $D^-(\mathbf{Z})$；
\item $H^i(X, G^\vee \otimes_{\mathcal{O}_X}^\mathbf{L} E) = 0$
对 $i \gg 0$ 成立；
\item $R\Gamma(X, G^\vee \otimes_{\mathcal{O}_X}^\mathbf{L} E)$
属于 $D^-(\mathbf{Z})$；
\item 对每个完美对象 $P$（属于 $D(\mathcal{O}_X)$），
\begin{enumerate}
\item 将 $G$ 换成 $P$ 后，断言 (2)、(3)、(4) 成立；且
\item $H^i(X, P \otimes_{\mathcal{O}_X}^\mathbf{L} E) = 0$
对 $i \gg 0$ 成立；
\item $R\Gamma(X, P \otimes_{\mathcal{O}_X}^\mathbf{L} E)$
属于 $D^-(\mathbf{Z})$。
\end{enumerate}
\end{enumerate}
\end{proposition}

\begin{proof}
假设 (1) 成立。由于
$\Hom_{D(\mathcal{O}_X)}(G[-i], E) = \Hom_{D(\mathcal{O}_X)}(G, E[i])$，
由引理 \ref{spaces-perfect-lemma-ext-from-perfect-into-bounded-QCoh}，
这在 $i \gg 0$ 时为零。这证明了 (1) 推出 (2)。

\medskip\noindent
由《位点上的上同调》第 \ref{sites-cohomology-section-global-RHom} 节
的讨论，(2)、(3)、(4) 等价。按定义，
$H^i(X, -) = H^i(R\Gamma(X, -))$，故 (5) 和 (6) 等价。
由《位点上的上同调》引理
\ref{sites-cohomology-lemma-dual-perfect-complex} 以及等式
$(G[-i])^\vee = G^\vee[i]$，等价条件 (2)、(3)、(4)
又等价于等价条件 (5)、(6)。

\medskip\noindent
显然 (7) 推出 (2)。反过来，我们证明等价条件 (2)--(6) 推出 (7)。
回顾注 \ref{spaces-perfect-remark-classical-generator}：$G$ 是
$D_{perf}(\mathcal{O}_X)$ 的经典生成元。对
$P \in D_{perf}(\mathcal{O}_X)$，令 $T(P)$ 表示断言
$R\Hom_X(P, E)$ 属于 $D^-(\mathbf{Z})$。
显然，$T$ 对直和、直和项和移位封闭，并且对特异三角满足两出三性质。
因此，由《导出范畴》注 \ref{derived-remark-check-on-generator}，
(4) 推出 $T$ 在整个 $D_{perf}(\mathcal{O}_X)$ 上成立。
同一论证适用于所有其他性质；但对性质 (7)(b) 和 (7)(c)，% P09-ERR-0503
还要用到 $P \mapsto P^\vee$ 是 $D_{perf}(\mathcal{O}_X)$ 的自等价。
略去这个小细节。

\medskip\noindent
我们将利用引理 \ref{spaces-perfect-lemma-induction-principle} 的归纳原理，
证明等价条件 (2)--(7) 推出 (1)。

\medskip\noindent
首先，我们证明 (2)--(7) $\Rightarrow$ (1) 在 $X$ 仿射时成立。
这由概形的情形得出；见《概形的导出范畴》命题
\ref{perfect-proposition-detecting-bounded-above}。

\medskip\noindent
现在假设 $(U \subset X, j : V \to X)$ 是 $S$ 上拟紧拟分离
代数空间的初等特异方形，并且已知蕴涵 (2)--(7) $\Rightarrow$ (1)
对 $U$、$V$ 和 $U \times_X V$ 成立。
为完成证明，必须说明蕴涵 (2)--(7) $\Rightarrow$ (1) 对 $X$ 成立。设
$E \in D_\QCoh(\mathcal{O}_X)$ 满足 (2)--(7)。
由引理 \ref{spaces-perfect-lemma-direct-summand-of-a-restriction} 和定理
\ref{spaces-perfect-theorem-bondal-van-den-Bergh}，存在完美复形 $Q$（在 $X$ 上），
使 $Q|_U$ 生成 $D_\QCoh (\mathcal{O}_U)$。

\medskip\noindent
写 $V = \Spec(A)$。令 $Z \subset V$ 为 $X \setminus U$ 的逆像
所对应的约化闭子概形；它同构地映到该逆像
（见定义 \ref{spaces-perfect-definition-elementary-distinguished-square}）。
这是 $V$ 的逆紧闭子集。取 $f_1, \ldots, f_r \in A$，使得
$Z = V(f_1, \ldots, f_r)$。令 $K \in D(\mathcal{O}_V)$ 为
由 $f_1, \ldots, f_r$ 在 $A$ 上的 Koszul 复形所对应的完美对象。
注意，由于 $K$ 支撑在 $Z$ 上，其正像
$K' = Rj_*K$ 是 $D(\mathcal{O}_X)$ 的完美对象，并且它限制到
$V$ 后是 $K$（见引理 \ref{spaces-perfect-lemma-pushforward-perfect} 和
\ref{spaces-perfect-lemma-pushforward-with-support-in-open}）。
由假设，$R\Hom_{\mathcal{O}_X}(Q, E)$ 和
$R\Hom_{\mathcal{O}_X}(K', E)$ 都有上界。

\medskip\noindent
由引理 \ref{spaces-perfect-lemma-pushforward-with-support-in-open}，有
$K' =  j_!K$，因而
$$
\Hom_{D(\mathcal{O}_X)}(K'[-i], E) = \Hom_{D(\mathcal{O}_V)}(K[-i], E|_V) = 0
$$
对 $i \gg 0$ 成立。因此，可将《概形的导出范畴》引理
\ref{perfect-lemma-orthogonal-koszul-first-variant} 应用于 $E|_V$，
得到整数 $a$，使得
$\tau_{\geq a}(E|_V) =
\tau_{\geq a} R (U \times_X V \to V)_* (E|_{U \times_X V})$。
于是 $\tau_{\geq a} E = \tau_{\geq a} R (U \to X)_* (E |_U)$
（限制到 $U$ 和 $V$ 后检验典范映射是同构）。
因而两次使用引理 \ref{spaces-perfect-lemma-bounded-truncation} 可得
$$
\Hom_{D(\mathcal{O}_U)}(Q|_U [-i], E|_U) =
\Hom_{D(\mathcal{O}_X)}(Q[-i], R (U \to X)_* (E|_U)) = 0
$$
对 $i \gg 0$ 成立。因为该命题对 $U$ 和生成元 $Q|_U$ 成立，
故 $E|_U \in D^-_\QCoh(\mathcal{O}_U)$。又因为函子
$R (U \to X)_*$ 保持 $D^-_\QCoh$
（由引理 \ref{spaces-perfect-lemma-quasi-coherence-direct-image}），所以
$\tau_{\geq a}E \in D^-_\QCoh(\mathcal{O}_X)$。因此
$E \in D^-_\QCoh (\mathcal{O}_X)$。
\end{proof}

\begin{proposition}
\label{spaces-perfect-proposition-detecting-bounded-below}
设 $S$ 为概形。设 $X$ 为 $S$ 上的拟紧拟分离代数空间。
设 $G \in D_{perf}(\mathcal{O}_X)$ 是生成
$D_\QCoh (\mathcal{O}_X)$ 的完美复形。设
$E \in D_\QCoh (\mathcal{O}_X)$。以下条件等价：
\begin{enumerate}
\item $E \in D^+_\QCoh (\mathcal{O}_X)$；
\item $\Hom_{D(\mathcal{O}_X)}(G[-i], E) = 0$
对 $i \ll 0$ 成立；
\item $\Ext^i_X(G, E) = 0$ 对 $i \ll 0$ 成立；
\item $R\Hom_X(G, E)$ 属于 $D^+(\mathbf{Z})$；
\item $H^i(X, G^\vee \otimes_{\mathcal{O}_X}^\mathbf{L} E) = 0$
对 $i \ll 0$ 成立；
\item $R\Gamma(X, G^\vee \otimes_{\mathcal{O}_X}^\mathbf{L} E)$
属于 $D^+(\mathbf{Z})$；
\item 对每个完美对象 $P$（属于 $D(\mathcal{O}_X)$），
\begin{enumerate}
\item 将 $G$ 换成 $P$ 后，断言 (2)、(3)、(4) 成立；且
\item $H^i(X, P \otimes_{\mathcal{O}_X}^\mathbf{L} E) = 0$
对 $i \ll 0$ 成立；
\item $R\Gamma(X, P \otimes_{\mathcal{O}_X}^\mathbf{L} E)$
属于 $D^+(\mathbf{Z})$。
\end{enumerate}
\end{enumerate}
\end{proposition}

\begin{proof}
假设 (1) 成立。由于
$\Hom_{D(\mathcal{O}_X)}(G[-i], E) = \Hom_{D(\mathcal{O}_X)}(G, E[i])$，
由引理 \ref{spaces-perfect-lemma-ext-from-perfect-into-bounded-QCoh}，
这在 $i \ll 0$ 时为零。这证明了 (1) 推出 (2)。

\medskip\noindent
由《位点上的上同调》第 \ref{sites-cohomology-section-global-RHom} 节
的讨论，(2)、(3)、(4) 等价。按定义，
$H^i(X, -) = H^i(R\Gamma(X, -))$，故 (5) 和 (6) 等价。
由《位点上的上同调》引理
\ref{sites-cohomology-lemma-dual-perfect-complex} 以及等式
$(G[-i])^\vee = G^\vee[i]$，等价条件 (2)、(3)、(4)
又等价于等价条件 (5)、(6)。

\medskip\noindent
显然 (7) 推出 (2)。反过来，我们证明等价条件 (2)--(6) 推出 (7)。
回顾注 \ref{spaces-perfect-remark-classical-generator}：$G$ 是
$D_{perf}(\mathcal{O}_X)$ 的经典生成元。对
$P \in D_{perf}(\mathcal{O}_X)$，令 $T(P)$ 表示断言
$R\Hom_X(P, E)$ 属于 $D^+(\mathbf{Z})$。
显然，$T$ 对直和、直和项和移位封闭，并且对特异三角满足两出三性质。
因此，由《导出范畴》注 \ref{derived-remark-check-on-generator}，
(4) 推出 $T$ 在整个 $D_{perf}(\mathcal{O}_X)$ 上成立。
同一论证适用于所有其他性质；但对性质 (7)(b) 和 (7)(c)，% P09-ERR-0503
还要用到 $P \mapsto P^\vee$ 是 $D_{perf}(\mathcal{O}_X)$ 的自等价。
略去这个小细节。

\medskip\noindent
我们将利用引理 \ref{spaces-perfect-lemma-induction-principle} 的归纳原理，
证明等价条件 (2)--(7) 推出 (1)。

\medskip\noindent
首先，我们证明 (2)--(7) $\Rightarrow$ (1) 在 $X$ 仿射时成立。
这由概形的情形得出；见《概形的导出范畴》命题
\ref{perfect-proposition-detecting-bounded-below}。

\medskip\noindent
现在假设 $(U \subset X, j : V \to X)$ 是 $S$ 上拟紧拟分离
代数空间的初等特异方形，并且已知蕴涵 (2)--(7) $\Rightarrow$ (1)
对 $U$、$V$ 和 $U \times_X V$ 成立。
为完成证明，必须说明蕴涵 (2)--(7) $\Rightarrow$ (1) 对 $X$ 成立。设
$E \in D_\QCoh(\mathcal{O}_X)$ 满足 (2)--(7)。
由引理 \ref{spaces-perfect-lemma-direct-summand-of-a-restriction} 和定理
\ref{spaces-perfect-theorem-bondal-van-den-Bergh}，存在完美复形 $Q$（在 $X$ 上），
使 $Q|_U$ 生成 $D_\QCoh (\mathcal{O}_U)$。

\medskip\noindent
写 $V = \Spec(A)$。令 $Z \subset V$ 为 $X \setminus U$ 的逆像
所对应的约化闭子概形；它同构地映到该逆像
（见定义 \ref{spaces-perfect-definition-elementary-distinguished-square}）。
这是 $V$ 的逆紧闭子集。取 $f_1, \ldots, f_r \in A$，使得
$Z = V(f_1, \ldots, f_r)$。令 $K \in D(\mathcal{O}_V)$ 为
由 $f_1, \ldots, f_r$ 在 $A$ 上的 Koszul 复形所对应的完美对象。
注意，由于 $K$ 支撑在 $Z$ 上，其正像
$K' = Rj_*K$ 是 $D(\mathcal{O}_X)$ 的完美对象，并且它限制到
$V$ 后是 $K$（见引理 \ref{spaces-perfect-lemma-pushforward-perfect} 和
\ref{spaces-perfect-lemma-pushforward-with-support-in-open}）。
由假设，$R\Hom_{\mathcal{O}_X}(Q, E)$ 和
$R\Hom_{\mathcal{O}_X}(K', E)$ 都有下界。

\medskip\noindent
由引理 \ref{spaces-perfect-lemma-pushforward-with-support-in-open}，有
$K' =  j_!K$，因而
$$
\Hom_{D(\mathcal{O}_X)}(K'[-i], E) = \Hom_{D(\mathcal{O}_V)}(K[-i], E|_V) = 0
$$
对 $i \ll 0$ 成立。因此，可将《概形的导出范畴》引理
\ref{perfect-lemma-orthogonal-koszul-second-variant} 应用于 $E|_V$，
得到整数 $a$，使得
$\tau_{\leq a}(E|_V) =
\tau_{\leq a} R (U \times_X V \to V)_* (E|_{U \times_X V})$。
于是 $\tau_{\leq a} E = \tau_{\leq a} R (U \to X)_* (E |_U)$
（限制到 $U$ 和 $V$ 后检验典范映射是同构）。
因而两次使用引理 \ref{spaces-perfect-lemma-bounded-below-truncation} 可得
$$
\Hom_{D(\mathcal{O}_U)}(Q|_U [-i], E|_U) =
\Hom_{D(\mathcal{O}_X)}(Q[-i], R (U \to X)_* (E|_U)) = 0
$$
对 $i \ll 0$ 成立。因为该命题对 $U$ 和生成元 $Q|_U$ 成立，
故 $E|_U \in D^+_\QCoh(\mathcal{O}_U)$。又因为函子
$R (U \to X)_*$ 保持 $D^+_\QCoh$
（由引理 \ref{spaces-perfect-lemma-quasi-coherence-direct-image}），所以
$\tau_{\leq a}E \in D^+_\QCoh(\mathcal{O}_X)$。因此
$E \in D^+_\QCoh (\mathcal{O}_X)$。
\end{proof}











\section{导出范畴中的拟凝聚对象}
\label{spaces-perfect-section-QC}

\noindent
设 $S$ 为概形，$X$ 为 $S$ 上的代数空间。回顾，
$X_{affine, \etale}$ 表示 $X_\etale$ 的仿射对象所成的范畴，
其拓扑由标准平展覆盖给出；见《空间的性质》定义
\ref{spaces-properties-definition-affine-etale-site}。
提醒读者，$X_{affine, \etale}$ 的拓扑斯就是 $X$ 的小平展拓扑斯；
见《空间的性质》引理
\ref{spaces-properties-lemma-alternative}。位点 $X_\etale$
带有结构层 $\mathcal{O}_X$；它限制到 $X_{affine, \etale}$ 后，
我们仍将其记作 $\mathcal{O}_X$。于是存在带环拓扑斯的等价
$$
(\Sh(X_{affine, \etale}), \mathcal{O}_X)
\longrightarrow
(\Sh(X_\etale), \mathcal{O}_X)
$$
见《空间上的下降》公式
(\ref{spaces-descent-equation-alternative-small-ringed})
以及《空间上的下降》第
\ref{spaces-descent-section-alternative-quasi-coherent} 节中的讨论。

\medskip\noindent
本节以 $X_{affine}$ 表示 $X_{affine, \etale}$ 的底层范畴，
并赋予它混沌拓扑，即层与预层相同的拓扑。特别地，结构层
$\mathcal{O}_X$ 也成为 $X_{affine}$ 上的层。
我们得到带环位点的态射 % P09-ERR-0504
$$
\epsilon :
(X_{affine, \etale}, \mathcal{O}_X)
\longrightarrow
(X_{affine}, \mathcal{O}_X)
$$
，如《位点上的上同调》第 \ref{sites-cohomology-section-compare} 节。
本节将把 $D_\QCoh(\mathcal{O}_X)$ 等同于
《位点上的上同调》第 \ref{sites-cohomology-section-modules-cohomology} 节
中引入的范畴 $\mathit{QC}(X_{affine}, \mathcal{O}_X)$。

\begin{lemma}
\label{spaces-perfect-lemma-DQCoh-alternative-small}
在上述情形中，下列三角范畴之间存在典范正合等价：
\begin{enumerate}
\item $D_\QCoh(\mathcal{O}_X)$；
\item $D_\QCoh(X_{affine, \etale}, \mathcal{O}_X)$；
\item $D_\QCoh(X_{affine}, \mathcal{O}_X)$；以及
\item $\mathit{QC}(X_{affine}, \mathcal{O}_X)$。
\end{enumerate}
\end{lemma}

\begin{proof}
若 $U \to V \to X$ 是平展态射且 $U$ 和 $V$ 仿射，
则环同态 $\mathcal{O}_X(V) \to \mathcal{O}_X(U)$ 平坦。
因此，(3) 与 (4) 之间的等价是《位点上的上同调》引理
\ref{sites-cohomology-lemma-cartesian-flat-quasi-coherent} 的特殊情形
（其证明也阐明了该陈述）。

\medskip\noindent
引理之前的讨论表明，我们有带环拓扑斯的等价
$(\Sh(X_{affine, \etale}), \mathcal{O}_X) \to (\Sh(X_\etale), \mathcal{O}_X)$，
从而有模的阿贝尔范畴之间的等价。由于拟凝聚模的概念是内蕴的
（《位点上的模》引理 \ref{sites-modules-lemma-special-locally-free}），
可知该等价保持拟凝聚模的子范畴。因而我们得到 (1) 与 (2)
中三角范畴之间的典范正合等价。

\medskip\noindent
为得到 (2) 与 (3) 中三角范畴之间的正合等价，我们把《位点上的上同调》
引理 \ref{sites-cohomology-lemma-compare-topologies-derived-adequate-modules}
应用于上面的态射 $\epsilon : (X_{affine, \etale}, \mathcal{O}_X) \to
(X_{affine}, \mathcal{O}_X)$。取
$\mathcal{B} = \Ob(X_{affine})$，并取
$\mathcal{A} \subset \textit{PMod}(X_{affine}, \mathcal{O})$
为满足下述条件的预层 $\mathcal{F}$ 所成的全子范畴：
$\mathcal{F}(V) \otimes_{\mathcal{O}_X(V)} \mathcal{O}_X(U)
\to \mathcal{F}(U)$ 是同构。注意，由《空间上的下降》引理
\ref{spaces-descent-lemma-quasi-coherent-alternative-small}，
$\mathcal{A}$ 的对象恰为平展拓扑中那些在
$(X_{affine, \etale}, \mathcal{O}_X)$ 上是拟凝聚模的层。
另一方面，由《位点上的模》引理
\ref{sites-modules-lemma-chaotic-quasi-coherent}，
$\mathcal{A}$ 的对象恰为 $(X_{affine}, \mathcal{O}_X)$ 上，
即混沌拓扑中的拟凝聚模。因此，若能说明《位点上的上同调》引理
\ref{sites-cohomology-lemma-compare-topologies-derived-adequate-modules}
适用，那么使用 $\epsilon^*$ 和 $R\epsilon_*$，确实得到
类别 (2) 与 (3) 之间的典范等价。

\medskip\noindent
我们必须验证四个条件：
\begin{enumerate}
\item $\mathcal{A}$ 的每个对象都是平展拓扑的层。这一点已在上面看到。
\item $\mathcal{A}$ 是
$\textit{Mod}(X_{affine, \etale}, \mathcal{O}_X)$ 的弱 Serre 子范畴。
上面已经看到
$\mathcal{A} = \QCoh(X_{affine, \etale}, \mathcal{O}_X)$，
并且也已看到，在等价
$\textit{Mod}(X_{affine, \etale}, \mathcal{O}) =
\textit{Mod}(X_\etale, \mathcal{O}_X)$ 下，它们对应于 $X$ 上的拟凝聚模。
所以结论由《空间的性质》引理
\ref{spaces-properties-lemma-properties-quasi-coherent} 和《同调》引理
\ref{homology-lemma-characterize-weak-serre-subcategory} 得出。% P09-ERR-0505
\item $X_{affine}$ 的每个对象在混沌拓扑中都有一个覆盖，其成员属于
$\mathcal{B}$。这是因为 $\mathcal{B}$ 包含所有对象。
\item 对每个对象 $U$（作为 $X_{affine}$ 的对象）和每个
$\mathcal{F}$（属于 $\mathcal{A}$），有
$H^p_{Zar}(U, \mathcal{F}) = 0$，其中 $p > 0$。这由仿射概形上拟凝聚模的上同调消失得出；
见《空间的上同调》第 \ref{spaces-cohomology-section-higher-direct-image} 节
中的讨论以及《概形的上同调》引理
\ref{coherent-lemma-quasi-coherent-affine-cohomology-zero}。
\end{enumerate}
证明完成。
\end{proof}

\begin{remark}
\label{spaces-perfect-remark-QC-compare}
设 $S$ 为概形，$X$ 为 $S$ 上的代数空间。稍后我们将证明
$\mathit{QC}((\textit{Aff}/X), \mathcal{O})$
也典范等价于 $D_\QCoh(\mathcal{O}_X)$。
见《栈上的层》命题 \ref{stacks-sheaves-proposition-QC-compare}。
\end{remark}
